% Môn: Vật lý
% Lớp: 12
% Chương: Chương IV. Vật lí hạt nhân
% Bài: L12C4 Bài 24. Công nghiệp hạt nhân
% Số câu: 189
% App đọc file này trực tiếp — sửa rồi Commit trên GitHub, bấm Đồng bộ GitHub .tex

% L12C4 Bài 24. Công nghiệp hạt nhân

% ===== Câu 1 =====
% ID: L12C4B24-01-DS
% Mức: TH
\dangbt{Cấu tạo và nguyên lí hoạt động của lò phản ứng hạt nhân}
\begin{ex}
% ID: L12C4B24-01-DS
% Mức: TH
Xét cấu tạo và nguyên lí hoạt động của lò phản ứng hạt nhân (tình huống 1). Hãy xác định tính đúng, sai của bốn phát biểu sau.
\choiceTF
{\True Lò phản ứng dùng nhiên liệu phân hạch, chất làm chậm, thanh điều khiển, chất tải nhiệt và hệ che chắn để duy trì phản ứng có kiểm soát.}
{Thanh điều khiển có nhiệm vụ làm tăng không giới hạn số neutron trong lò.}
{\True Khi giải cần đọc đúng dữ kiện, dùng hệ đơn vị thống nhất và kiểm tra điều kiện áp dụng.}
{Có thể bỏ qua dữ kiện, đơn vị và ý nghĩa vật lí của kết quả.}
\loigiai{
\begin{itemchoice}
\item (Đúng) Lò phản ứng dùng nhiên liệu phân hạch, chất làm chậm, thanh điều khiển, chất tải nhiệt và hệ che chắn để duy trì phản ứng có kiểm soát. (Vì): Lò phản ứng dùng nhiên liệu phân hạch, chất làm chậm, thanh điều khiển, chất tải nhiệt và hệ che chắn để duy trì phản ứng có kiểm soát. b) Sai: Thanh điều khiển có nhiệm vụ làm tăng không giới hạn số neutron trong lò. c) Đúng: Khi giải cần đọc đúng dữ kiện, dùng hệ đơn vị thống nhất và kiểm tra điều kiện áp dụng. d) Sai: Có thể bỏ qua dữ kiện, đơn vị và ý nghĩa vật lí của kết quả. Kiến thức cốt lõi: Lò phản ứng dùng nhiên liệu phân hạch, chất làm chậm, thanh điều khiển, chất tải nhiệt và hệ che chắn để duy trì phản ứng có kiểm soát
\item (Sai) Thanh điều khiển có nhiệm vụ làm tăng không giới hạn số neutron trong lò.
\item (Đúng) Khi giải cần đọc đúng dữ kiện, dùng hệ đơn vị thống nhất và kiểm tra điều kiện áp dụng.
\item (Sai) Có thể bỏ qua dữ kiện, đơn vị và ý nghĩa vật lí của kết quả.
\end{itemchoice}
}
\end{ex}

% ===== Câu 2 =====
% ID: L12C4B24-01-TL
% Mức: TH
\begin{bt}
% ID: L12C4B24-01-TL
% Mức: TH
Trình bày kiến thức cốt lõi của cấu tạo và nguyên lí hoạt động của lò phản ứng hạt nhân và nêu điều kiện áp dụng.
\loigiai{
Cần nêu được: Lò phản ứng dùng nhiên liệu phân hạch, chất làm chậm, thanh điều khiển, chất tải nhiệt và hệ che chắn để duy trì phản ứng có kiểm soát. Khi vận dụng phải xác định đúng dữ kiện, điều kiện áp dụng, hệ đơn vị và kiểm tra tính hợp lí của kết quả. Nhận định “Thanh điều khiển có nhiệm vụ làm tăng không giới hạn số neutron trong lò.” sai vì trái với kiến thức trên.
}
\end{bt}

% ===== Câu 3 =====
% ID: L12C4B24-01-TLN
% Mức: TH
\begin{ex}
% ID: L12C4B24-01-TLN
% Mức: TH
Trong bốn phát biểu sau về cấu tạo và nguyên lí hoạt động của lò phản ứng hạt nhân, có bao nhiêu phát biểu đúng? Chỉ ghi một số nguyên. (1) Lò phản ứng dùng nhiên liệu phân hạch, chất làm chậm, thanh điều khiển, chất tải nhiệt và hệ che chắn để duy trì phản ứng có kiểm soát. (2) Thanh điều khiển có nhiệm vụ làm tăng không giới hạn số neutron trong lò. (3) Cần kiểm tra điều kiện áp dụng trước khi tính toán. (4) Có thể bỏ qua đơn vị của kết quả.
\shortans{2}
\loigiai{
Đối chiếu từng phát biểu với kiến thức cốt lõi: Lò phản ứng dùng nhiên liệu phân hạch, chất làm chậm, thanh điều khiển, chất tải nhiệt và hệ che chắn để duy trì phản ứng có kiểm soát. Có 2 phát biểu đúng.
}
\end{ex}

% ===== Câu 4 =====
% ID: L12C4B24-01-TN
% Mức: NB
\begin{ex}
% ID: L12C4B24-01-TN
% Mức: NB
Phát biểu nào sau đây đúng về cấu tạo và nguyên lí hoạt động của lò phản ứng hạt nhân?
\choice
{\True Lò phản ứng dùng nhiên liệu phân hạch, chất làm chậm, thanh điều khiển, chất tải nhiệt và hệ che chắn để duy trì phản ứng có kiểm soát.}
{Thanh điều khiển có nhiệm vụ làm tăng không giới hạn số neutron trong lò.}
{Có thể bỏ qua điều kiện áp dụng và đơn vị trong mọi trường hợp.}
{Kết quả của bài toán không phụ thuộc dữ kiện đã cho.}
\loigiai{
Lò phản ứng dùng nhiên liệu phân hạch, chất làm chậm, thanh điều khiển, chất tải nhiệt và hệ che chắn để duy trì phản ứng có kiểm soát. Vì vậy phương án $A$ đúng; các phương án còn lại trái với kiến thức hoặc điều kiện áp dụng.
}
\end{ex}

% ===== Câu 5 =====
% ID: L12C4B24-02-DS
% Mức: TH
\begin{ex}
% ID: L12C4B24-02-DS
% Mức: TH
Xét cấu tạo và nguyên lí hoạt động của lò phản ứng hạt nhân (tình huống 2). Hãy xác định tính đúng, sai của bốn phát biểu sau.
\choiceTF
{Thanh điều khiển có nhiệm vụ làm tăng không giới hạn số neutron trong lò.}
{\True Lò phản ứng dùng nhiên liệu phân hạch, chất làm chậm, thanh điều khiển, chất tải nhiệt và hệ che chắn để duy trì phản ứng có kiểm soát.}
{Có thể bỏ qua dữ kiện, đơn vị và ý nghĩa vật lí của kết quả.}
{\True Khi giải cần đọc đúng dữ kiện, dùng hệ đơn vị thống nhất và kiểm tra điều kiện áp dụng.}
\loigiai{
\begin{itemchoice}
\item (Sai) Thanh điều khiển có nhiệm vụ làm tăng không giới hạn số neutron trong lò. (Vì): Thanh điều khiển có nhiệm vụ làm tăng không giới hạn số neutron trong lò. b) Đúng: Lò phản ứng dùng nhiên liệu phân hạch, chất làm chậm, thanh điều khiển, chất tải nhiệt và hệ che chắn để duy trì phản ứng có kiểm soát. c) Sai: Có thể bỏ qua dữ kiện, đơn vị và ý nghĩa vật lí của kết quả. d) Đúng: Khi giải cần đọc đúng dữ kiện, dùng hệ đơn vị thống nhất và kiểm tra điều kiện áp dụng. Kiến thức cốt lõi: Lò phản ứng dùng nhiên liệu phân hạch, chất làm chậm, thanh điều khiển, chất tải nhiệt và hệ che chắn để duy trì phản ứng có kiểm soát
\item (Đúng) Lò phản ứng dùng nhiên liệu phân hạch, chất làm chậm, thanh điều khiển, chất tải nhiệt và hệ che chắn để duy trì phản ứng có kiểm soát.
\item (Sai) Có thể bỏ qua dữ kiện, đơn vị và ý nghĩa vật lí của kết quả.
\item (Đúng) Khi giải cần đọc đúng dữ kiện, dùng hệ đơn vị thống nhất và kiểm tra điều kiện áp dụng.
\end{itemchoice}
}
\end{ex}

% ===== Câu 6 =====
% ID: L12C4B24-02-TL
% Mức: TH
\begin{bt}
% ID: L12C4B24-02-TL
% Mức: TH
Giải thích vì sao nhận định sau đúng: “Lò phản ứng dùng nhiên liệu phân hạch, chất làm chậm, thanh điều khiển, chất tải nhiệt và hệ che chắn để duy trì phản ứng có kiểm soát.”
\loigiai{
Cần nêu được: Lò phản ứng dùng nhiên liệu phân hạch, chất làm chậm, thanh điều khiển, chất tải nhiệt và hệ che chắn để duy trì phản ứng có kiểm soát. Khi vận dụng phải xác định đúng dữ kiện, điều kiện áp dụng, hệ đơn vị và kiểm tra tính hợp lí của kết quả. Nhận định “Thanh điều khiển có nhiệm vụ làm tăng không giới hạn số neutron trong lò.” sai vì trái với kiến thức trên.
}
\end{bt}

% ===== Câu 7 =====
% ID: L12C4B24-02-TLN
% Mức: TH
\begin{ex}
% ID: L12C4B24-02-TLN
% Mức: TH
Trong bốn phát biểu sau về cấu tạo và nguyên lí hoạt động của lò phản ứng hạt nhân, có bao nhiêu phát biểu đúng? Chỉ ghi một số nguyên. (1) Thanh điều khiển có nhiệm vụ làm tăng không giới hạn số neutron trong lò. (2) Lò phản ứng dùng nhiên liệu phân hạch, chất làm chậm, thanh điều khiển, chất tải nhiệt và hệ che chắn để duy trì phản ứng có kiểm soát. (3) Kết quả phải phù hợp ý nghĩa vật lí. (4) Mọi đại lượng đều có cùng bản chất.
\shortans{2}
\loigiai{
Đối chiếu từng phát biểu với kiến thức cốt lõi: Lò phản ứng dùng nhiên liệu phân hạch, chất làm chậm, thanh điều khiển, chất tải nhiệt và hệ che chắn để duy trì phản ứng có kiểm soát. Có 2 phát biểu đúng.
}
\end{ex}

% ===== Câu 8 =====
% ID: L12C4B24-02-TN
% Mức: NB
\begin{ex}
% ID: L12C4B24-02-TN
% Mức: NB
Nhận định nào phù hợp với kiến thức Vật lí về cấu tạo và nguyên lí hoạt động của lò phản ứng hạt nhân?
\choice
{Thanh điều khiển có nhiệm vụ làm tăng không giới hạn số neutron trong lò.}
{\True Lò phản ứng dùng nhiên liệu phân hạch, chất làm chậm, thanh điều khiển, chất tải nhiệt và hệ che chắn để duy trì phản ứng có kiểm soát.}
{Có thể bỏ qua điều kiện áp dụng và đơn vị trong mọi trường hợp.}
{Kết quả của bài toán không phụ thuộc dữ kiện đã cho.}
\loigiai{
Lò phản ứng dùng nhiên liệu phân hạch, chất làm chậm, thanh điều khiển, chất tải nhiệt và hệ che chắn để duy trì phản ứng có kiểm soát. Vì vậy phương án $B$ đúng; các phương án còn lại trái với kiến thức hoặc điều kiện áp dụng.
}
\end{ex}

% ===== Câu 9 =====
% ID: L12C4B24-03-DS
% Mức: VD
\begin{ex}
% ID: L12C4B24-03-DS
% Mức: VD
Xét cấu tạo và nguyên lí hoạt động của lò phản ứng hạt nhân (tình huống 3). Hãy xác định tính đúng, sai của bốn phát biểu sau.
\choiceTF
{\True Lò phản ứng dùng nhiên liệu phân hạch, chất làm chậm, thanh điều khiển, chất tải nhiệt và hệ che chắn để duy trì phản ứng có kiểm soát.}
{Thanh điều khiển có nhiệm vụ làm tăng không giới hạn số neutron trong lò.}
{\True Khi giải cần đọc đúng dữ kiện, dùng hệ đơn vị thống nhất và kiểm tra điều kiện áp dụng.}
{Có thể bỏ qua dữ kiện, đơn vị và ý nghĩa vật lí của kết quả.}
\loigiai{
\begin{itemchoice}
\item (Đúng) Lò phản ứng dùng nhiên liệu phân hạch, chất làm chậm, thanh điều khiển, chất tải nhiệt và hệ che chắn để duy trì phản ứng có kiểm soát. (Vì): Lò phản ứng dùng nhiên liệu phân hạch, chất làm chậm, thanh điều khiển, chất tải nhiệt và hệ che chắn để duy trì phản ứng có kiểm soát. b) Sai: Thanh điều khiển có nhiệm vụ làm tăng không giới hạn số neutron trong lò. c) Đúng: Khi giải cần đọc đúng dữ kiện, dùng hệ đơn vị thống nhất và kiểm tra điều kiện áp dụng. d) Sai: Có thể bỏ qua dữ kiện, đơn vị và ý nghĩa vật lí của kết quả. Kiến thức cốt lõi: Lò phản ứng dùng nhiên liệu phân hạch, chất làm chậm, thanh điều khiển, chất tải nhiệt và hệ che chắn để duy trì phản ứng có kiểm soát
\item (Sai) Thanh điều khiển có nhiệm vụ làm tăng không giới hạn số neutron trong lò.
\item (Đúng) Khi giải cần đọc đúng dữ kiện, dùng hệ đơn vị thống nhất và kiểm tra điều kiện áp dụng.
\item (Sai) Có thể bỏ qua dữ kiện, đơn vị và ý nghĩa vật lí của kết quả.
\end{itemchoice}
}
\end{ex}

% ===== Câu 10 =====
% ID: L12C4B24-03-TL
% Mức: VD
\begin{bt}
% ID: L12C4B24-03-TL
% Mức: VD
Phân tích sai lầm trong nhận định: “Thanh điều khiển có nhiệm vụ làm tăng không giới hạn số neutron trong lò.”
\loigiai{
Cần nêu được: Lò phản ứng dùng nhiên liệu phân hạch, chất làm chậm, thanh điều khiển, chất tải nhiệt và hệ che chắn để duy trì phản ứng có kiểm soát. Khi vận dụng phải xác định đúng dữ kiện, điều kiện áp dụng, hệ đơn vị và kiểm tra tính hợp lí của kết quả. Nhận định “Thanh điều khiển có nhiệm vụ làm tăng không giới hạn số neutron trong lò.” sai vì trái với kiến thức trên.
}
\end{bt}

% ===== Câu 11 =====
% ID: L12C4B24-03-TLN
% Mức: TH
\begin{ex}
% ID: L12C4B24-03-TLN
% Mức: TH
Trong bốn phát biểu sau về cấu tạo và nguyên lí hoạt động của lò phản ứng hạt nhân, có bao nhiêu phát biểu đúng? Chỉ ghi một số nguyên. (1) Cần đổi các đại lượng về hệ đơn vị thống nhất. (2) Lò phản ứng dùng nhiên liệu phân hạch, chất làm chậm, thanh điều khiển, chất tải nhiệt và hệ che chắn để duy trì phản ứng có kiểm soát. (3) Thanh điều khiển có nhiệm vụ làm tăng không giới hạn số neutron trong lò. (4) Không cần kiểm tra dữ kiện.
\shortans{2}
\loigiai{
Đối chiếu từng phát biểu với kiến thức cốt lõi: Lò phản ứng dùng nhiên liệu phân hạch, chất làm chậm, thanh điều khiển, chất tải nhiệt và hệ che chắn để duy trì phản ứng có kiểm soát. Có 2 phát biểu đúng.
}
\end{ex}

% ===== Câu 12 =====
% ID: L12C4B24-03-TN
% Mức: NB
\begin{ex}
% ID: L12C4B24-03-TN
% Mức: NB
Chọn kết luận chính xác về cấu tạo và nguyên lí hoạt động của lò phản ứng hạt nhân?
\choice
{Thanh điều khiển có nhiệm vụ làm tăng không giới hạn số neutron trong lò.}
{Có thể bỏ qua điều kiện áp dụng và đơn vị trong mọi trường hợp.}
{\True Lò phản ứng dùng nhiên liệu phân hạch, chất làm chậm, thanh điều khiển, chất tải nhiệt và hệ che chắn để duy trì phản ứng có kiểm soát.}
{Kết quả của bài toán không phụ thuộc dữ kiện đã cho.}
\loigiai{
Lò phản ứng dùng nhiên liệu phân hạch, chất làm chậm, thanh điều khiển, chất tải nhiệt và hệ che chắn để duy trì phản ứng có kiểm soát. Vì vậy phương án $C$ đúng; các phương án còn lại trái với kiến thức hoặc điều kiện áp dụng.
}
\end{ex}

% ===== Câu 13 =====
% ID: L12C4B24-04-DS
% Mức: VD
\begin{ex}
% ID: L12C4B24-04-DS
% Mức: VD
Xét cấu tạo và nguyên lí hoạt động của lò phản ứng hạt nhân (tình huống 4). Hãy xác định tính đúng, sai của bốn phát biểu sau.
\choiceTF
{Thanh điều khiển có nhiệm vụ làm tăng không giới hạn số neutron trong lò.}
{\True Lò phản ứng dùng nhiên liệu phân hạch, chất làm chậm, thanh điều khiển, chất tải nhiệt và hệ che chắn để duy trì phản ứng có kiểm soát.}
{Có thể bỏ qua dữ kiện, đơn vị và ý nghĩa vật lí của kết quả.}
{\True Khi giải cần đọc đúng dữ kiện, dùng hệ đơn vị thống nhất và kiểm tra điều kiện áp dụng.}
\loigiai{
\begin{itemchoice}
\item (Sai) Thanh điều khiển có nhiệm vụ làm tăng không giới hạn số neutron trong lò. (Vì): Thanh điều khiển có nhiệm vụ làm tăng không giới hạn số neutron trong lò. b) Đúng: Lò phản ứng dùng nhiên liệu phân hạch, chất làm chậm, thanh điều khiển, chất tải nhiệt và hệ che chắn để duy trì phản ứng có kiểm soát. c) Sai: Có thể bỏ qua dữ kiện, đơn vị và ý nghĩa vật lí của kết quả. d) Đúng: Khi giải cần đọc đúng dữ kiện, dùng hệ đơn vị thống nhất và kiểm tra điều kiện áp dụng. Kiến thức cốt lõi: Lò phản ứng dùng nhiên liệu phân hạch, chất làm chậm, thanh điều khiển, chất tải nhiệt và hệ che chắn để duy trì phản ứng có kiểm soát
\item (Đúng) Lò phản ứng dùng nhiên liệu phân hạch, chất làm chậm, thanh điều khiển, chất tải nhiệt và hệ che chắn để duy trì phản ứng có kiểm soát.
\item (Sai) Có thể bỏ qua dữ kiện, đơn vị và ý nghĩa vật lí của kết quả.
\item (Đúng) Khi giải cần đọc đúng dữ kiện, dùng hệ đơn vị thống nhất và kiểm tra điều kiện áp dụng.
\end{itemchoice}
}
\end{ex}

% ===== Câu 14 =====
% ID: L12C4B24-04-TL
% Mức: VD
\begin{bt}
% ID: L12C4B24-04-TL
% Mức: VD
Đề xuất các bước giải một bài toán thuộc dạng cấu tạo và nguyên lí hoạt động của lò phản ứng hạt nhân.
\loigiai{
Cần nêu được: Lò phản ứng dùng nhiên liệu phân hạch, chất làm chậm, thanh điều khiển, chất tải nhiệt và hệ che chắn để duy trì phản ứng có kiểm soát. Khi vận dụng phải xác định đúng dữ kiện, điều kiện áp dụng, hệ đơn vị và kiểm tra tính hợp lí của kết quả. Nhận định “Thanh điều khiển có nhiệm vụ làm tăng không giới hạn số neutron trong lò.” sai vì trái với kiến thức trên.
}
\end{bt}

% ===== Câu 15 =====
% ID: L12C4B24-04-TLN
% Mức: VD
\begin{ex}
% ID: L12C4B24-04-TLN
% Mức: VD
Trong bốn phát biểu sau về cấu tạo và nguyên lí hoạt động của lò phản ứng hạt nhân, có bao nhiêu phát biểu đúng? Chỉ ghi một số nguyên. (1) Lò phản ứng dùng nhiên liệu phân hạch, chất làm chậm, thanh điều khiển, chất tải nhiệt và hệ che chắn để duy trì phản ứng có kiểm soát. (2) Cần đối chiếu kết quả với điều kiện bài toán. (3) Thanh điều khiển có nhiệm vụ làm tăng không giới hạn số neutron trong lò. (4) Mọi trường hợp đều cho cùng một kết quả.
\shortans{2}
\loigiai{
Đối chiếu từng phát biểu với kiến thức cốt lõi: Lò phản ứng dùng nhiên liệu phân hạch, chất làm chậm, thanh điều khiển, chất tải nhiệt và hệ che chắn để duy trì phản ứng có kiểm soát. Có 2 phát biểu đúng.
}
\end{ex}

% ===== Câu 16 =====
% ID: L12C4B24-04-TN
% Mức: TH
\begin{ex}
% ID: L12C4B24-04-TN
% Mức: TH
Nội dung nào mô tả đúng nhất về cấu tạo và nguyên lí hoạt động của lò phản ứng hạt nhân?
\choice
{Thanh điều khiển có nhiệm vụ làm tăng không giới hạn số neutron trong lò.}
{Có thể bỏ qua điều kiện áp dụng và đơn vị trong mọi trường hợp.}
{Kết quả của bài toán không phụ thuộc dữ kiện đã cho.}
{\True Lò phản ứng dùng nhiên liệu phân hạch, chất làm chậm, thanh điều khiển, chất tải nhiệt và hệ che chắn để duy trì phản ứng có kiểm soát.}
\loigiai{
Lò phản ứng dùng nhiên liệu phân hạch, chất làm chậm, thanh điều khiển, chất tải nhiệt và hệ che chắn để duy trì phản ứng có kiểm soát. Vì vậy phương án $D$ đúng; các phương án còn lại trái với kiến thức hoặc điều kiện áp dụng.
}
\end{ex}

% ===== Câu 17 =====
% ID: L12C4B24-05-DS
% Mức: VD
\begin{ex}
% ID: L12C4B24-05-DS
% Mức: VD
Xét cấu tạo và nguyên lí hoạt động của lò phản ứng hạt nhân (tình huống 5). Hãy xác định tính đúng, sai của bốn phát biểu sau.
\choiceTF
{\True Lò phản ứng dùng nhiên liệu phân hạch, chất làm chậm, thanh điều khiển, chất tải nhiệt và hệ che chắn để duy trì phản ứng có kiểm soát.}
{Có thể bỏ qua dữ kiện, đơn vị và ý nghĩa vật lí của kết quả.}
{Thanh điều khiển có nhiệm vụ làm tăng không giới hạn số neutron trong lò.}
{Có thể bỏ qua dữ kiện, đơn vị và ý nghĩa vật lí của kết quả.}
\loigiai{
\begin{itemchoice}
\item (Đúng) Lò phản ứng dùng nhiên liệu phân hạch, chất làm chậm, thanh điều khiển, chất tải nhiệt và hệ che chắn để duy trì phản ứng có kiểm soát. (Vì): Lò phản ứng dùng nhiên liệu phân hạch, chất làm chậm, thanh điều khiển, chất tải nhiệt và hệ che chắn để duy trì phản ứng có kiểm soát. b) Sai: Có thể bỏ qua dữ kiện, đơn vị và ý nghĩa vật lí của kết quả. c) Sai: Thanh điều khiển có nhiệm vụ làm tăng không giới hạn số neutron trong lò. d) Sai: Có thể bỏ qua dữ kiện, đơn vị và ý nghĩa vật lí của kết quả. Kiến thức cốt lõi: Lò phản ứng dùng nhiên liệu phân hạch, chất làm chậm, thanh điều khiển, chất tải nhiệt và hệ che chắn để duy trì phản ứng có kiểm soát
\item (Sai) Có thể bỏ qua dữ kiện, đơn vị và ý nghĩa vật lí của kết quả.
\item (Sai) Thanh điều khiển có nhiệm vụ làm tăng không giới hạn số neutron trong lò.
\item (Sai) Có thể bỏ qua dữ kiện, đơn vị và ý nghĩa vật lí của kết quả.
\end{itemchoice}
}
\end{ex}

% ===== Câu 18 =====
% ID: L12C4B24-05-TL
% Mức: VD
\begin{bt}
% ID: L12C4B24-05-TL
% Mức: VD
Nêu một tình huống thực tiễn liên quan đến cấu tạo và nguyên lí hoạt động của lò phản ứng hạt nhân và giải thích bằng kiến thức Vật lí.
\loigiai{
Cần nêu được: Lò phản ứng dùng nhiên liệu phân hạch, chất làm chậm, thanh điều khiển, chất tải nhiệt và hệ che chắn để duy trì phản ứng có kiểm soát. Khi vận dụng phải xác định đúng dữ kiện, điều kiện áp dụng, hệ đơn vị và kiểm tra tính hợp lí của kết quả. Nhận định “Thanh điều khiển có nhiệm vụ làm tăng không giới hạn số neutron trong lò.” sai vì trái với kiến thức trên.
}
\end{bt}

% ===== Câu 19 =====
% ID: L12C4B24-05-TLN
% Mức: VD
\begin{ex}
% ID: L12C4B24-05-TLN
% Mức: VD
Trong bốn phát biểu sau về cấu tạo và nguyên lí hoạt động của lò phản ứng hạt nhân, có bao nhiêu phát biểu đúng? Chỉ ghi một số nguyên. (1) Thanh điều khiển có nhiệm vụ làm tăng không giới hạn số neutron trong lò. (2) Cần đọc đúng dữ kiện và giả thiết. (3) Lò phản ứng dùng nhiên liệu phân hạch, chất làm chậm, thanh điều khiển, chất tải nhiệt và hệ che chắn để duy trì phản ứng có kiểm soát. (4) Có thể thay số trước khi chọn công thức.
\shortans{2}
\loigiai{
Đối chiếu từng phát biểu với kiến thức cốt lõi: Lò phản ứng dùng nhiên liệu phân hạch, chất làm chậm, thanh điều khiển, chất tải nhiệt và hệ che chắn để duy trì phản ứng có kiểm soát. Có 2 phát biểu đúng.
}
\end{ex}

% ===== Câu 20 =====
% ID: L12C4B24-05-TN
% Mức: TH
\begin{ex}
% ID: L12C4B24-05-TN
% Mức: TH
Khi phân tích, phát biểu nào đúng về cấu tạo và nguyên lí hoạt động của lò phản ứng hạt nhân?
\choice
{\True Lò phản ứng dùng nhiên liệu phân hạch, chất làm chậm, thanh điều khiển, chất tải nhiệt và hệ che chắn để duy trì phản ứng có kiểm soát.}
{Thanh điều khiển có nhiệm vụ làm tăng không giới hạn số neutron trong lò.}
{Có thể bỏ qua điều kiện áp dụng và đơn vị trong mọi trường hợp.}
{Kết quả của bài toán không phụ thuộc dữ kiện đã cho.}
\loigiai{
Lò phản ứng dùng nhiên liệu phân hạch, chất làm chậm, thanh điều khiển, chất tải nhiệt và hệ che chắn để duy trì phản ứng có kiểm soát. Vì vậy phương án $A$ đúng; các phương án còn lại trái với kiến thức hoặc điều kiện áp dụng.
}
\end{ex}

% ===== Câu 21 =====
% ID: L12C4B24-06-DS
% Mức: TH
\dangbt{Nhận biết phản ứng phân hạch và phản ứng dây chuyền}
\begin{ex}
% ID: L12C4B24-06-DS
% Mức: TH
Xét nhận biết phản ứng phân hạch và phản ứng dây chuyền (tình huống 1). Hãy xác định tính đúng, sai của bốn phát biểu sau.
\choiceTF
{\True Phân hạch là sự vỡ của hạt nhân nặng thành các hạt nhân nhẹ hơn, thường kèm neutron và năng lượng; neutron có thể duy trì phản ứng dây chuyền.}
{Phân hạch là sự kết hợp hai hạt nhân nhẹ thành hạt nhân nặng.}
{\True Khi giải cần đọc đúng dữ kiện, dùng hệ đơn vị thống nhất và kiểm tra điều kiện áp dụng.}
{Có thể bỏ qua dữ kiện, đơn vị và ý nghĩa vật lí của kết quả.}
\loigiai{
\begin{itemchoice}
\item (Đúng) Phân hạch là sự vỡ của hạt nhân nặng thành các hạt nhân nhẹ hơn, thường kèm neutron và năng lượng; neutron có thể duy trì phản ứng dây chuyền. (Vì): Phân hạch là sự vỡ của hạt nhân nặng thành các hạt nhân nhẹ hơn, thường kèm neutron và năng lượng; neutron có thể duy trì phản ứng dây chuyền. b) Sai: Phân hạch là sự kết hợp hai hạt nhân nhẹ thành hạt nhân nặng. c) Đúng: Khi giải cần đọc đúng dữ kiện, dùng hệ đơn vị thống nhất và kiểm tra điều kiện áp dụng. d) Sai: Có thể bỏ qua dữ kiện, đơn vị và ý nghĩa vật lí của kết quả. Kiến thức cốt lõi: Phân hạch là sự vỡ của hạt nhân nặng thành các hạt nhân nhẹ hơn, thường kèm neutron và năng lượng; neutron có thể duy trì phản ứng dây chuyền
\item (Sai) Phân hạch là sự kết hợp hai hạt nhân nhẹ thành hạt nhân nặng.
\item (Đúng) Khi giải cần đọc đúng dữ kiện, dùng hệ đơn vị thống nhất và kiểm tra điều kiện áp dụng.
\item (Sai) Có thể bỏ qua dữ kiện, đơn vị và ý nghĩa vật lí của kết quả.
\end{itemchoice}
}
\end{ex}

% ===== Câu 22 =====
% ID: L12C4B24-06-TL
% Mức: VDC
\dangbt{Cấu tạo và nguyên lí hoạt động của lò phản ứng hạt nhân}
\begin{bt}
% ID: L12C4B24-06-TL
% Mức: VDC
So sánh nhận định đúng “Lò phản ứng dùng nhiên liệu phân hạch, chất làm chậm, thanh điều khiển, chất tải nhiệt và hệ che chắn để duy trì phản ứng có kiểm soát.” với nhận định sai “Thanh điều khiển có nhiệm vụ làm tăng không giới hạn số neutron trong lò.”; chỉ rõ căn cứ Vật lí.
\loigiai{
Cần nêu được: Lò phản ứng dùng nhiên liệu phân hạch, chất làm chậm, thanh điều khiển, chất tải nhiệt và hệ che chắn để duy trì phản ứng có kiểm soát. Khi vận dụng phải xác định đúng dữ kiện, điều kiện áp dụng, hệ đơn vị và kiểm tra tính hợp lí của kết quả. Nhận định “Thanh điều khiển có nhiệm vụ làm tăng không giới hạn số neutron trong lò.” sai vì trái với kiến thức trên.
}
\end{bt}

% ===== Câu 23 =====
% ID: L12C4B24-06-TLN
% Mức: VDC
\begin{ex}
% ID: L12C4B24-06-TLN
% Mức: VDC
Trong bốn phát biểu sau về cấu tạo và nguyên lí hoạt động của lò phản ứng hạt nhân, có bao nhiêu phát biểu đúng? Chỉ ghi một số nguyên. (1) Kết quả cần có ý nghĩa vật lí. (2) Thanh điều khiển có nhiệm vụ làm tăng không giới hạn số neutron trong lò. (3) Phải dùng đúng định luật cho tình huống. (4) Lò phản ứng dùng nhiên liệu phân hạch, chất làm chậm, thanh điều khiển, chất tải nhiệt và hệ che chắn để duy trì phản ứng có kiểm soát.
\shortans{3}
\loigiai{
Đối chiếu từng phát biểu với kiến thức cốt lõi: Lò phản ứng dùng nhiên liệu phân hạch, chất làm chậm, thanh điều khiển, chất tải nhiệt và hệ che chắn để duy trì phản ứng có kiểm soát. Có 3 phát biểu đúng.
}
\end{ex}

% ===== Câu 24 =====
% ID: L12C4B24-06-TN
% Mức: TH
\begin{ex}
% ID: L12C4B24-06-TN
% Mức: TH
Kết luận nào của học sinh là đúng về cấu tạo và nguyên lí hoạt động của lò phản ứng hạt nhân?
\choice
{Thanh điều khiển có nhiệm vụ làm tăng không giới hạn số neutron trong lò.}
{\True Lò phản ứng dùng nhiên liệu phân hạch, chất làm chậm, thanh điều khiển, chất tải nhiệt và hệ che chắn để duy trì phản ứng có kiểm soát.}
{Có thể bỏ qua điều kiện áp dụng và đơn vị trong mọi trường hợp.}
{Kết quả của bài toán không phụ thuộc dữ kiện đã cho.}
\loigiai{
Lò phản ứng dùng nhiên liệu phân hạch, chất làm chậm, thanh điều khiển, chất tải nhiệt và hệ che chắn để duy trì phản ứng có kiểm soát. Vì vậy phương án $B$ đúng; các phương án còn lại trái với kiến thức hoặc điều kiện áp dụng.
}
\end{ex}

% ===== Câu 25 =====
% ID: L12C4B24-07-DS
% Mức: TH
\dangbt{Nhận biết phản ứng phân hạch và phản ứng dây chuyền}
\begin{ex}
% ID: L12C4B24-07-DS
% Mức: TH
Xét nhận biết phản ứng phân hạch và phản ứng dây chuyền (tình huống 2). Hãy xác định tính đúng, sai của bốn phát biểu sau.
\choiceTF
{Phân hạch là sự kết hợp hai hạt nhân nhẹ thành hạt nhân nặng.}
{\True Phân hạch là sự vỡ của hạt nhân nặng thành các hạt nhân nhẹ hơn, thường kèm neutron và năng lượng; neutron có thể duy trì phản ứng dây chuyền.}
{Có thể bỏ qua dữ kiện, đơn vị và ý nghĩa vật lí của kết quả.}
{\True Khi giải cần đọc đúng dữ kiện, dùng hệ đơn vị thống nhất và kiểm tra điều kiện áp dụng.}
\loigiai{
\begin{itemchoice}
\item (Sai) Phân hạch là sự kết hợp hai hạt nhân nhẹ thành hạt nhân nặng. (Vì): Phân hạch là sự kết hợp hai hạt nhân nhẹ thành hạt nhân nặng. b) Đúng: Phân hạch là sự vỡ của hạt nhân nặng thành các hạt nhân nhẹ hơn, thường kèm neutron và năng lượng; neutron có thể duy trì phản ứng dây chuyền. c) Sai: Có thể bỏ qua dữ kiện, đơn vị và ý nghĩa vật lí của kết quả. d) Đúng: Khi giải cần đọc đúng dữ kiện, dùng hệ đơn vị thống nhất và kiểm tra điều kiện áp dụng. Kiến thức cốt lõi: Phân hạch là sự vỡ của hạt nhân nặng thành các hạt nhân nhẹ hơn, thường kèm neutron và năng lượng; neutron có thể duy trì phản ứng dây chuyền
\item (Đúng) Phân hạch là sự vỡ của hạt nhân nặng thành các hạt nhân nhẹ hơn, thường kèm neutron và năng lượng; neutron có thể duy trì phản ứng dây chuyền.
\item (Sai) Có thể bỏ qua dữ kiện, đơn vị và ý nghĩa vật lí của kết quả.
\item (Đúng) Khi giải cần đọc đúng dữ kiện, dùng hệ đơn vị thống nhất và kiểm tra điều kiện áp dụng.
\end{itemchoice}
}
\end{ex}

% ===== Câu 26 =====
% ID: L12C4B24-07-TL
% Mức: TH
\begin{bt}
% ID: L12C4B24-07-TL
% Mức: TH
Trình bày kiến thức cốt lõi của nhận biết phản ứng phân hạch và phản ứng dây chuyền và nêu điều kiện áp dụng.
\loigiai{
Cần nêu được: Phân hạch là sự vỡ của hạt nhân nặng thành các hạt nhân nhẹ hơn, thường kèm neutron và năng lượng; neutron có thể duy trì phản ứng dây chuyền. Khi vận dụng phải xác định đúng dữ kiện, điều kiện áp dụng, hệ đơn vị và kiểm tra tính hợp lí của kết quả. Nhận định “Phân hạch là sự kết hợp hai hạt nhân nhẹ thành hạt nhân nặng.” sai vì trái với kiến thức trên.
}
\end{bt}

% ===== Câu 27 =====
% ID: L12C4B24-07-TLN
% Mức: TH
\begin{ex}
% ID: L12C4B24-07-TLN
% Mức: TH
Trong bốn phát biểu sau về nhận biết phản ứng phân hạch và phản ứng dây chuyền, có bao nhiêu phát biểu đúng? Chỉ ghi một số nguyên. (1) Phân hạch là sự vỡ của hạt nhân nặng thành các hạt nhân nhẹ hơn, thường kèm neutron và năng lượng; neutron có thể duy trì phản ứng dây chuyền. (2) Phân hạch là sự kết hợp hai hạt nhân nhẹ thành hạt nhân nặng. (3) Cần kiểm tra điều kiện áp dụng trước khi tính toán. (4) Có thể bỏ qua đơn vị của kết quả.
\shortans{2}
\loigiai{
Đối chiếu từng phát biểu với kiến thức cốt lõi: Phân hạch là sự vỡ của hạt nhân nặng thành các hạt nhân nhẹ hơn, thường kèm neutron và năng lượng; neutron có thể duy trì phản ứng dây chuyền. Có 2 phát biểu đúng.
}
\end{ex}

% ===== Câu 28 =====
% ID: L12C4B24-07-TN
% Mức: TH
\dangbt{Cấu tạo và nguyên lí hoạt động của lò phản ứng hạt nhân}
\begin{ex}
% ID: L12C4B24-07-TN
% Mức: TH
Chọn phát biểu không trái với kiến thức về cấu tạo và nguyên lí hoạt động của lò phản ứng hạt nhân?
\choice
{Thanh điều khiển có nhiệm vụ làm tăng không giới hạn số neutron trong lò.}
{Có thể bỏ qua điều kiện áp dụng và đơn vị trong mọi trường hợp.}
{\True Lò phản ứng dùng nhiên liệu phân hạch, chất làm chậm, thanh điều khiển, chất tải nhiệt và hệ che chắn để duy trì phản ứng có kiểm soát.}
{Kết quả của bài toán không phụ thuộc dữ kiện đã cho.}
\loigiai{
Lò phản ứng dùng nhiên liệu phân hạch, chất làm chậm, thanh điều khiển, chất tải nhiệt và hệ che chắn để duy trì phản ứng có kiểm soát. Vì vậy phương án $C$ đúng; các phương án còn lại trái với kiến thức hoặc điều kiện áp dụng.
}
\end{ex}

% ===== Câu 29 =====
% ID: L12C4B24-08-DS
% Mức: VD
\dangbt{Nhận biết phản ứng phân hạch và phản ứng dây chuyền}
\begin{ex}
% ID: L12C4B24-08-DS
% Mức: VD
Xét nhận biết phản ứng phân hạch và phản ứng dây chuyền (tình huống 3). Hãy xác định tính đúng, sai của bốn phát biểu sau.
\choiceTF
{\True Phân hạch là sự vỡ của hạt nhân nặng thành các hạt nhân nhẹ hơn, thường kèm neutron và năng lượng; neutron có thể duy trì phản ứng dây chuyền.}
{Phân hạch là sự kết hợp hai hạt nhân nhẹ thành hạt nhân nặng.}
{\True Khi giải cần đọc đúng dữ kiện, dùng hệ đơn vị thống nhất và kiểm tra điều kiện áp dụng.}
{Có thể bỏ qua dữ kiện, đơn vị và ý nghĩa vật lí của kết quả.}
\loigiai{
\begin{itemchoice}
\item (Đúng) Phân hạch là sự vỡ của hạt nhân nặng thành các hạt nhân nhẹ hơn, thường kèm neutron và năng lượng; neutron có thể duy trì phản ứng dây chuyền. (Vì): Phân hạch là sự vỡ của hạt nhân nặng thành các hạt nhân nhẹ hơn, thường kèm neutron và năng lượng; neutron có thể duy trì phản ứng dây chuyền. b) Sai: Phân hạch là sự kết hợp hai hạt nhân nhẹ thành hạt nhân nặng. c) Đúng: Khi giải cần đọc đúng dữ kiện, dùng hệ đơn vị thống nhất và kiểm tra điều kiện áp dụng. d) Sai: Có thể bỏ qua dữ kiện, đơn vị và ý nghĩa vật lí của kết quả. Kiến thức cốt lõi: Phân hạch là sự vỡ của hạt nhân nặng thành các hạt nhân nhẹ hơn, thường kèm neutron và năng lượng; neutron có thể duy trì phản ứng dây chuyền
\item (Sai) Phân hạch là sự kết hợp hai hạt nhân nhẹ thành hạt nhân nặng.
\item (Đúng) Khi giải cần đọc đúng dữ kiện, dùng hệ đơn vị thống nhất và kiểm tra điều kiện áp dụng.
\item (Sai) Có thể bỏ qua dữ kiện, đơn vị và ý nghĩa vật lí của kết quả.
\end{itemchoice}
}
\end{ex}

% ===== Câu 30 =====
% ID: L12C4B24-08-TL
% Mức: TH
\begin{bt}
% ID: L12C4B24-08-TL
% Mức: TH
Giải thích vì sao nhận định sau đúng: “Phân hạch là sự vỡ của hạt nhân nặng thành các hạt nhân nhẹ hơn, thường kèm neutron và năng lượng; neutron có thể duy trì phản ứng dây chuyền.”
\loigiai{
Cần nêu được: Phân hạch là sự vỡ của hạt nhân nặng thành các hạt nhân nhẹ hơn, thường kèm neutron và năng lượng; neutron có thể duy trì phản ứng dây chuyền. Khi vận dụng phải xác định đúng dữ kiện, điều kiện áp dụng, hệ đơn vị và kiểm tra tính hợp lí của kết quả. Nhận định “Phân hạch là sự kết hợp hai hạt nhân nhẹ thành hạt nhân nặng.” sai vì trái với kiến thức trên.
}
\end{bt}

% ===== Câu 31 =====
% ID: L12C4B24-08-TLN
% Mức: TH
\begin{ex}
% ID: L12C4B24-08-TLN
% Mức: TH
Trong bốn phát biểu sau về nhận biết phản ứng phân hạch và phản ứng dây chuyền, có bao nhiêu phát biểu đúng? Chỉ ghi một số nguyên. (1) Phân hạch là sự kết hợp hai hạt nhân nhẹ thành hạt nhân nặng. (2) Phân hạch là sự vỡ của hạt nhân nặng thành các hạt nhân nhẹ hơn, thường kèm neutron và năng lượng; neutron có thể duy trì phản ứng dây chuyền. (3) Kết quả phải phù hợp ý nghĩa vật lí. (4) Mọi đại lượng đều có cùng bản chất.
\shortans{2}
\loigiai{
Đối chiếu từng phát biểu với kiến thức cốt lõi: Phân hạch là sự vỡ của hạt nhân nặng thành các hạt nhân nhẹ hơn, thường kèm neutron và năng lượng; neutron có thể duy trì phản ứng dây chuyền. Có 2 phát biểu đúng.
}
\end{ex}

% ===== Câu 32 =====
% ID: L12C4B24-08-TN
% Mức: VD
\dangbt{Cấu tạo và nguyên lí hoạt động của lò phản ứng hạt nhân}
\begin{ex}
% ID: L12C4B24-08-TN
% Mức: VD
Cơ sở Vật lí đúng để giải bài là gì về cấu tạo và nguyên lí hoạt động của lò phản ứng hạt nhân?
\choice
{Thanh điều khiển có nhiệm vụ làm tăng không giới hạn số neutron trong lò.}
{Có thể bỏ qua điều kiện áp dụng và đơn vị trong mọi trường hợp.}
{Kết quả của bài toán không phụ thuộc dữ kiện đã cho.}
{\True Lò phản ứng dùng nhiên liệu phân hạch, chất làm chậm, thanh điều khiển, chất tải nhiệt và hệ che chắn để duy trì phản ứng có kiểm soát.}
\loigiai{
Lò phản ứng dùng nhiên liệu phân hạch, chất làm chậm, thanh điều khiển, chất tải nhiệt và hệ che chắn để duy trì phản ứng có kiểm soát. Vì vậy phương án $D$ đúng; các phương án còn lại trái với kiến thức hoặc điều kiện áp dụng.
}
\end{ex}

% ===== Câu 33 =====
% ID: L12C4B24-09-DS
% Mức: VD
\dangbt{Nhận biết phản ứng phân hạch và phản ứng dây chuyền}
\begin{ex}
% ID: L12C4B24-09-DS
% Mức: VD
Xét nhận biết phản ứng phân hạch và phản ứng dây chuyền (tình huống 4). Hãy xác định tính đúng, sai của bốn phát biểu sau.
\choiceTF
{Phân hạch là sự kết hợp hai hạt nhân nhẹ thành hạt nhân nặng.}
{\True Phân hạch là sự vỡ của hạt nhân nặng thành các hạt nhân nhẹ hơn, thường kèm neutron và năng lượng; neutron có thể duy trì phản ứng dây chuyền.}
{Có thể bỏ qua dữ kiện, đơn vị và ý nghĩa vật lí của kết quả.}
{\True Khi giải cần đọc đúng dữ kiện, dùng hệ đơn vị thống nhất và kiểm tra điều kiện áp dụng.}
\loigiai{
\begin{itemchoice}
\item (Sai) Phân hạch là sự kết hợp hai hạt nhân nhẹ thành hạt nhân nặng. (Vì): Phân hạch là sự kết hợp hai hạt nhân nhẹ thành hạt nhân nặng. b) Đúng: Phân hạch là sự vỡ của hạt nhân nặng thành các hạt nhân nhẹ hơn, thường kèm neutron và năng lượng; neutron có thể duy trì phản ứng dây chuyền. c) Sai: Có thể bỏ qua dữ kiện, đơn vị và ý nghĩa vật lí của kết quả. d) Đúng: Khi giải cần đọc đúng dữ kiện, dùng hệ đơn vị thống nhất và kiểm tra điều kiện áp dụng. Kiến thức cốt lõi: Phân hạch là sự vỡ của hạt nhân nặng thành các hạt nhân nhẹ hơn, thường kèm neutron và năng lượng; neutron có thể duy trì phản ứng dây chuyền
\item (Đúng) Phân hạch là sự vỡ của hạt nhân nặng thành các hạt nhân nhẹ hơn, thường kèm neutron và năng lượng; neutron có thể duy trì phản ứng dây chuyền.
\item (Sai) Có thể bỏ qua dữ kiện, đơn vị và ý nghĩa vật lí của kết quả.
\item (Đúng) Khi giải cần đọc đúng dữ kiện, dùng hệ đơn vị thống nhất và kiểm tra điều kiện áp dụng.
\end{itemchoice}
}
\end{ex}

% ===== Câu 34 =====
% ID: L12C4B24-09-TL
% Mức: VD
\begin{bt}
% ID: L12C4B24-09-TL
% Mức: VD
Phân tích sai lầm trong nhận định: “Phân hạch là sự kết hợp hai hạt nhân nhẹ thành hạt nhân nặng.”
\loigiai{
Cần nêu được: Phân hạch là sự vỡ của hạt nhân nặng thành các hạt nhân nhẹ hơn, thường kèm neutron và năng lượng; neutron có thể duy trì phản ứng dây chuyền. Khi vận dụng phải xác định đúng dữ kiện, điều kiện áp dụng, hệ đơn vị và kiểm tra tính hợp lí của kết quả. Nhận định “Phân hạch là sự kết hợp hai hạt nhân nhẹ thành hạt nhân nặng.” sai vì trái với kiến thức trên.
}
\end{bt}

% ===== Câu 35 =====
% ID: L12C4B24-09-TLN
% Mức: TH
\begin{ex}
% ID: L12C4B24-09-TLN
% Mức: TH
Trong bốn phát biểu sau về nhận biết phản ứng phân hạch và phản ứng dây chuyền, có bao nhiêu phát biểu đúng? Chỉ ghi một số nguyên. (1) Cần đổi các đại lượng về hệ đơn vị thống nhất. (2) Phân hạch là sự vỡ của hạt nhân nặng thành các hạt nhân nhẹ hơn, thường kèm neutron và năng lượng; neutron có thể duy trì phản ứng dây chuyền. (3) Phân hạch là sự kết hợp hai hạt nhân nhẹ thành hạt nhân nặng. (4) Không cần kiểm tra dữ kiện.
\shortans{2}
\loigiai{
Đối chiếu từng phát biểu với kiến thức cốt lõi: Phân hạch là sự vỡ của hạt nhân nặng thành các hạt nhân nhẹ hơn, thường kèm neutron và năng lượng; neutron có thể duy trì phản ứng dây chuyền. Có 2 phát biểu đúng.
}
\end{ex}

% ===== Câu 36 =====
% ID: L12C4B24-09-TN
% Mức: VD
\dangbt{Cấu tạo và nguyên lí hoạt động của lò phản ứng hạt nhân}
\begin{ex}
% ID: L12C4B24-09-TN
% Mức: VD
Trong một bài thực hành, nhận xét nào đúng về cấu tạo và nguyên lí hoạt động của lò phản ứng hạt nhân?
\choice
{\True Lò phản ứng dùng nhiên liệu phân hạch, chất làm chậm, thanh điều khiển, chất tải nhiệt và hệ che chắn để duy trì phản ứng có kiểm soát.}
{Thanh điều khiển có nhiệm vụ làm tăng không giới hạn số neutron trong lò.}
{Có thể bỏ qua điều kiện áp dụng và đơn vị trong mọi trường hợp.}
{Kết quả của bài toán không phụ thuộc dữ kiện đã cho.}
\loigiai{
Lò phản ứng dùng nhiên liệu phân hạch, chất làm chậm, thanh điều khiển, chất tải nhiệt và hệ che chắn để duy trì phản ứng có kiểm soát. Vì vậy phương án $A$ đúng; các phương án còn lại trái với kiến thức hoặc điều kiện áp dụng.
}
\end{ex}

% ===== Câu 37 =====
% ID: L12C4B24-10-DS
% Mức: VD
\dangbt{Nhận biết phản ứng phân hạch và phản ứng dây chuyền}
\begin{ex}
% ID: L12C4B24-10-DS
% Mức: VD
Xét nhận biết phản ứng phân hạch và phản ứng dây chuyền (tình huống 5). Hãy xác định tính đúng, sai của bốn phát biểu sau.
\choiceTF
{\True Phân hạch là sự vỡ của hạt nhân nặng thành các hạt nhân nhẹ hơn, thường kèm neutron và năng lượng; neutron có thể duy trì phản ứng dây chuyền.}
{Có thể bỏ qua dữ kiện, đơn vị và ý nghĩa vật lí của kết quả.}
{Phân hạch là sự kết hợp hai hạt nhân nhẹ thành hạt nhân nặng.}
{Có thể bỏ qua dữ kiện, đơn vị và ý nghĩa vật lí của kết quả.}
\loigiai{
\begin{itemchoice}
\item (Đúng) Phân hạch là sự vỡ của hạt nhân nặng thành các hạt nhân nhẹ hơn, thường kèm neutron và năng lượng; neutron có thể duy trì phản ứng dây chuyền. (Vì): Phân hạch là sự vỡ của hạt nhân nặng thành các hạt nhân nhẹ hơn, thường kèm neutron và năng lượng; neutron có thể duy trì phản ứng dây chuyền. b) Sai: Có thể bỏ qua dữ kiện, đơn vị và ý nghĩa vật lí của kết quả. c) Sai: Phân hạch là sự kết hợp hai hạt nhân nhẹ thành hạt nhân nặng. d) Sai: Có thể bỏ qua dữ kiện, đơn vị và ý nghĩa vật lí của kết quả. Kiến thức cốt lõi: Phân hạch là sự vỡ của hạt nhân nặng thành các hạt nhân nhẹ hơn, thường kèm neutron và năng lượng; neutron có thể duy trì phản ứng dây chuyền
\item (Sai) Có thể bỏ qua dữ kiện, đơn vị và ý nghĩa vật lí của kết quả.
\item (Sai) Phân hạch là sự kết hợp hai hạt nhân nhẹ thành hạt nhân nặng.
\item (Sai) Có thể bỏ qua dữ kiện, đơn vị và ý nghĩa vật lí của kết quả.
\end{itemchoice}
}
\end{ex}

% ===== Câu 38 =====
% ID: L12C4B24-10-TL
% Mức: VD
\begin{bt}
% ID: L12C4B24-10-TL
% Mức: VD
Đề xuất các bước giải một bài toán thuộc dạng nhận biết phản ứng phân hạch và phản ứng dây chuyền.
\loigiai{
Cần nêu được: Phân hạch là sự vỡ của hạt nhân nặng thành các hạt nhân nhẹ hơn, thường kèm neutron và năng lượng; neutron có thể duy trì phản ứng dây chuyền. Khi vận dụng phải xác định đúng dữ kiện, điều kiện áp dụng, hệ đơn vị và kiểm tra tính hợp lí của kết quả. Nhận định “Phân hạch là sự kết hợp hai hạt nhân nhẹ thành hạt nhân nặng.” sai vì trái với kiến thức trên.
}
\end{bt}

% ===== Câu 39 =====
% ID: L12C4B24-10-TLN
% Mức: VD
\begin{ex}
% ID: L12C4B24-10-TLN
% Mức: VD
Trong bốn phát biểu sau về nhận biết phản ứng phân hạch và phản ứng dây chuyền, có bao nhiêu phát biểu đúng? Chỉ ghi một số nguyên. (1) Phân hạch là sự vỡ của hạt nhân nặng thành các hạt nhân nhẹ hơn, thường kèm neutron và năng lượng; neutron có thể duy trì phản ứng dây chuyền. (2) Cần đối chiếu kết quả với điều kiện bài toán. (3) Phân hạch là sự kết hợp hai hạt nhân nhẹ thành hạt nhân nặng. (4) Mọi trường hợp đều cho cùng một kết quả.
\shortans{2}
\loigiai{
Đối chiếu từng phát biểu với kiến thức cốt lõi: Phân hạch là sự vỡ của hạt nhân nặng thành các hạt nhân nhẹ hơn, thường kèm neutron và năng lượng; neutron có thể duy trì phản ứng dây chuyền. Có 2 phát biểu đúng.
}
\end{ex}

% ===== Câu 40 =====
% ID: L12C4B24-10-TN
% Mức: VD
\dangbt{Cấu tạo và nguyên lí hoạt động của lò phản ứng hạt nhân}
\begin{ex}
% ID: L12C4B24-10-TN
% Mức: VD
Kết luận đúng khi vận dụng là về cấu tạo và nguyên lí hoạt động của lò phản ứng hạt nhân?
\choice
{Thanh điều khiển có nhiệm vụ làm tăng không giới hạn số neutron trong lò.}
{\True Lò phản ứng dùng nhiên liệu phân hạch, chất làm chậm, thanh điều khiển, chất tải nhiệt và hệ che chắn để duy trì phản ứng có kiểm soát.}
{Có thể bỏ qua điều kiện áp dụng và đơn vị trong mọi trường hợp.}
{Kết quả của bài toán không phụ thuộc dữ kiện đã cho.}
\loigiai{
Lò phản ứng dùng nhiên liệu phân hạch, chất làm chậm, thanh điều khiển, chất tải nhiệt và hệ che chắn để duy trì phản ứng có kiểm soát. Vì vậy phương án $B$ đúng; các phương án còn lại trái với kiến thức hoặc điều kiện áp dụng.
}
\end{ex}

% ===== Câu 41 =====
% ID: L12C4B24-11-DS
% Mức: TH
\dangbt{Nhận biết phản ứng nhiệt hạch và điều kiện xảy ra}
\begin{ex}
% ID: L12C4B24-11-DS
% Mức: TH
Xét nhận biết phản ứng nhiệt hạch và điều kiện xảy ra (tình huống 1). Hãy xác định tính đúng, sai của bốn phát biểu sau.
\choiceTF
{\True Nhiệt hạch là sự kết hợp các hạt nhân nhẹ, cần nhiệt độ rất cao để thắng lực đẩy Coulomb và có thể tỏa năng lượng lớn.}
{Nhiệt hạch dễ xảy ra ở nhiệt độ phòng.}
{\True Khi giải cần đọc đúng dữ kiện, dùng hệ đơn vị thống nhất và kiểm tra điều kiện áp dụng.}
{Có thể bỏ qua dữ kiện, đơn vị và ý nghĩa vật lí của kết quả.}
\loigiai{
\begin{itemchoice}
\item (Đúng) Nhiệt hạch là sự kết hợp các hạt nhân nhẹ, cần nhiệt độ rất cao để thắng lực đẩy Coulomb và có thể tỏa năng lượng lớn. (Vì): Nhiệt hạch là sự kết hợp các hạt nhân nhẹ, cần nhiệt độ rất cao để thắng lực đẩy Coulomb và có thể tỏa năng lượng lớn. b) Sai: Nhiệt hạch dễ xảy ra ở nhiệt độ phòng. c) Đúng: Khi giải cần đọc đúng dữ kiện, dùng hệ đơn vị thống nhất và kiểm tra điều kiện áp dụng. d) Sai: Có thể bỏ qua dữ kiện, đơn vị và ý nghĩa vật lí của kết quả. Kiến thức cốt lõi: Nhiệt hạch là sự kết hợp các hạt nhân nhẹ, cần nhiệt độ rất cao để thắng lực đẩy Coulomb và có thể tỏa năng lượng lớn
\item (Sai) Nhiệt hạch dễ xảy ra ở nhiệt độ phòng.
\item (Đúng) Khi giải cần đọc đúng dữ kiện, dùng hệ đơn vị thống nhất và kiểm tra điều kiện áp dụng.
\item (Sai) Có thể bỏ qua dữ kiện, đơn vị và ý nghĩa vật lí của kết quả.
\end{itemchoice}
}
\end{ex}

% ===== Câu 42 =====
% ID: L12C4B24-11-TL
% Mức: VD
\dangbt{Nhận biết phản ứng phân hạch và phản ứng dây chuyền}
\begin{bt}
% ID: L12C4B24-11-TL
% Mức: VD
Nêu một tình huống thực tiễn liên quan đến nhận biết phản ứng phân hạch và phản ứng dây chuyền và giải thích bằng kiến thức Vật lí.
\loigiai{
Cần nêu được: Phân hạch là sự vỡ của hạt nhân nặng thành các hạt nhân nhẹ hơn, thường kèm neutron và năng lượng; neutron có thể duy trì phản ứng dây chuyền. Khi vận dụng phải xác định đúng dữ kiện, điều kiện áp dụng, hệ đơn vị và kiểm tra tính hợp lí của kết quả. Nhận định “Phân hạch là sự kết hợp hai hạt nhân nhẹ thành hạt nhân nặng.” sai vì trái với kiến thức trên.
}
\end{bt}

% ===== Câu 43 =====
% ID: L12C4B24-11-TLN
% Mức: VD
\begin{ex}
% ID: L12C4B24-11-TLN
% Mức: VD
Trong bốn phát biểu sau về nhận biết phản ứng phân hạch và phản ứng dây chuyền, có bao nhiêu phát biểu đúng? Chỉ ghi một số nguyên. (1) Phân hạch là sự kết hợp hai hạt nhân nhẹ thành hạt nhân nặng. (2) Cần đọc đúng dữ kiện và giả thiết. (3) Phân hạch là sự vỡ của hạt nhân nặng thành các hạt nhân nhẹ hơn, thường kèm neutron và năng lượng; neutron có thể duy trì phản ứng dây chuyền. (4) Có thể thay số trước khi chọn công thức.
\shortans{2}
\loigiai{
Đối chiếu từng phát biểu với kiến thức cốt lõi: Phân hạch là sự vỡ của hạt nhân nặng thành các hạt nhân nhẹ hơn, thường kèm neutron và năng lượng; neutron có thể duy trì phản ứng dây chuyền. Có 2 phát biểu đúng.
}
\end{ex}

% ===== Câu 44 =====
% ID: L12C4B24-11-TN
% Mức: NB
\begin{ex}
% ID: L12C4B24-11-TN
% Mức: NB
Phát biểu nào sau đây đúng về nhận biết phản ứng phân hạch và phản ứng dây chuyền?
\choice
{\True Phân hạch là sự vỡ của hạt nhân nặng thành các hạt nhân nhẹ hơn, thường kèm neutron và năng lượng; neutron có thể duy trì phản ứng dây chuyền.}
{Phân hạch là sự kết hợp hai hạt nhân nhẹ thành hạt nhân nặng.}
{Có thể bỏ qua điều kiện áp dụng và đơn vị trong mọi trường hợp.}
{Kết quả của bài toán không phụ thuộc dữ kiện đã cho.}
\loigiai{
Phân hạch là sự vỡ của hạt nhân nặng thành các hạt nhân nhẹ hơn, thường kèm neutron và năng lượng; neutron có thể duy trì phản ứng dây chuyền. Vì vậy phương án $A$ đúng; các phương án còn lại trái với kiến thức hoặc điều kiện áp dụng.
}
\end{ex}

% ===== Câu 45 =====
% ID: L12C4B24-12-DS
% Mức: TH
\dangbt{Nhận biết phản ứng nhiệt hạch và điều kiện xảy ra}
\begin{ex}
% ID: L12C4B24-12-DS
% Mức: TH
Xét nhận biết phản ứng nhiệt hạch và điều kiện xảy ra (tình huống 2). Hãy xác định tính đúng, sai của bốn phát biểu sau.
\choiceTF
{Nhiệt hạch dễ xảy ra ở nhiệt độ phòng.}
{\True Nhiệt hạch là sự kết hợp các hạt nhân nhẹ, cần nhiệt độ rất cao để thắng lực đẩy Coulomb và có thể tỏa năng lượng lớn.}
{Có thể bỏ qua dữ kiện, đơn vị và ý nghĩa vật lí của kết quả.}
{\True Khi giải cần đọc đúng dữ kiện, dùng hệ đơn vị thống nhất và kiểm tra điều kiện áp dụng.}
\loigiai{
\begin{itemchoice}
\item (Sai) Nhiệt hạch dễ xảy ra ở nhiệt độ phòng. (Vì): Nhiệt hạch dễ xảy ra ở nhiệt độ phòng. b) Đúng: Nhiệt hạch là sự kết hợp các hạt nhân nhẹ, cần nhiệt độ rất cao để thắng lực đẩy Coulomb và có thể tỏa năng lượng lớn. c) Sai: Có thể bỏ qua dữ kiện, đơn vị và ý nghĩa vật lí của kết quả. d) Đúng: Khi giải cần đọc đúng dữ kiện, dùng hệ đơn vị thống nhất và kiểm tra điều kiện áp dụng. Kiến thức cốt lõi: Nhiệt hạch là sự kết hợp các hạt nhân nhẹ, cần nhiệt độ rất cao để thắng lực đẩy Coulomb và có thể tỏa năng lượng lớn
\item (Đúng) Nhiệt hạch là sự kết hợp các hạt nhân nhẹ, cần nhiệt độ rất cao để thắng lực đẩy Coulomb và có thể tỏa năng lượng lớn.
\item (Sai) Có thể bỏ qua dữ kiện, đơn vị và ý nghĩa vật lí của kết quả.
\item (Đúng) Khi giải cần đọc đúng dữ kiện, dùng hệ đơn vị thống nhất và kiểm tra điều kiện áp dụng.
\end{itemchoice}
}
\end{ex}

% ===== Câu 46 =====
% ID: L12C4B24-12-TL
% Mức: VDC
\dangbt{Nhận biết phản ứng phân hạch và phản ứng dây chuyền}
\begin{bt}
% ID: L12C4B24-12-TL
% Mức: VDC
So sánh nhận định đúng “Phân hạch là sự vỡ của hạt nhân nặng thành các hạt nhân nhẹ hơn, thường kèm neutron và năng lượng; neutron có thể duy trì phản ứng dây chuyền.” với nhận định sai “Phân hạch là sự kết hợp hai hạt nhân nhẹ thành hạt nhân nặng.”; chỉ rõ căn cứ Vật lí.
\loigiai{
Cần nêu được: Phân hạch là sự vỡ của hạt nhân nặng thành các hạt nhân nhẹ hơn, thường kèm neutron và năng lượng; neutron có thể duy trì phản ứng dây chuyền. Khi vận dụng phải xác định đúng dữ kiện, điều kiện áp dụng, hệ đơn vị và kiểm tra tính hợp lí của kết quả. Nhận định “Phân hạch là sự kết hợp hai hạt nhân nhẹ thành hạt nhân nặng.” sai vì trái với kiến thức trên.
}
\end{bt}

% ===== Câu 47 =====
% ID: L12C4B24-12-TLN
% Mức: VDC
\begin{ex}
% ID: L12C4B24-12-TLN
% Mức: VDC
Trong bốn phát biểu sau về nhận biết phản ứng phân hạch và phản ứng dây chuyền, có bao nhiêu phát biểu đúng? Chỉ ghi một số nguyên. (1) Kết quả cần có ý nghĩa vật lí. (2) Phân hạch là sự kết hợp hai hạt nhân nhẹ thành hạt nhân nặng. (3) Phải dùng đúng định luật cho tình huống. (4) Phân hạch là sự vỡ của hạt nhân nặng thành các hạt nhân nhẹ hơn, thường kèm neutron và năng lượng; neutron có thể duy trì phản ứng dây chuyền.
\shortans{3}
\loigiai{
Đối chiếu từng phát biểu với kiến thức cốt lõi: Phân hạch là sự vỡ của hạt nhân nặng thành các hạt nhân nhẹ hơn, thường kèm neutron và năng lượng; neutron có thể duy trì phản ứng dây chuyền. Có 3 phát biểu đúng.
}
\end{ex}

% ===== Câu 48 =====
% ID: L12C4B24-12-TN
% Mức: NB
\begin{ex}
% ID: L12C4B24-12-TN
% Mức: NB
Nhận định nào phù hợp với kiến thức Vật lí về nhận biết phản ứng phân hạch và phản ứng dây chuyền?
\choice
{Phân hạch là sự kết hợp hai hạt nhân nhẹ thành hạt nhân nặng.}
{\True Phân hạch là sự vỡ của hạt nhân nặng thành các hạt nhân nhẹ hơn, thường kèm neutron và năng lượng; neutron có thể duy trì phản ứng dây chuyền.}
{Có thể bỏ qua điều kiện áp dụng và đơn vị trong mọi trường hợp.}
{Kết quả của bài toán không phụ thuộc dữ kiện đã cho.}
\loigiai{
Phân hạch là sự vỡ của hạt nhân nặng thành các hạt nhân nhẹ hơn, thường kèm neutron và năng lượng; neutron có thể duy trì phản ứng dây chuyền. Vì vậy phương án $B$ đúng; các phương án còn lại trái với kiến thức hoặc điều kiện áp dụng.
}
\end{ex}

% ===== Câu 49 =====
% ID: L12C4B24-13-DS
% Mức: VD
\dangbt{Nhận biết phản ứng nhiệt hạch và điều kiện xảy ra}
\begin{ex}
% ID: L12C4B24-13-DS
% Mức: VD
Xét nhận biết phản ứng nhiệt hạch và điều kiện xảy ra (tình huống 3). Hãy xác định tính đúng, sai của bốn phát biểu sau.
\choiceTF
{\True Nhiệt hạch là sự kết hợp các hạt nhân nhẹ, cần nhiệt độ rất cao để thắng lực đẩy Coulomb và có thể tỏa năng lượng lớn.}
{Nhiệt hạch dễ xảy ra ở nhiệt độ phòng.}
{\True Khi giải cần đọc đúng dữ kiện, dùng hệ đơn vị thống nhất và kiểm tra điều kiện áp dụng.}
{Có thể bỏ qua dữ kiện, đơn vị và ý nghĩa vật lí của kết quả.}
\loigiai{
\begin{itemchoice}
\item (Đúng) Nhiệt hạch là sự kết hợp các hạt nhân nhẹ, cần nhiệt độ rất cao để thắng lực đẩy Coulomb và có thể tỏa năng lượng lớn. (Vì): Nhiệt hạch là sự kết hợp các hạt nhân nhẹ, cần nhiệt độ rất cao để thắng lực đẩy Coulomb và có thể tỏa năng lượng lớn. b) Sai: Nhiệt hạch dễ xảy ra ở nhiệt độ phòng. c) Đúng: Khi giải cần đọc đúng dữ kiện, dùng hệ đơn vị thống nhất và kiểm tra điều kiện áp dụng. d) Sai: Có thể bỏ qua dữ kiện, đơn vị và ý nghĩa vật lí của kết quả. Kiến thức cốt lõi: Nhiệt hạch là sự kết hợp các hạt nhân nhẹ, cần nhiệt độ rất cao để thắng lực đẩy Coulomb và có thể tỏa năng lượng lớn
\item (Sai) Nhiệt hạch dễ xảy ra ở nhiệt độ phòng.
\item (Đúng) Khi giải cần đọc đúng dữ kiện, dùng hệ đơn vị thống nhất và kiểm tra điều kiện áp dụng.
\item (Sai) Có thể bỏ qua dữ kiện, đơn vị và ý nghĩa vật lí của kết quả.
\end{itemchoice}
}
\end{ex}

% ===== Câu 50 =====
% ID: L12C4B24-13-TL
% Mức: TH
\begin{bt}
% ID: L12C4B24-13-TL
% Mức: TH
Trình bày kiến thức cốt lõi của nhận biết phản ứng nhiệt hạch và điều kiện xảy ra và nêu điều kiện áp dụng.
\loigiai{
Cần nêu được: Nhiệt hạch là sự kết hợp các hạt nhân nhẹ, cần nhiệt độ rất cao để thắng lực đẩy Coulomb và có thể tỏa năng lượng lớn. Khi vận dụng phải xác định đúng dữ kiện, điều kiện áp dụng, hệ đơn vị và kiểm tra tính hợp lí của kết quả. Nhận định “Nhiệt hạch dễ xảy ra ở nhiệt độ phòng.” sai vì trái với kiến thức trên.
}
\end{bt}

% ===== Câu 51 =====
% ID: L12C4B24-13-TLN
% Mức: TH
\begin{ex}
% ID: L12C4B24-13-TLN
% Mức: TH
Trong bốn phát biểu sau về nhận biết phản ứng nhiệt hạch và điều kiện xảy ra, có bao nhiêu phát biểu đúng? Chỉ ghi một số nguyên. (1) Nhiệt hạch là sự kết hợp các hạt nhân nhẹ, cần nhiệt độ rất cao để thắng lực đẩy Coulomb và có thể tỏa năng lượng lớn. (2) Nhiệt hạch dễ xảy ra ở nhiệt độ phòng. (3) Cần kiểm tra điều kiện áp dụng trước khi tính toán. (4) Có thể bỏ qua đơn vị của kết quả.
\shortans{2}
\loigiai{
Đối chiếu từng phát biểu với kiến thức cốt lõi: Nhiệt hạch là sự kết hợp các hạt nhân nhẹ, cần nhiệt độ rất cao để thắng lực đẩy Coulomb và có thể tỏa năng lượng lớn. Có 2 phát biểu đúng.
}
\end{ex}

% ===== Câu 52 =====
% ID: L12C4B24-13-TN
% Mức: NB
\dangbt{Nhận biết phản ứng phân hạch và phản ứng dây chuyền}
\begin{ex}
% ID: L12C4B24-13-TN
% Mức: NB
Chọn kết luận chính xác về nhận biết phản ứng phân hạch và phản ứng dây chuyền?
\choice
{Phân hạch là sự kết hợp hai hạt nhân nhẹ thành hạt nhân nặng.}
{Có thể bỏ qua điều kiện áp dụng và đơn vị trong mọi trường hợp.}
{\True Phân hạch là sự vỡ của hạt nhân nặng thành các hạt nhân nhẹ hơn, thường kèm neutron và năng lượng; neutron có thể duy trì phản ứng dây chuyền.}
{Kết quả của bài toán không phụ thuộc dữ kiện đã cho.}
\loigiai{
Phân hạch là sự vỡ của hạt nhân nặng thành các hạt nhân nhẹ hơn, thường kèm neutron và năng lượng; neutron có thể duy trì phản ứng dây chuyền. Vì vậy phương án $C$ đúng; các phương án còn lại trái với kiến thức hoặc điều kiện áp dụng.
}
\end{ex}

% ===== Câu 53 =====
% ID: L12C4B24-14-DS
% Mức: VD
\dangbt{Nhận biết phản ứng nhiệt hạch và điều kiện xảy ra}
\begin{ex}
% ID: L12C4B24-14-DS
% Mức: VD
Xét nhận biết phản ứng nhiệt hạch và điều kiện xảy ra (tình huống 4). Hãy xác định tính đúng, sai của bốn phát biểu sau.
\choiceTF
{Nhiệt hạch dễ xảy ra ở nhiệt độ phòng.}
{\True Nhiệt hạch là sự kết hợp các hạt nhân nhẹ, cần nhiệt độ rất cao để thắng lực đẩy Coulomb và có thể tỏa năng lượng lớn.}
{Có thể bỏ qua dữ kiện, đơn vị và ý nghĩa vật lí của kết quả.}
{\True Khi giải cần đọc đúng dữ kiện, dùng hệ đơn vị thống nhất và kiểm tra điều kiện áp dụng.}
\loigiai{
\begin{itemchoice}
\item (Sai) Nhiệt hạch dễ xảy ra ở nhiệt độ phòng. (Vì): Nhiệt hạch dễ xảy ra ở nhiệt độ phòng. b) Đúng: Nhiệt hạch là sự kết hợp các hạt nhân nhẹ, cần nhiệt độ rất cao để thắng lực đẩy Coulomb và có thể tỏa năng lượng lớn. c) Sai: Có thể bỏ qua dữ kiện, đơn vị và ý nghĩa vật lí của kết quả. d) Đúng: Khi giải cần đọc đúng dữ kiện, dùng hệ đơn vị thống nhất và kiểm tra điều kiện áp dụng. Kiến thức cốt lõi: Nhiệt hạch là sự kết hợp các hạt nhân nhẹ, cần nhiệt độ rất cao để thắng lực đẩy Coulomb và có thể tỏa năng lượng lớn
\item (Đúng) Nhiệt hạch là sự kết hợp các hạt nhân nhẹ, cần nhiệt độ rất cao để thắng lực đẩy Coulomb và có thể tỏa năng lượng lớn.
\item (Sai) Có thể bỏ qua dữ kiện, đơn vị và ý nghĩa vật lí của kết quả.
\item (Đúng) Khi giải cần đọc đúng dữ kiện, dùng hệ đơn vị thống nhất và kiểm tra điều kiện áp dụng.
\end{itemchoice}
}
\end{ex}

% ===== Câu 54 =====
% ID: L12C4B24-14-TL
% Mức: TH
\begin{bt}
% ID: L12C4B24-14-TL
% Mức: TH
Giải thích vì sao nhận định sau đúng: “Nhiệt hạch là sự kết hợp các hạt nhân nhẹ, cần nhiệt độ rất cao để thắng lực đẩy Coulomb và có thể tỏa năng lượng lớn.”
\loigiai{
Cần nêu được: Nhiệt hạch là sự kết hợp các hạt nhân nhẹ, cần nhiệt độ rất cao để thắng lực đẩy Coulomb và có thể tỏa năng lượng lớn. Khi vận dụng phải xác định đúng dữ kiện, điều kiện áp dụng, hệ đơn vị và kiểm tra tính hợp lí của kết quả. Nhận định “Nhiệt hạch dễ xảy ra ở nhiệt độ phòng.” sai vì trái với kiến thức trên.
}
\end{bt}

% ===== Câu 55 =====
% ID: L12C4B24-14-TLN
% Mức: TH
\begin{ex}
% ID: L12C4B24-14-TLN
% Mức: TH
Trong bốn phát biểu sau về nhận biết phản ứng nhiệt hạch và điều kiện xảy ra, có bao nhiêu phát biểu đúng? Chỉ ghi một số nguyên. (1) Nhiệt hạch dễ xảy ra ở nhiệt độ phòng. (2) Nhiệt hạch là sự kết hợp các hạt nhân nhẹ, cần nhiệt độ rất cao để thắng lực đẩy Coulomb và có thể tỏa năng lượng lớn. (3) Kết quả phải phù hợp ý nghĩa vật lí. (4) Mọi đại lượng đều có cùng bản chất.
\shortans{2}
\loigiai{
Đối chiếu từng phát biểu với kiến thức cốt lõi: Nhiệt hạch là sự kết hợp các hạt nhân nhẹ, cần nhiệt độ rất cao để thắng lực đẩy Coulomb và có thể tỏa năng lượng lớn. Có 2 phát biểu đúng.
}
\end{ex}

% ===== Câu 56 =====
% ID: L12C4B24-14-TN
% Mức: TH
\dangbt{Nhận biết phản ứng phân hạch và phản ứng dây chuyền}
\begin{ex}
% ID: L12C4B24-14-TN
% Mức: TH
Nội dung nào mô tả đúng nhất về nhận biết phản ứng phân hạch và phản ứng dây chuyền?
\choice
{Phân hạch là sự kết hợp hai hạt nhân nhẹ thành hạt nhân nặng.}
{Có thể bỏ qua điều kiện áp dụng và đơn vị trong mọi trường hợp.}
{Kết quả của bài toán không phụ thuộc dữ kiện đã cho.}
{\True Phân hạch là sự vỡ của hạt nhân nặng thành các hạt nhân nhẹ hơn, thường kèm neutron và năng lượng; neutron có thể duy trì phản ứng dây chuyền.}
\loigiai{
Phân hạch là sự vỡ của hạt nhân nặng thành các hạt nhân nhẹ hơn, thường kèm neutron và năng lượng; neutron có thể duy trì phản ứng dây chuyền. Vì vậy phương án $D$ đúng; các phương án còn lại trái với kiến thức hoặc điều kiện áp dụng.
}
\end{ex}

% ===== Câu 57 =====
% ID: L12C4B24-15-DS
% Mức: VD
\dangbt{Nhận biết phản ứng nhiệt hạch và điều kiện xảy ra}
\begin{ex}
% ID: L12C4B24-15-DS
% Mức: VD
Xét nhận biết phản ứng nhiệt hạch và điều kiện xảy ra (tình huống 5). Hãy xác định tính đúng, sai của bốn phát biểu sau.
\choiceTF
{\True Nhiệt hạch là sự kết hợp các hạt nhân nhẹ, cần nhiệt độ rất cao để thắng lực đẩy Coulomb và có thể tỏa năng lượng lớn.}
{Có thể bỏ qua dữ kiện, đơn vị và ý nghĩa vật lí của kết quả.}
{Nhiệt hạch dễ xảy ra ở nhiệt độ phòng.}
{Có thể bỏ qua dữ kiện, đơn vị và ý nghĩa vật lí của kết quả.}
\loigiai{
\begin{itemchoice}
\item (Đúng) Nhiệt hạch là sự kết hợp các hạt nhân nhẹ, cần nhiệt độ rất cao để thắng lực đẩy Coulomb và có thể tỏa năng lượng lớn. (Vì): Nhiệt hạch là sự kết hợp các hạt nhân nhẹ, cần nhiệt độ rất cao để thắng lực đẩy Coulomb và có thể tỏa năng lượng lớn. b) Sai: Có thể bỏ qua dữ kiện, đơn vị và ý nghĩa vật lí của kết quả. c) Sai: Nhiệt hạch dễ xảy ra ở nhiệt độ phòng. d) Sai: Có thể bỏ qua dữ kiện, đơn vị và ý nghĩa vật lí của kết quả. Kiến thức cốt lõi: Nhiệt hạch là sự kết hợp các hạt nhân nhẹ, cần nhiệt độ rất cao để thắng lực đẩy Coulomb và có thể tỏa năng lượng lớn
\item (Sai) Có thể bỏ qua dữ kiện, đơn vị và ý nghĩa vật lí của kết quả.
\item (Sai) Nhiệt hạch dễ xảy ra ở nhiệt độ phòng.
\item (Sai) Có thể bỏ qua dữ kiện, đơn vị và ý nghĩa vật lí của kết quả.
\end{itemchoice}
}
\end{ex}

% ===== Câu 58 =====
% ID: L12C4B24-15-TL
% Mức: VD
\begin{bt}
% ID: L12C4B24-15-TL
% Mức: VD
Phân tích sai lầm trong nhận định: “Nhiệt hạch dễ xảy ra ở nhiệt độ phòng.”
\loigiai{
Cần nêu được: Nhiệt hạch là sự kết hợp các hạt nhân nhẹ, cần nhiệt độ rất cao để thắng lực đẩy Coulomb và có thể tỏa năng lượng lớn. Khi vận dụng phải xác định đúng dữ kiện, điều kiện áp dụng, hệ đơn vị và kiểm tra tính hợp lí của kết quả. Nhận định “Nhiệt hạch dễ xảy ra ở nhiệt độ phòng.” sai vì trái với kiến thức trên.
}
\end{bt}

% ===== Câu 59 =====
% ID: L12C4B24-15-TLN
% Mức: TH
\begin{ex}
% ID: L12C4B24-15-TLN
% Mức: TH
Trong bốn phát biểu sau về nhận biết phản ứng nhiệt hạch và điều kiện xảy ra, có bao nhiêu phát biểu đúng? Chỉ ghi một số nguyên. (1) Cần đổi các đại lượng về hệ đơn vị thống nhất. (2) Nhiệt hạch là sự kết hợp các hạt nhân nhẹ, cần nhiệt độ rất cao để thắng lực đẩy Coulomb và có thể tỏa năng lượng lớn. (3) Nhiệt hạch dễ xảy ra ở nhiệt độ phòng. (4) Không cần kiểm tra dữ kiện.
\shortans{2}
\loigiai{
Đối chiếu từng phát biểu với kiến thức cốt lõi: Nhiệt hạch là sự kết hợp các hạt nhân nhẹ, cần nhiệt độ rất cao để thắng lực đẩy Coulomb và có thể tỏa năng lượng lớn. Có 2 phát biểu đúng.
}
\end{ex}

% ===== Câu 60 =====
% ID: L12C4B24-15-TN
% Mức: TH
\dangbt{Nhận biết phản ứng phân hạch và phản ứng dây chuyền}
\begin{ex}
% ID: L12C4B24-15-TN
% Mức: TH
Khi phân tích, phát biểu nào đúng về nhận biết phản ứng phân hạch và phản ứng dây chuyền?
\choice
{\True Phân hạch là sự vỡ của hạt nhân nặng thành các hạt nhân nhẹ hơn, thường kèm neutron và năng lượng; neutron có thể duy trì phản ứng dây chuyền.}
{Phân hạch là sự kết hợp hai hạt nhân nhẹ thành hạt nhân nặng.}
{Có thể bỏ qua điều kiện áp dụng và đơn vị trong mọi trường hợp.}
{Kết quả của bài toán không phụ thuộc dữ kiện đã cho.}
\loigiai{
Phân hạch là sự vỡ của hạt nhân nặng thành các hạt nhân nhẹ hơn, thường kèm neutron và năng lượng; neutron có thể duy trì phản ứng dây chuyền. Vì vậy phương án $A$ đúng; các phương án còn lại trái với kiến thức hoặc điều kiện áp dụng.
}
\end{ex}

% ===== Câu 61 =====
% ID: L12C4B24-16-DS
% Mức: TH
\dangbt{Tính năng lượng trong phân hạch và nhiệt hạch}
\begin{ex}
% ID: L12C4B24-16-DS
% Mức: TH
Xét tính năng lượng trong phân hạch và nhiệt hạch (tình huống 1). Hãy xác định tính đúng, sai của bốn phát biểu sau.
\choiceTF
{\True Năng lượng toàn phần bằng năng lượng mỗi phản ứng nhân số phản ứng; số phản ứng được suy từ số hạt tham gia.}
{Năng lượng toàn phần không phụ thuộc số hạt đã phản ứng.}
{\True Khi giải cần đọc đúng dữ kiện, dùng hệ đơn vị thống nhất và kiểm tra điều kiện áp dụng.}
{Có thể bỏ qua dữ kiện, đơn vị và ý nghĩa vật lí của kết quả.}
\loigiai{
\begin{itemchoice}
\item (Đúng) Năng lượng toàn phần bằng năng lượng mỗi phản ứng nhân số phản ứng; số phản ứng được suy từ số hạt tham gia. (Vì): Năng lượng toàn phần bằng năng lượng mỗi phản ứng nhân số phản ứng; số phản ứng được suy từ số hạt tham gia. b) Sai: Năng lượng toàn phần không phụ thuộc số hạt đã phản ứng. c) Đúng: Khi giải cần đọc đúng dữ kiện, dùng hệ đơn vị thống nhất và kiểm tra điều kiện áp dụng. d) Sai: Có thể bỏ qua dữ kiện, đơn vị và ý nghĩa vật lí của kết quả. Kiến thức cốt lõi: Năng lượng toàn phần bằng năng lượng mỗi phản ứng nhân số phản ứng; số phản ứng được suy từ số hạt tham gia
\item (Sai) Năng lượng toàn phần không phụ thuộc số hạt đã phản ứng.
\item (Đúng) Khi giải cần đọc đúng dữ kiện, dùng hệ đơn vị thống nhất và kiểm tra điều kiện áp dụng.
\item (Sai) Có thể bỏ qua dữ kiện, đơn vị và ý nghĩa vật lí của kết quả.
\end{itemchoice}
}
\end{ex}

% ===== Câu 62 =====
% ID: L12C4B24-16-TL
% Mức: VD
\dangbt{Nhận biết phản ứng nhiệt hạch và điều kiện xảy ra}
\begin{bt}
% ID: L12C4B24-16-TL
% Mức: VD
Đề xuất các bước giải một bài toán thuộc dạng nhận biết phản ứng nhiệt hạch và điều kiện xảy ra.
\loigiai{
Cần nêu được: Nhiệt hạch là sự kết hợp các hạt nhân nhẹ, cần nhiệt độ rất cao để thắng lực đẩy Coulomb và có thể tỏa năng lượng lớn. Khi vận dụng phải xác định đúng dữ kiện, điều kiện áp dụng, hệ đơn vị và kiểm tra tính hợp lí của kết quả. Nhận định “Nhiệt hạch dễ xảy ra ở nhiệt độ phòng.” sai vì trái với kiến thức trên.
}
\end{bt}

% ===== Câu 63 =====
% ID: L12C4B24-16-TLN
% Mức: VD
\begin{ex}
% ID: L12C4B24-16-TLN
% Mức: VD
Trong bốn phát biểu sau về nhận biết phản ứng nhiệt hạch và điều kiện xảy ra, có bao nhiêu phát biểu đúng? Chỉ ghi một số nguyên. (1) Nhiệt hạch là sự kết hợp các hạt nhân nhẹ, cần nhiệt độ rất cao để thắng lực đẩy Coulomb và có thể tỏa năng lượng lớn. (2) Cần đối chiếu kết quả với điều kiện bài toán. (3) Nhiệt hạch dễ xảy ra ở nhiệt độ phòng. (4) Mọi trường hợp đều cho cùng một kết quả.
\shortans{2}
\loigiai{
Đối chiếu từng phát biểu với kiến thức cốt lõi: Nhiệt hạch là sự kết hợp các hạt nhân nhẹ, cần nhiệt độ rất cao để thắng lực đẩy Coulomb và có thể tỏa năng lượng lớn. Có 2 phát biểu đúng.
}
\end{ex}

% ===== Câu 64 =====
% ID: L12C4B24-16-TN
% Mức: TH
\dangbt{Nhận biết phản ứng phân hạch và phản ứng dây chuyền}
\begin{ex}
% ID: L12C4B24-16-TN
% Mức: TH
Kết luận nào của học sinh là đúng về nhận biết phản ứng phân hạch và phản ứng dây chuyền?
\choice
{Phân hạch là sự kết hợp hai hạt nhân nhẹ thành hạt nhân nặng.}
{\True Phân hạch là sự vỡ của hạt nhân nặng thành các hạt nhân nhẹ hơn, thường kèm neutron và năng lượng; neutron có thể duy trì phản ứng dây chuyền.}
{Có thể bỏ qua điều kiện áp dụng và đơn vị trong mọi trường hợp.}
{Kết quả của bài toán không phụ thuộc dữ kiện đã cho.}
\loigiai{
Phân hạch là sự vỡ của hạt nhân nặng thành các hạt nhân nhẹ hơn, thường kèm neutron và năng lượng; neutron có thể duy trì phản ứng dây chuyền. Vì vậy phương án $B$ đúng; các phương án còn lại trái với kiến thức hoặc điều kiện áp dụng.
}
\end{ex}

% ===== Câu 65 =====
% ID: L12C4B24-17-DS
% Mức: TH
\dangbt{Tính năng lượng trong phân hạch và nhiệt hạch}
\begin{ex}
% ID: L12C4B24-17-DS
% Mức: TH
Xét tính năng lượng trong phân hạch và nhiệt hạch (tình huống 2). Hãy xác định tính đúng, sai của bốn phát biểu sau.
\choiceTF
{Năng lượng toàn phần không phụ thuộc số hạt đã phản ứng.}
{\True Năng lượng toàn phần bằng năng lượng mỗi phản ứng nhân số phản ứng; số phản ứng được suy từ số hạt tham gia.}
{Có thể bỏ qua dữ kiện, đơn vị và ý nghĩa vật lí của kết quả.}
{\True Khi giải cần đọc đúng dữ kiện, dùng hệ đơn vị thống nhất và kiểm tra điều kiện áp dụng.}
\loigiai{
\begin{itemchoice}
\item (Sai) Năng lượng toàn phần không phụ thuộc số hạt đã phản ứng. (Vì): Năng lượng toàn phần không phụ thuộc số hạt đã phản ứng. b) Đúng: Năng lượng toàn phần bằng năng lượng mỗi phản ứng nhân số phản ứng; số phản ứng được suy từ số hạt tham gia. c) Sai: Có thể bỏ qua dữ kiện, đơn vị và ý nghĩa vật lí của kết quả. d) Đúng: Khi giải cần đọc đúng dữ kiện, dùng hệ đơn vị thống nhất và kiểm tra điều kiện áp dụng. Kiến thức cốt lõi: Năng lượng toàn phần bằng năng lượng mỗi phản ứng nhân số phản ứng; số phản ứng được suy từ số hạt tham gia
\item (Đúng) Năng lượng toàn phần bằng năng lượng mỗi phản ứng nhân số phản ứng; số phản ứng được suy từ số hạt tham gia.
\item (Sai) Có thể bỏ qua dữ kiện, đơn vị và ý nghĩa vật lí của kết quả.
\item (Đúng) Khi giải cần đọc đúng dữ kiện, dùng hệ đơn vị thống nhất và kiểm tra điều kiện áp dụng.
\end{itemchoice}
}
\end{ex}

% ===== Câu 66 =====
% ID: L12C4B24-17-TL
% Mức: VD
\dangbt{Nhận biết phản ứng nhiệt hạch và điều kiện xảy ra}
\begin{bt}
% ID: L12C4B24-17-TL
% Mức: VD
Nêu một tình huống thực tiễn liên quan đến nhận biết phản ứng nhiệt hạch và điều kiện xảy ra và giải thích bằng kiến thức Vật lí.
\loigiai{
Cần nêu được: Nhiệt hạch là sự kết hợp các hạt nhân nhẹ, cần nhiệt độ rất cao để thắng lực đẩy Coulomb và có thể tỏa năng lượng lớn. Khi vận dụng phải xác định đúng dữ kiện, điều kiện áp dụng, hệ đơn vị và kiểm tra tính hợp lí của kết quả. Nhận định “Nhiệt hạch dễ xảy ra ở nhiệt độ phòng.” sai vì trái với kiến thức trên.
}
\end{bt}

% ===== Câu 67 =====
% ID: L12C4B24-17-TLN
% Mức: VD
\begin{ex}
% ID: L12C4B24-17-TLN
% Mức: VD
Trong bốn phát biểu sau về nhận biết phản ứng nhiệt hạch và điều kiện xảy ra, có bao nhiêu phát biểu đúng? Chỉ ghi một số nguyên. (1) Nhiệt hạch dễ xảy ra ở nhiệt độ phòng. (2) Cần đọc đúng dữ kiện và giả thiết. (3) Nhiệt hạch là sự kết hợp các hạt nhân nhẹ, cần nhiệt độ rất cao để thắng lực đẩy Coulomb và có thể tỏa năng lượng lớn. (4) Có thể thay số trước khi chọn công thức.
\shortans{2}
\loigiai{
Đối chiếu từng phát biểu với kiến thức cốt lõi: Nhiệt hạch là sự kết hợp các hạt nhân nhẹ, cần nhiệt độ rất cao để thắng lực đẩy Coulomb và có thể tỏa năng lượng lớn. Có 2 phát biểu đúng.
}
\end{ex}

% ===== Câu 68 =====
% ID: L12C4B24-17-TN
% Mức: TH
\dangbt{Nhận biết phản ứng phân hạch và phản ứng dây chuyền}
\begin{ex}
% ID: L12C4B24-17-TN
% Mức: TH
Chọn phát biểu không trái với kiến thức về nhận biết phản ứng phân hạch và phản ứng dây chuyền?
\choice
{Phân hạch là sự kết hợp hai hạt nhân nhẹ thành hạt nhân nặng.}
{Có thể bỏ qua điều kiện áp dụng và đơn vị trong mọi trường hợp.}
{\True Phân hạch là sự vỡ của hạt nhân nặng thành các hạt nhân nhẹ hơn, thường kèm neutron và năng lượng; neutron có thể duy trì phản ứng dây chuyền.}
{Kết quả của bài toán không phụ thuộc dữ kiện đã cho.}
\loigiai{
Phân hạch là sự vỡ của hạt nhân nặng thành các hạt nhân nhẹ hơn, thường kèm neutron và năng lượng; neutron có thể duy trì phản ứng dây chuyền. Vì vậy phương án $C$ đúng; các phương án còn lại trái với kiến thức hoặc điều kiện áp dụng.
}
\end{ex}

% ===== Câu 69 =====
% ID: L12C4B24-18-DS
% Mức: VD
\dangbt{Tính năng lượng trong phân hạch và nhiệt hạch}
\begin{ex}
% ID: L12C4B24-18-DS
% Mức: VD
Xét tính năng lượng trong phân hạch và nhiệt hạch (tình huống 3). Hãy xác định tính đúng, sai của bốn phát biểu sau.
\choiceTF
{\True Năng lượng toàn phần bằng năng lượng mỗi phản ứng nhân số phản ứng; số phản ứng được suy từ số hạt tham gia.}
{Năng lượng toàn phần không phụ thuộc số hạt đã phản ứng.}
{\True Khi giải cần đọc đúng dữ kiện, dùng hệ đơn vị thống nhất và kiểm tra điều kiện áp dụng.}
{Có thể bỏ qua dữ kiện, đơn vị và ý nghĩa vật lí của kết quả.}
\loigiai{
\begin{itemchoice}
\item (Đúng) Năng lượng toàn phần bằng năng lượng mỗi phản ứng nhân số phản ứng; số phản ứng được suy từ số hạt tham gia. (Vì): Năng lượng toàn phần bằng năng lượng mỗi phản ứng nhân số phản ứng; số phản ứng được suy từ số hạt tham gia. b) Sai: Năng lượng toàn phần không phụ thuộc số hạt đã phản ứng. c) Đúng: Khi giải cần đọc đúng dữ kiện, dùng hệ đơn vị thống nhất và kiểm tra điều kiện áp dụng. d) Sai: Có thể bỏ qua dữ kiện, đơn vị và ý nghĩa vật lí của kết quả. Kiến thức cốt lõi: Năng lượng toàn phần bằng năng lượng mỗi phản ứng nhân số phản ứng; số phản ứng được suy từ số hạt tham gia
\item (Sai) Năng lượng toàn phần không phụ thuộc số hạt đã phản ứng.
\item (Đúng) Khi giải cần đọc đúng dữ kiện, dùng hệ đơn vị thống nhất và kiểm tra điều kiện áp dụng.
\item (Sai) Có thể bỏ qua dữ kiện, đơn vị và ý nghĩa vật lí của kết quả.
\end{itemchoice}
}
\end{ex}

% ===== Câu 70 =====
% ID: L12C4B24-18-TL
% Mức: VDC
\dangbt{Nhận biết phản ứng nhiệt hạch và điều kiện xảy ra}
\begin{bt}
% ID: L12C4B24-18-TL
% Mức: VDC
So sánh nhận định đúng “Nhiệt hạch là sự kết hợp các hạt nhân nhẹ, cần nhiệt độ rất cao để thắng lực đẩy Coulomb và có thể tỏa năng lượng lớn.” với nhận định sai “Nhiệt hạch dễ xảy ra ở nhiệt độ phòng.”; chỉ rõ căn cứ Vật lí.
\loigiai{
Cần nêu được: Nhiệt hạch là sự kết hợp các hạt nhân nhẹ, cần nhiệt độ rất cao để thắng lực đẩy Coulomb và có thể tỏa năng lượng lớn. Khi vận dụng phải xác định đúng dữ kiện, điều kiện áp dụng, hệ đơn vị và kiểm tra tính hợp lí của kết quả. Nhận định “Nhiệt hạch dễ xảy ra ở nhiệt độ phòng.” sai vì trái với kiến thức trên.
}
\end{bt}

% ===== Câu 71 =====
% ID: L12C4B24-18-TLN
% Mức: VDC
\begin{ex}
% ID: L12C4B24-18-TLN
% Mức: VDC
Trong bốn phát biểu sau về nhận biết phản ứng nhiệt hạch và điều kiện xảy ra, có bao nhiêu phát biểu đúng? Chỉ ghi một số nguyên. (1) Kết quả cần có ý nghĩa vật lí. (2) Nhiệt hạch dễ xảy ra ở nhiệt độ phòng. (3) Phải dùng đúng định luật cho tình huống. (4) Nhiệt hạch là sự kết hợp các hạt nhân nhẹ, cần nhiệt độ rất cao để thắng lực đẩy Coulomb và có thể tỏa năng lượng lớn.
\shortans{3}
\loigiai{
Đối chiếu từng phát biểu với kiến thức cốt lõi: Nhiệt hạch là sự kết hợp các hạt nhân nhẹ, cần nhiệt độ rất cao để thắng lực đẩy Coulomb và có thể tỏa năng lượng lớn. Có 3 phát biểu đúng.
}
\end{ex}

% ===== Câu 72 =====
% ID: L12C4B24-18-TN
% Mức: VD
\dangbt{Nhận biết phản ứng phân hạch và phản ứng dây chuyền}
\begin{ex}
% ID: L12C4B24-18-TN
% Mức: VD
Cơ sở Vật lí đúng để giải bài là gì về nhận biết phản ứng phân hạch và phản ứng dây chuyền?
\choice
{Phân hạch là sự kết hợp hai hạt nhân nhẹ thành hạt nhân nặng.}
{Có thể bỏ qua điều kiện áp dụng và đơn vị trong mọi trường hợp.}
{Kết quả của bài toán không phụ thuộc dữ kiện đã cho.}
{\True Phân hạch là sự vỡ của hạt nhân nặng thành các hạt nhân nhẹ hơn, thường kèm neutron và năng lượng; neutron có thể duy trì phản ứng dây chuyền.}
\loigiai{
Phân hạch là sự vỡ của hạt nhân nặng thành các hạt nhân nhẹ hơn, thường kèm neutron và năng lượng; neutron có thể duy trì phản ứng dây chuyền. Vì vậy phương án $D$ đúng; các phương án còn lại trái với kiến thức hoặc điều kiện áp dụng.
}
\end{ex}

% ===== Câu 73 =====
% ID: L12C4B24-19-DS
% Mức: VD
\dangbt{Tính năng lượng trong phân hạch và nhiệt hạch}
\begin{ex}
% ID: L12C4B24-19-DS
% Mức: VD
Xét tính năng lượng trong phân hạch và nhiệt hạch (tình huống 4). Hãy xác định tính đúng, sai của bốn phát biểu sau.
\choiceTF
{Năng lượng toàn phần không phụ thuộc số hạt đã phản ứng.}
{\True Năng lượng toàn phần bằng năng lượng mỗi phản ứng nhân số phản ứng; số phản ứng được suy từ số hạt tham gia.}
{Có thể bỏ qua dữ kiện, đơn vị và ý nghĩa vật lí của kết quả.}
{\True Khi giải cần đọc đúng dữ kiện, dùng hệ đơn vị thống nhất và kiểm tra điều kiện áp dụng.}
\loigiai{
\begin{itemchoice}
\item (Sai) Năng lượng toàn phần không phụ thuộc số hạt đã phản ứng. (Vì): Năng lượng toàn phần không phụ thuộc số hạt đã phản ứng. b) Đúng: Năng lượng toàn phần bằng năng lượng mỗi phản ứng nhân số phản ứng; số phản ứng được suy từ số hạt tham gia. c) Sai: Có thể bỏ qua dữ kiện, đơn vị và ý nghĩa vật lí của kết quả. d) Đúng: Khi giải cần đọc đúng dữ kiện, dùng hệ đơn vị thống nhất và kiểm tra điều kiện áp dụng. Kiến thức cốt lõi: Năng lượng toàn phần bằng năng lượng mỗi phản ứng nhân số phản ứng; số phản ứng được suy từ số hạt tham gia
\item (Đúng) Năng lượng toàn phần bằng năng lượng mỗi phản ứng nhân số phản ứng; số phản ứng được suy từ số hạt tham gia.
\item (Sai) Có thể bỏ qua dữ kiện, đơn vị và ý nghĩa vật lí của kết quả.
\item (Đúng) Khi giải cần đọc đúng dữ kiện, dùng hệ đơn vị thống nhất và kiểm tra điều kiện áp dụng.
\end{itemchoice}
}
\end{ex}

% ===== Câu 74 =====
% ID: L12C4B24-19-TL
% Mức: TH
\begin{bt}
% ID: L12C4B24-19-TL
% Mức: TH
Trình bày kiến thức cốt lõi của tính năng lượng trong phân hạch và nhiệt hạch và nêu điều kiện áp dụng.
\loigiai{
Cần nêu được: Năng lượng toàn phần bằng năng lượng mỗi phản ứng nhân số phản ứng; số phản ứng được suy từ số hạt tham gia. Khi vận dụng phải xác định đúng dữ kiện, điều kiện áp dụng, hệ đơn vị và kiểm tra tính hợp lí của kết quả. Nhận định “Năng lượng toàn phần không phụ thuộc số hạt đã phản ứng.” sai vì trái với kiến thức trên.
}
\end{bt}

% ===== Câu 75 =====
% ID: L12C4B24-19-TLN
% Mức: TH
\begin{ex}
% ID: L12C4B24-19-TLN
% Mức: TH
Trong bốn phát biểu sau về tính năng lượng trong phân hạch và nhiệt hạch, có bao nhiêu phát biểu đúng? Chỉ ghi một số nguyên. (1) Năng lượng toàn phần bằng năng lượng mỗi phản ứng nhân số phản ứng; số phản ứng được suy từ số hạt tham gia. (2) Năng lượng toàn phần không phụ thuộc số hạt đã phản ứng. (3) Cần kiểm tra điều kiện áp dụng trước khi tính toán. (4) Có thể bỏ qua đơn vị của kết quả.
\shortans{2}
\loigiai{
Đối chiếu từng phát biểu với kiến thức cốt lõi: Năng lượng toàn phần bằng năng lượng mỗi phản ứng nhân số phản ứng; số phản ứng được suy từ số hạt tham gia. Có 2 phát biểu đúng.
}
\end{ex}

% ===== Câu 76 =====
% ID: L12C4B24-19-TN
% Mức: VD
\dangbt{Nhận biết phản ứng phân hạch và phản ứng dây chuyền}
\begin{ex}
% ID: L12C4B24-19-TN
% Mức: VD
Trong một bài thực hành, nhận xét nào đúng về nhận biết phản ứng phân hạch và phản ứng dây chuyền?
\choice
{\True Phân hạch là sự vỡ của hạt nhân nặng thành các hạt nhân nhẹ hơn, thường kèm neutron và năng lượng; neutron có thể duy trì phản ứng dây chuyền.}
{Phân hạch là sự kết hợp hai hạt nhân nhẹ thành hạt nhân nặng.}
{Có thể bỏ qua điều kiện áp dụng và đơn vị trong mọi trường hợp.}
{Kết quả của bài toán không phụ thuộc dữ kiện đã cho.}
\loigiai{
Phân hạch là sự vỡ của hạt nhân nặng thành các hạt nhân nhẹ hơn, thường kèm neutron và năng lượng; neutron có thể duy trì phản ứng dây chuyền. Vì vậy phương án $A$ đúng; các phương án còn lại trái với kiến thức hoặc điều kiện áp dụng.
}
\end{ex}

% ===== Câu 77 =====
% ID: L12C4B24-20-DS
% Mức: VD
\dangbt{Tính năng lượng trong phân hạch và nhiệt hạch}
\begin{ex}
% ID: L12C4B24-20-DS
% Mức: VD
Xét tính năng lượng trong phân hạch và nhiệt hạch (tình huống 5). Hãy xác định tính đúng, sai của bốn phát biểu sau.
\choiceTF
{\True Năng lượng toàn phần bằng năng lượng mỗi phản ứng nhân số phản ứng; số phản ứng được suy từ số hạt tham gia.}
{Có thể bỏ qua dữ kiện, đơn vị và ý nghĩa vật lí của kết quả.}
{Năng lượng toàn phần không phụ thuộc số hạt đã phản ứng.}
{Có thể bỏ qua dữ kiện, đơn vị và ý nghĩa vật lí của kết quả.}
\loigiai{
\begin{itemchoice}
\item (Đúng) Năng lượng toàn phần bằng năng lượng mỗi phản ứng nhân số phản ứng; số phản ứng được suy từ số hạt tham gia. (Vì): Năng lượng toàn phần bằng năng lượng mỗi phản ứng nhân số phản ứng; số phản ứng được suy từ số hạt tham gia. b) Sai: Có thể bỏ qua dữ kiện, đơn vị và ý nghĩa vật lí của kết quả. c) Sai: Năng lượng toàn phần không phụ thuộc số hạt đã phản ứng. d) Sai: Có thể bỏ qua dữ kiện, đơn vị và ý nghĩa vật lí của kết quả. Kiến thức cốt lõi: Năng lượng toàn phần bằng năng lượng mỗi phản ứng nhân số phản ứng; số phản ứng được suy từ số hạt tham gia
\item (Sai) Có thể bỏ qua dữ kiện, đơn vị và ý nghĩa vật lí của kết quả.
\item (Sai) Năng lượng toàn phần không phụ thuộc số hạt đã phản ứng.
\item (Sai) Có thể bỏ qua dữ kiện, đơn vị và ý nghĩa vật lí của kết quả.
\end{itemchoice}
}
\end{ex}

% ===== Câu 78 =====
% ID: L12C4B24-20-TL
% Mức: TH
\begin{bt}
% ID: L12C4B24-20-TL
% Mức: TH
Giải thích vì sao nhận định sau đúng: “Năng lượng toàn phần bằng năng lượng mỗi phản ứng nhân số phản ứng; số phản ứng được suy từ số hạt tham gia.”
\loigiai{
Cần nêu được: Năng lượng toàn phần bằng năng lượng mỗi phản ứng nhân số phản ứng; số phản ứng được suy từ số hạt tham gia. Khi vận dụng phải xác định đúng dữ kiện, điều kiện áp dụng, hệ đơn vị và kiểm tra tính hợp lí của kết quả. Nhận định “Năng lượng toàn phần không phụ thuộc số hạt đã phản ứng.” sai vì trái với kiến thức trên.
}
\end{bt}

% ===== Câu 79 =====
% ID: L12C4B24-20-TLN
% Mức: TH
\begin{ex}
% ID: L12C4B24-20-TLN
% Mức: TH
Trong bốn phát biểu sau về tính năng lượng trong phân hạch và nhiệt hạch, có bao nhiêu phát biểu đúng? Chỉ ghi một số nguyên. (1) Năng lượng toàn phần không phụ thuộc số hạt đã phản ứng. (2) Năng lượng toàn phần bằng năng lượng mỗi phản ứng nhân số phản ứng; số phản ứng được suy từ số hạt tham gia. (3) Kết quả phải phù hợp ý nghĩa vật lí. (4) Mọi đại lượng đều có cùng bản chất.
\shortans{2}
\loigiai{
Đối chiếu từng phát biểu với kiến thức cốt lõi: Năng lượng toàn phần bằng năng lượng mỗi phản ứng nhân số phản ứng; số phản ứng được suy từ số hạt tham gia. Có 2 phát biểu đúng.
}
\end{ex}

% ===== Câu 80 =====
% ID: L12C4B24-20-TN
% Mức: VD
\dangbt{Nhận biết phản ứng phân hạch và phản ứng dây chuyền}
\begin{ex}
% ID: L12C4B24-20-TN
% Mức: VD
Kết luận đúng khi vận dụng là về nhận biết phản ứng phân hạch và phản ứng dây chuyền?
\choice
{Phân hạch là sự kết hợp hai hạt nhân nhẹ thành hạt nhân nặng.}
{\True Phân hạch là sự vỡ của hạt nhân nặng thành các hạt nhân nhẹ hơn, thường kèm neutron và năng lượng; neutron có thể duy trì phản ứng dây chuyền.}
{Có thể bỏ qua điều kiện áp dụng và đơn vị trong mọi trường hợp.}
{Kết quả của bài toán không phụ thuộc dữ kiện đã cho.}
\loigiai{
Phân hạch là sự vỡ của hạt nhân nặng thành các hạt nhân nhẹ hơn, thường kèm neutron và năng lượng; neutron có thể duy trì phản ứng dây chuyền. Vì vậy phương án $B$ đúng; các phương án còn lại trái với kiến thức hoặc điều kiện áp dụng.
}
\end{ex}

% ===== Câu 81 =====
% ID: L12C4B24-21-DS
% Mức: TH
\dangbt{Cấu tạo và nguyên lí hoạt động của lò phản ứng hạt nhân}
\begin{ex}
% ID: L12C4B24-21-DS
% Mức: TH
Xét cấu tạo và nguyên lí hoạt động của lò phản ứng hạt nhân (tình huống 1). Hãy xác định tính đúng, sai của bốn phát biểu sau.
\choiceTF
{\True Lò phản ứng dùng nhiên liệu phân hạch, chất làm chậm, thanh điều khiển, chất tải nhiệt và hệ che chắn để duy trì phản ứng có kiểm soát.}
{Thanh điều khiển có nhiệm vụ làm tăng không giới hạn số neutron trong lò.}
{\True Khi giải cần đọc đúng dữ kiện, dùng hệ đơn vị thống nhất và kiểm tra điều kiện áp dụng.}
{Có thể bỏ qua dữ kiện, đơn vị và ý nghĩa vật lí của kết quả.}
\loigiai{
\begin{itemchoice}
\item (Đúng) Lò phản ứng dùng nhiên liệu phân hạch, chất làm chậm, thanh điều khiển, chất tải nhiệt và hệ che chắn để duy trì phản ứng có kiểm soát. (Vì): Lò phản ứng dùng nhiên liệu phân hạch, chất làm chậm, thanh điều khiển, chất tải nhiệt và hệ che chắn để duy trì phản ứng có kiểm soát. b) Sai: Thanh điều khiển có nhiệm vụ làm tăng không giới hạn số neutron trong lò. c) Đúng: Khi giải cần đọc đúng dữ kiện, dùng hệ đơn vị thống nhất và kiểm tra điều kiện áp dụng. d) Sai: Có thể bỏ qua dữ kiện, đơn vị và ý nghĩa vật lí của kết quả. Kiến thức cốt lõi: Lò phản ứng dùng nhiên liệu phân hạch, chất làm chậm, thanh điều khiển, chất tải nhiệt và hệ che chắn để duy trì phản ứng có kiểm soát
\item (Sai) Thanh điều khiển có nhiệm vụ làm tăng không giới hạn số neutron trong lò.
\item (Đúng) Khi giải cần đọc đúng dữ kiện, dùng hệ đơn vị thống nhất và kiểm tra điều kiện áp dụng.
\item (Sai) Có thể bỏ qua dữ kiện, đơn vị và ý nghĩa vật lí của kết quả.
\end{itemchoice}
}
\end{ex}

% ===== Câu 82 =====
% ID: L12C4B24-21-TL
% Mức: VD
\dangbt{Tính năng lượng trong phân hạch và nhiệt hạch}
\begin{bt}
% ID: L12C4B24-21-TL
% Mức: VD
Phân tích sai lầm trong nhận định: “Năng lượng toàn phần không phụ thuộc số hạt đã phản ứng.”
\loigiai{
Cần nêu được: Năng lượng toàn phần bằng năng lượng mỗi phản ứng nhân số phản ứng; số phản ứng được suy từ số hạt tham gia. Khi vận dụng phải xác định đúng dữ kiện, điều kiện áp dụng, hệ đơn vị và kiểm tra tính hợp lí của kết quả. Nhận định “Năng lượng toàn phần không phụ thuộc số hạt đã phản ứng.” sai vì trái với kiến thức trên.
}
\end{bt}

% ===== Câu 83 =====
% ID: L12C4B24-21-TLN
% Mức: TH
\begin{ex}
% ID: L12C4B24-21-TLN
% Mức: TH
Trong bốn phát biểu sau về tính năng lượng trong phân hạch và nhiệt hạch, có bao nhiêu phát biểu đúng? Chỉ ghi một số nguyên. (1) Cần đổi các đại lượng về hệ đơn vị thống nhất. (2) Năng lượng toàn phần bằng năng lượng mỗi phản ứng nhân số phản ứng; số phản ứng được suy từ số hạt tham gia. (3) Năng lượng toàn phần không phụ thuộc số hạt đã phản ứng. (4) Không cần kiểm tra dữ kiện.
\shortans{2}
\loigiai{
Đối chiếu từng phát biểu với kiến thức cốt lõi: Năng lượng toàn phần bằng năng lượng mỗi phản ứng nhân số phản ứng; số phản ứng được suy từ số hạt tham gia. Có 2 phát biểu đúng.
}
\end{ex}

% ===== Câu 84 =====
% ID: L12C4B24-21-TN
% Mức: NB
\dangbt{Nhận biết phản ứng nhiệt hạch và điều kiện xảy ra}
\begin{ex}
% ID: L12C4B24-21-TN
% Mức: NB
Phát biểu nào sau đây đúng về nhận biết phản ứng nhiệt hạch và điều kiện xảy ra?
\choice
{\True Nhiệt hạch là sự kết hợp các hạt nhân nhẹ, cần nhiệt độ rất cao để thắng lực đẩy Coulomb và có thể tỏa năng lượng lớn.}
{Nhiệt hạch dễ xảy ra ở nhiệt độ phòng.}
{Có thể bỏ qua điều kiện áp dụng và đơn vị trong mọi trường hợp.}
{Kết quả của bài toán không phụ thuộc dữ kiện đã cho.}
\loigiai{
Nhiệt hạch là sự kết hợp các hạt nhân nhẹ, cần nhiệt độ rất cao để thắng lực đẩy Coulomb và có thể tỏa năng lượng lớn. Vì vậy phương án $A$ đúng; các phương án còn lại trái với kiến thức hoặc điều kiện áp dụng.
}
\end{ex}

% ===== Câu 85 =====
% ID: L12C4B24-22-DS
% Mức: TH
\dangbt{Cấu tạo và nguyên lí hoạt động của lò phản ứng hạt nhân}
\begin{ex}
% ID: L12C4B24-22-DS
% Mức: TH
Xét cấu tạo và nguyên lí hoạt động của lò phản ứng hạt nhân (tình huống 2). Hãy xác định tính đúng, sai của bốn phát biểu sau.
\choiceTF
{Thanh điều khiển có nhiệm vụ làm tăng không giới hạn số neutron trong lò.}
{\True Lò phản ứng dùng nhiên liệu phân hạch, chất làm chậm, thanh điều khiển, chất tải nhiệt và hệ che chắn để duy trì phản ứng có kiểm soát.}
{Có thể bỏ qua dữ kiện, đơn vị và ý nghĩa vật lí của kết quả.}
{\True Khi giải cần đọc đúng dữ kiện, dùng hệ đơn vị thống nhất và kiểm tra điều kiện áp dụng.}
\loigiai{
\begin{itemchoice}
\item (Sai) Thanh điều khiển có nhiệm vụ làm tăng không giới hạn số neutron trong lò. (Vì): Thanh điều khiển có nhiệm vụ làm tăng không giới hạn số neutron trong lò. b) Đúng: Lò phản ứng dùng nhiên liệu phân hạch, chất làm chậm, thanh điều khiển, chất tải nhiệt và hệ che chắn để duy trì phản ứng có kiểm soát. c) Sai: Có thể bỏ qua dữ kiện, đơn vị và ý nghĩa vật lí của kết quả. d) Đúng: Khi giải cần đọc đúng dữ kiện, dùng hệ đơn vị thống nhất và kiểm tra điều kiện áp dụng. Kiến thức cốt lõi: Lò phản ứng dùng nhiên liệu phân hạch, chất làm chậm, thanh điều khiển, chất tải nhiệt và hệ che chắn để duy trì phản ứng có kiểm soát
\item (Đúng) Lò phản ứng dùng nhiên liệu phân hạch, chất làm chậm, thanh điều khiển, chất tải nhiệt và hệ che chắn để duy trì phản ứng có kiểm soát.
\item (Sai) Có thể bỏ qua dữ kiện, đơn vị và ý nghĩa vật lí của kết quả.
\item (Đúng) Khi giải cần đọc đúng dữ kiện, dùng hệ đơn vị thống nhất và kiểm tra điều kiện áp dụng.
\end{itemchoice}
}
\end{ex}

% ===== Câu 86 =====
% ID: L12C4B24-22-TL
% Mức: VD
\dangbt{Tính năng lượng trong phân hạch và nhiệt hạch}
\begin{bt}
% ID: L12C4B24-22-TL
% Mức: VD
Đề xuất các bước giải một bài toán thuộc dạng tính năng lượng trong phân hạch và nhiệt hạch.
\loigiai{
Cần nêu được: Năng lượng toàn phần bằng năng lượng mỗi phản ứng nhân số phản ứng; số phản ứng được suy từ số hạt tham gia. Khi vận dụng phải xác định đúng dữ kiện, điều kiện áp dụng, hệ đơn vị và kiểm tra tính hợp lí của kết quả. Nhận định “Năng lượng toàn phần không phụ thuộc số hạt đã phản ứng.” sai vì trái với kiến thức trên.
}
\end{bt}

% ===== Câu 87 =====
% ID: L12C4B24-22-TLN
% Mức: VD
\begin{ex}
% ID: L12C4B24-22-TLN
% Mức: VD
Trong bốn phát biểu sau về tính năng lượng trong phân hạch và nhiệt hạch, có bao nhiêu phát biểu đúng? Chỉ ghi một số nguyên. (1) Năng lượng toàn phần bằng năng lượng mỗi phản ứng nhân số phản ứng; số phản ứng được suy từ số hạt tham gia. (2) Cần đối chiếu kết quả với điều kiện bài toán. (3) Năng lượng toàn phần không phụ thuộc số hạt đã phản ứng. (4) Mọi trường hợp đều cho cùng một kết quả.
\shortans{2}
\loigiai{
Đối chiếu từng phát biểu với kiến thức cốt lõi: Năng lượng toàn phần bằng năng lượng mỗi phản ứng nhân số phản ứng; số phản ứng được suy từ số hạt tham gia. Có 2 phát biểu đúng.
}
\end{ex}

% ===== Câu 88 =====
% ID: L12C4B24-22-TN
% Mức: NB
\dangbt{Nhận biết phản ứng nhiệt hạch và điều kiện xảy ra}
\begin{ex}
% ID: L12C4B24-22-TN
% Mức: NB
Nhận định nào phù hợp với kiến thức Vật lí về nhận biết phản ứng nhiệt hạch và điều kiện xảy ra?
\choice
{Nhiệt hạch dễ xảy ra ở nhiệt độ phòng.}
{\True Nhiệt hạch là sự kết hợp các hạt nhân nhẹ, cần nhiệt độ rất cao để thắng lực đẩy Coulomb và có thể tỏa năng lượng lớn.}
{Có thể bỏ qua điều kiện áp dụng và đơn vị trong mọi trường hợp.}
{Kết quả của bài toán không phụ thuộc dữ kiện đã cho.}
\loigiai{
Nhiệt hạch là sự kết hợp các hạt nhân nhẹ, cần nhiệt độ rất cao để thắng lực đẩy Coulomb và có thể tỏa năng lượng lớn. Vì vậy phương án $B$ đúng; các phương án còn lại trái với kiến thức hoặc điều kiện áp dụng.
}
\end{ex}

% ===== Câu 89 =====
% ID: L12C4B24-23-DS
% Mức: VD
\dangbt{Cấu tạo và nguyên lí hoạt động của lò phản ứng hạt nhân}
\begin{ex}
% ID: L12C4B24-23-DS
% Mức: VD
Xét cấu tạo và nguyên lí hoạt động của lò phản ứng hạt nhân (tình huống 3). Hãy xác định tính đúng, sai của bốn phát biểu sau.
\choiceTF
{\True Lò phản ứng dùng nhiên liệu phân hạch, chất làm chậm, thanh điều khiển, chất tải nhiệt và hệ che chắn để duy trì phản ứng có kiểm soát.}
{Thanh điều khiển có nhiệm vụ làm tăng không giới hạn số neutron trong lò.}
{\True Khi giải cần đọc đúng dữ kiện, dùng hệ đơn vị thống nhất và kiểm tra điều kiện áp dụng.}
{Có thể bỏ qua dữ kiện, đơn vị và ý nghĩa vật lí của kết quả.}
\loigiai{
\begin{itemchoice}
\item (Đúng) Lò phản ứng dùng nhiên liệu phân hạch, chất làm chậm, thanh điều khiển, chất tải nhiệt và hệ che chắn để duy trì phản ứng có kiểm soát. (Vì): Lò phản ứng dùng nhiên liệu phân hạch, chất làm chậm, thanh điều khiển, chất tải nhiệt và hệ che chắn để duy trì phản ứng có kiểm soát. b) Sai: Thanh điều khiển có nhiệm vụ làm tăng không giới hạn số neutron trong lò. c) Đúng: Khi giải cần đọc đúng dữ kiện, dùng hệ đơn vị thống nhất và kiểm tra điều kiện áp dụng. d) Sai: Có thể bỏ qua dữ kiện, đơn vị và ý nghĩa vật lí của kết quả. Kiến thức cốt lõi: Lò phản ứng dùng nhiên liệu phân hạch, chất làm chậm, thanh điều khiển, chất tải nhiệt và hệ che chắn để duy trì phản ứng có kiểm soát
\item (Sai) Thanh điều khiển có nhiệm vụ làm tăng không giới hạn số neutron trong lò.
\item (Đúng) Khi giải cần đọc đúng dữ kiện, dùng hệ đơn vị thống nhất và kiểm tra điều kiện áp dụng.
\item (Sai) Có thể bỏ qua dữ kiện, đơn vị và ý nghĩa vật lí của kết quả.
\end{itemchoice}
}
\end{ex}

% ===== Câu 90 =====
% ID: L12C4B24-23-TL
% Mức: VD
\dangbt{Tính năng lượng trong phân hạch và nhiệt hạch}
\begin{bt}
% ID: L12C4B24-23-TL
% Mức: VD
Nêu một tình huống thực tiễn liên quan đến tính năng lượng trong phân hạch và nhiệt hạch và giải thích bằng kiến thức Vật lí.
\loigiai{
Cần nêu được: Năng lượng toàn phần bằng năng lượng mỗi phản ứng nhân số phản ứng; số phản ứng được suy từ số hạt tham gia. Khi vận dụng phải xác định đúng dữ kiện, điều kiện áp dụng, hệ đơn vị và kiểm tra tính hợp lí của kết quả. Nhận định “Năng lượng toàn phần không phụ thuộc số hạt đã phản ứng.” sai vì trái với kiến thức trên.
}
\end{bt}

% ===== Câu 91 =====
% ID: L12C4B24-23-TLN
% Mức: VD
\begin{ex}
% ID: L12C4B24-23-TLN
% Mức: VD
Trong bốn phát biểu sau về tính năng lượng trong phân hạch và nhiệt hạch, có bao nhiêu phát biểu đúng? Chỉ ghi một số nguyên. (1) Năng lượng toàn phần không phụ thuộc số hạt đã phản ứng. (2) Cần đọc đúng dữ kiện và giả thiết. (3) Năng lượng toàn phần bằng năng lượng mỗi phản ứng nhân số phản ứng; số phản ứng được suy từ số hạt tham gia. (4) Có thể thay số trước khi chọn công thức.
\shortans{2}
\loigiai{
Đối chiếu từng phát biểu với kiến thức cốt lõi: Năng lượng toàn phần bằng năng lượng mỗi phản ứng nhân số phản ứng; số phản ứng được suy từ số hạt tham gia. Có 2 phát biểu đúng.
}
\end{ex}

% ===== Câu 92 =====
% ID: L12C4B24-23-TN
% Mức: NB
\dangbt{Nhận biết phản ứng nhiệt hạch và điều kiện xảy ra}
\begin{ex}
% ID: L12C4B24-23-TN
% Mức: NB
Chọn kết luận chính xác về nhận biết phản ứng nhiệt hạch và điều kiện xảy ra?
\choice
{Nhiệt hạch dễ xảy ra ở nhiệt độ phòng.}
{Có thể bỏ qua điều kiện áp dụng và đơn vị trong mọi trường hợp.}
{\True Nhiệt hạch là sự kết hợp các hạt nhân nhẹ, cần nhiệt độ rất cao để thắng lực đẩy Coulomb và có thể tỏa năng lượng lớn.}
{Kết quả của bài toán không phụ thuộc dữ kiện đã cho.}
\loigiai{
Nhiệt hạch là sự kết hợp các hạt nhân nhẹ, cần nhiệt độ rất cao để thắng lực đẩy Coulomb và có thể tỏa năng lượng lớn. Vì vậy phương án $C$ đúng; các phương án còn lại trái với kiến thức hoặc điều kiện áp dụng.
}
\end{ex}

% ===== Câu 93 =====
% ID: L12C4B24-24-DS
% Mức: VD
\dangbt{Cấu tạo và nguyên lí hoạt động của lò phản ứng hạt nhân}
\begin{ex}
% ID: L12C4B24-24-DS
% Mức: VD
Xét cấu tạo và nguyên lí hoạt động của lò phản ứng hạt nhân (tình huống 4). Hãy xác định tính đúng, sai của bốn phát biểu sau.
\choiceTF
{Thanh điều khiển có nhiệm vụ làm tăng không giới hạn số neutron trong lò.}
{\True Lò phản ứng dùng nhiên liệu phân hạch, chất làm chậm, thanh điều khiển, chất tải nhiệt và hệ che chắn để duy trì phản ứng có kiểm soát.}
{Có thể bỏ qua dữ kiện, đơn vị và ý nghĩa vật lí của kết quả.}
{\True Khi giải cần đọc đúng dữ kiện, dùng hệ đơn vị thống nhất và kiểm tra điều kiện áp dụng.}
\loigiai{
\begin{itemchoice}
\item (Sai) Thanh điều khiển có nhiệm vụ làm tăng không giới hạn số neutron trong lò. (Vì): Thanh điều khiển có nhiệm vụ làm tăng không giới hạn số neutron trong lò. b) Đúng: Lò phản ứng dùng nhiên liệu phân hạch, chất làm chậm, thanh điều khiển, chất tải nhiệt và hệ che chắn để duy trì phản ứng có kiểm soát. c) Sai: Có thể bỏ qua dữ kiện, đơn vị và ý nghĩa vật lí của kết quả. d) Đúng: Khi giải cần đọc đúng dữ kiện, dùng hệ đơn vị thống nhất và kiểm tra điều kiện áp dụng. Kiến thức cốt lõi: Lò phản ứng dùng nhiên liệu phân hạch, chất làm chậm, thanh điều khiển, chất tải nhiệt và hệ che chắn để duy trì phản ứng có kiểm soát
\item (Đúng) Lò phản ứng dùng nhiên liệu phân hạch, chất làm chậm, thanh điều khiển, chất tải nhiệt và hệ che chắn để duy trì phản ứng có kiểm soát.
\item (Sai) Có thể bỏ qua dữ kiện, đơn vị và ý nghĩa vật lí của kết quả.
\item (Đúng) Khi giải cần đọc đúng dữ kiện, dùng hệ đơn vị thống nhất và kiểm tra điều kiện áp dụng.
\end{itemchoice}
}
\end{ex}

% ===== Câu 94 =====
% ID: L12C4B24-24-TL
% Mức: VDC
\dangbt{Tính năng lượng trong phân hạch và nhiệt hạch}
\begin{bt}
% ID: L12C4B24-24-TL
% Mức: VDC
So sánh nhận định đúng “Năng lượng toàn phần bằng năng lượng mỗi phản ứng nhân số phản ứng; số phản ứng được suy từ số hạt tham gia.” với nhận định sai “Năng lượng toàn phần không phụ thuộc số hạt đã phản ứng.”; chỉ rõ căn cứ Vật lí.
\loigiai{
Cần nêu được: Năng lượng toàn phần bằng năng lượng mỗi phản ứng nhân số phản ứng; số phản ứng được suy từ số hạt tham gia. Khi vận dụng phải xác định đúng dữ kiện, điều kiện áp dụng, hệ đơn vị và kiểm tra tính hợp lí của kết quả. Nhận định “Năng lượng toàn phần không phụ thuộc số hạt đã phản ứng.” sai vì trái với kiến thức trên.
}
\end{bt}

% ===== Câu 95 =====
% ID: L12C4B24-24-TLN
% Mức: VDC
\begin{ex}
% ID: L12C4B24-24-TLN
% Mức: VDC
Trong bốn phát biểu sau về tính năng lượng trong phân hạch và nhiệt hạch, có bao nhiêu phát biểu đúng? Chỉ ghi một số nguyên. (1) Kết quả cần có ý nghĩa vật lí. (2) Năng lượng toàn phần không phụ thuộc số hạt đã phản ứng. (3) Phải dùng đúng định luật cho tình huống. (4) Năng lượng toàn phần bằng năng lượng mỗi phản ứng nhân số phản ứng; số phản ứng được suy từ số hạt tham gia.
\shortans{3}
\loigiai{
Đối chiếu từng phát biểu với kiến thức cốt lõi: Năng lượng toàn phần bằng năng lượng mỗi phản ứng nhân số phản ứng; số phản ứng được suy từ số hạt tham gia. Có 3 phát biểu đúng.
}
\end{ex}

% ===== Câu 96 =====
% ID: L12C4B24-24-TN
% Mức: TH
\dangbt{Nhận biết phản ứng nhiệt hạch và điều kiện xảy ra}
\begin{ex}
% ID: L12C4B24-24-TN
% Mức: TH
Nội dung nào mô tả đúng nhất về nhận biết phản ứng nhiệt hạch và điều kiện xảy ra?
\choice
{Nhiệt hạch dễ xảy ra ở nhiệt độ phòng.}
{Có thể bỏ qua điều kiện áp dụng và đơn vị trong mọi trường hợp.}
{Kết quả của bài toán không phụ thuộc dữ kiện đã cho.}
{\True Nhiệt hạch là sự kết hợp các hạt nhân nhẹ, cần nhiệt độ rất cao để thắng lực đẩy Coulomb và có thể tỏa năng lượng lớn.}
\loigiai{
Nhiệt hạch là sự kết hợp các hạt nhân nhẹ, cần nhiệt độ rất cao để thắng lực đẩy Coulomb và có thể tỏa năng lượng lớn. Vì vậy phương án $D$ đúng; các phương án còn lại trái với kiến thức hoặc điều kiện áp dụng.
}
\end{ex}

% ===== Câu 97 =====
% ID: L12C4B24-25-DS
% Mức: VD
\dangbt{Cấu tạo và nguyên lí hoạt động của lò phản ứng hạt nhân}
\begin{ex}
% ID: L12C4B24-25-DS
% Mức: VD
Xét cấu tạo và nguyên lí hoạt động của lò phản ứng hạt nhân (tình huống 5). Hãy xác định tính đúng, sai của bốn phát biểu sau.
\choiceTF
{\True Lò phản ứng dùng nhiên liệu phân hạch, chất làm chậm, thanh điều khiển, chất tải nhiệt và hệ che chắn để duy trì phản ứng có kiểm soát.}
{Có thể bỏ qua dữ kiện, đơn vị và ý nghĩa vật lí của kết quả.}
{Thanh điều khiển có nhiệm vụ làm tăng không giới hạn số neutron trong lò.}
{Có thể bỏ qua dữ kiện, đơn vị và ý nghĩa vật lí của kết quả.}
\loigiai{
\begin{itemchoice}
\item (Đúng) Lò phản ứng dùng nhiên liệu phân hạch, chất làm chậm, thanh điều khiển, chất tải nhiệt và hệ che chắn để duy trì phản ứng có kiểm soát. (Vì): Lò phản ứng dùng nhiên liệu phân hạch, chất làm chậm, thanh điều khiển, chất tải nhiệt và hệ che chắn để duy trì phản ứng có kiểm soát. b) Sai: Có thể bỏ qua dữ kiện, đơn vị và ý nghĩa vật lí của kết quả. c) Sai: Thanh điều khiển có nhiệm vụ làm tăng không giới hạn số neutron trong lò. d) Sai: Có thể bỏ qua dữ kiện, đơn vị và ý nghĩa vật lí của kết quả. Kiến thức cốt lõi: Lò phản ứng dùng nhiên liệu phân hạch, chất làm chậm, thanh điều khiển, chất tải nhiệt và hệ che chắn để duy trì phản ứng có kiểm soát
\item (Sai) Có thể bỏ qua dữ kiện, đơn vị và ý nghĩa vật lí của kết quả.
\item (Sai) Thanh điều khiển có nhiệm vụ làm tăng không giới hạn số neutron trong lò.
\item (Sai) Có thể bỏ qua dữ kiện, đơn vị và ý nghĩa vật lí của kết quả.
\end{itemchoice}
}
\end{ex}

% ===== Câu 98 =====
% ID: L12C4B24-25-TL
% Mức: TH
\begin{bt}
% ID: L12C4B24-25-TL
% Mức: TH
Trình bày kiến thức cốt lõi của cấu tạo và nguyên lí hoạt động của lò phản ứng hạt nhân và nêu điều kiện áp dụng.
\loigiai{
Cần nêu được: Lò phản ứng dùng nhiên liệu phân hạch, chất làm chậm, thanh điều khiển, chất tải nhiệt và hệ che chắn để duy trì phản ứng có kiểm soát. Khi vận dụng phải xác định đúng dữ kiện, điều kiện áp dụng, hệ đơn vị và kiểm tra tính hợp lí của kết quả. Nhận định “Thanh điều khiển có nhiệm vụ làm tăng không giới hạn số neutron trong lò.” sai vì trái với kiến thức trên.
}
\end{bt}

% ===== Câu 99 =====
% ID: L12C4B24-25-TLN
% Mức: TH
\begin{ex}
% ID: L12C4B24-25-TLN
% Mức: TH
Trong bốn phát biểu sau về cấu tạo và nguyên lí hoạt động của lò phản ứng hạt nhân, có bao nhiêu phát biểu đúng? Chỉ ghi một số nguyên. (1) Lò phản ứng dùng nhiên liệu phân hạch, chất làm chậm, thanh điều khiển, chất tải nhiệt và hệ che chắn để duy trì phản ứng có kiểm soát. (2) Thanh điều khiển có nhiệm vụ làm tăng không giới hạn số neutron trong lò. (3) Cần kiểm tra điều kiện áp dụng trước khi tính toán. (4) Có thể bỏ qua đơn vị của kết quả.
\shortans{2}
\loigiai{
Đối chiếu từng phát biểu với kiến thức cốt lõi: Lò phản ứng dùng nhiên liệu phân hạch, chất làm chậm, thanh điều khiển, chất tải nhiệt và hệ che chắn để duy trì phản ứng có kiểm soát. Có 2 phát biểu đúng.
}
\end{ex}

% ===== Câu 100 =====
% ID: L12C4B24-25-TN
% Mức: TH
\dangbt{Nhận biết phản ứng nhiệt hạch và điều kiện xảy ra}
\begin{ex}
% ID: L12C4B24-25-TN
% Mức: TH
Khi phân tích, phát biểu nào đúng về nhận biết phản ứng nhiệt hạch và điều kiện xảy ra?
\choice
{\True Nhiệt hạch là sự kết hợp các hạt nhân nhẹ, cần nhiệt độ rất cao để thắng lực đẩy Coulomb và có thể tỏa năng lượng lớn.}
{Nhiệt hạch dễ xảy ra ở nhiệt độ phòng.}
{Có thể bỏ qua điều kiện áp dụng và đơn vị trong mọi trường hợp.}
{Kết quả của bài toán không phụ thuộc dữ kiện đã cho.}
\loigiai{
Nhiệt hạch là sự kết hợp các hạt nhân nhẹ, cần nhiệt độ rất cao để thắng lực đẩy Coulomb và có thể tỏa năng lượng lớn. Vì vậy phương án $A$ đúng; các phương án còn lại trái với kiến thức hoặc điều kiện áp dụng.
}
\end{ex}

% ===== Câu 101 =====
% ID: L12C4B24-26-DS
% Mức: TH
\dangbt{Điện hạt nhân, nhiên liệu hạt nhân và xử lí chất thải}
\begin{ex}
% ID: L12C4B24-26-DS
% Mức: TH
Xét điện hạt nhân, nhiên liệu hạt nhân và xử lí chất thải (tình huống 1). Hãy xác định tính đúng, sai của bốn phát biểu sau.
\choiceTF
{\True Nhà máy điện hạt nhân biến năng lượng hạt nhân thành nhiệt, cơ năng rồi điện năng; nhiên liệu và chất thải phải được quản lí nghiêm ngặt.}
{Chất thải phóng xạ có thể xử lí như rác sinh hoạt thông thường.}
{\True Khi giải cần đọc đúng dữ kiện, dùng hệ đơn vị thống nhất và kiểm tra điều kiện áp dụng.}
{Có thể bỏ qua dữ kiện, đơn vị và ý nghĩa vật lí của kết quả.}
\loigiai{
\begin{itemchoice}
\item (Đúng) Nhà máy điện hạt nhân biến năng lượng hạt nhân thành nhiệt, cơ năng rồi điện năng; nhiên liệu và chất thải phải được quản lí nghiêm ngặt. (Vì): Nhà máy điện hạt nhân biến năng lượng hạt nhân thành nhiệt, cơ năng rồi điện năng; nhiên liệu và chất thải phải được quản lí nghiêm ngặt. b) Sai: Chất thải phóng xạ có thể xử lí như rác sinh hoạt thông thường. c) Đúng: Khi giải cần đọc đúng dữ kiện, dùng hệ đơn vị thống nhất và kiểm tra điều kiện áp dụng. d) Sai: Có thể bỏ qua dữ kiện, đơn vị và ý nghĩa vật lí của kết quả. Kiến thức cốt lõi: Nhà máy điện hạt nhân biến năng lượng hạt nhân thành nhiệt, cơ năng rồi điện năng; nhiên liệu và chất thải phải được quản lí nghiêm ngặt
\item (Sai) Chất thải phóng xạ có thể xử lí như rác sinh hoạt thông thường.
\item (Đúng) Khi giải cần đọc đúng dữ kiện, dùng hệ đơn vị thống nhất và kiểm tra điều kiện áp dụng.
\item (Sai) Có thể bỏ qua dữ kiện, đơn vị và ý nghĩa vật lí của kết quả.
\end{itemchoice}
}
\end{ex}

% ===== Câu 102 =====
% ID: L12C4B24-26-TL
% Mức: TH
\dangbt{Cấu tạo và nguyên lí hoạt động của lò phản ứng hạt nhân}
\begin{bt}
% ID: L12C4B24-26-TL
% Mức: TH
Giải thích vì sao nhận định sau đúng: “Lò phản ứng dùng nhiên liệu phân hạch, chất làm chậm, thanh điều khiển, chất tải nhiệt và hệ che chắn để duy trì phản ứng có kiểm soát.”
\loigiai{
Cần nêu được: Lò phản ứng dùng nhiên liệu phân hạch, chất làm chậm, thanh điều khiển, chất tải nhiệt và hệ che chắn để duy trì phản ứng có kiểm soát. Khi vận dụng phải xác định đúng dữ kiện, điều kiện áp dụng, hệ đơn vị và kiểm tra tính hợp lí của kết quả. Nhận định “Thanh điều khiển có nhiệm vụ làm tăng không giới hạn số neutron trong lò.” sai vì trái với kiến thức trên.
}
\end{bt}

% ===== Câu 103 =====
% ID: L12C4B24-26-TLN
% Mức: TH
\begin{ex}
% ID: L12C4B24-26-TLN
% Mức: TH
Trong bốn phát biểu sau về cấu tạo và nguyên lí hoạt động của lò phản ứng hạt nhân, có bao nhiêu phát biểu đúng? Chỉ ghi một số nguyên. (1) Thanh điều khiển có nhiệm vụ làm tăng không giới hạn số neutron trong lò. (2) Lò phản ứng dùng nhiên liệu phân hạch, chất làm chậm, thanh điều khiển, chất tải nhiệt và hệ che chắn để duy trì phản ứng có kiểm soát. (3) Kết quả phải phù hợp ý nghĩa vật lí. (4) Mọi đại lượng đều có cùng bản chất.
\shortans{2}
\loigiai{
Đối chiếu từng phát biểu với kiến thức cốt lõi: Lò phản ứng dùng nhiên liệu phân hạch, chất làm chậm, thanh điều khiển, chất tải nhiệt và hệ che chắn để duy trì phản ứng có kiểm soát. Có 2 phát biểu đúng.
}
\end{ex}

% ===== Câu 104 =====
% ID: L12C4B24-26-TN
% Mức: TH
\dangbt{Nhận biết phản ứng nhiệt hạch và điều kiện xảy ra}
\begin{ex}
% ID: L12C4B24-26-TN
% Mức: TH
Kết luận nào của học sinh là đúng về nhận biết phản ứng nhiệt hạch và điều kiện xảy ra?
\choice
{Nhiệt hạch dễ xảy ra ở nhiệt độ phòng.}
{\True Nhiệt hạch là sự kết hợp các hạt nhân nhẹ, cần nhiệt độ rất cao để thắng lực đẩy Coulomb và có thể tỏa năng lượng lớn.}
{Có thể bỏ qua điều kiện áp dụng và đơn vị trong mọi trường hợp.}
{Kết quả của bài toán không phụ thuộc dữ kiện đã cho.}
\loigiai{
Nhiệt hạch là sự kết hợp các hạt nhân nhẹ, cần nhiệt độ rất cao để thắng lực đẩy Coulomb và có thể tỏa năng lượng lớn. Vì vậy phương án $B$ đúng; các phương án còn lại trái với kiến thức hoặc điều kiện áp dụng.
}
\end{ex}

% ===== Câu 105 =====
% ID: L12C4B24-27-DS
% Mức: TH
\dangbt{Điện hạt nhân, nhiên liệu hạt nhân và xử lí chất thải}
\begin{ex}
% ID: L12C4B24-27-DS
% Mức: TH
Xét điện hạt nhân, nhiên liệu hạt nhân và xử lí chất thải (tình huống 2). Hãy xác định tính đúng, sai của bốn phát biểu sau.
\choiceTF
{Chất thải phóng xạ có thể xử lí như rác sinh hoạt thông thường.}
{\True Nhà máy điện hạt nhân biến năng lượng hạt nhân thành nhiệt, cơ năng rồi điện năng; nhiên liệu và chất thải phải được quản lí nghiêm ngặt.}
{Có thể bỏ qua dữ kiện, đơn vị và ý nghĩa vật lí của kết quả.}
{\True Khi giải cần đọc đúng dữ kiện, dùng hệ đơn vị thống nhất và kiểm tra điều kiện áp dụng.}
\loigiai{
\begin{itemchoice}
\item (Sai) Chất thải phóng xạ có thể xử lí như rác sinh hoạt thông thường. (Vì): Chất thải phóng xạ có thể xử lí như rác sinh hoạt thông thường. b) Đúng: Nhà máy điện hạt nhân biến năng lượng hạt nhân thành nhiệt, cơ năng rồi điện năng; nhiên liệu và chất thải phải được quản lí nghiêm ngặt. c) Sai: Có thể bỏ qua dữ kiện, đơn vị và ý nghĩa vật lí của kết quả. d) Đúng: Khi giải cần đọc đúng dữ kiện, dùng hệ đơn vị thống nhất và kiểm tra điều kiện áp dụng. Kiến thức cốt lõi: Nhà máy điện hạt nhân biến năng lượng hạt nhân thành nhiệt, cơ năng rồi điện năng; nhiên liệu và chất thải phải được quản lí nghiêm ngặt
\item (Đúng) Nhà máy điện hạt nhân biến năng lượng hạt nhân thành nhiệt, cơ năng rồi điện năng; nhiên liệu và chất thải phải được quản lí nghiêm ngặt.
\item (Sai) Có thể bỏ qua dữ kiện, đơn vị và ý nghĩa vật lí của kết quả.
\item (Đúng) Khi giải cần đọc đúng dữ kiện, dùng hệ đơn vị thống nhất và kiểm tra điều kiện áp dụng.
\end{itemchoice}
}
\end{ex}

% ===== Câu 106 =====
% ID: L12C4B24-27-TL
% Mức: VD
\dangbt{Cấu tạo và nguyên lí hoạt động của lò phản ứng hạt nhân}
\begin{bt}
% ID: L12C4B24-27-TL
% Mức: VD
Phân tích sai lầm trong nhận định: “Thanh điều khiển có nhiệm vụ làm tăng không giới hạn số neutron trong lò.”
\loigiai{
Cần nêu được: Lò phản ứng dùng nhiên liệu phân hạch, chất làm chậm, thanh điều khiển, chất tải nhiệt và hệ che chắn để duy trì phản ứng có kiểm soát. Khi vận dụng phải xác định đúng dữ kiện, điều kiện áp dụng, hệ đơn vị và kiểm tra tính hợp lí của kết quả. Nhận định “Thanh điều khiển có nhiệm vụ làm tăng không giới hạn số neutron trong lò.” sai vì trái với kiến thức trên.
}
\end{bt}

% ===== Câu 107 =====
% ID: L12C4B24-27-TLN
% Mức: TH
\begin{ex}
% ID: L12C4B24-27-TLN
% Mức: TH
Trong bốn phát biểu sau về cấu tạo và nguyên lí hoạt động của lò phản ứng hạt nhân, có bao nhiêu phát biểu đúng? Chỉ ghi một số nguyên. (1) Cần đổi các đại lượng về hệ đơn vị thống nhất. (2) Lò phản ứng dùng nhiên liệu phân hạch, chất làm chậm, thanh điều khiển, chất tải nhiệt và hệ che chắn để duy trì phản ứng có kiểm soát. (3) Thanh điều khiển có nhiệm vụ làm tăng không giới hạn số neutron trong lò. (4) Không cần kiểm tra dữ kiện.
\shortans{2}
\loigiai{
Đối chiếu từng phát biểu với kiến thức cốt lõi: Lò phản ứng dùng nhiên liệu phân hạch, chất làm chậm, thanh điều khiển, chất tải nhiệt và hệ che chắn để duy trì phản ứng có kiểm soát. Có 2 phát biểu đúng.
}
\end{ex}

% ===== Câu 108 =====
% ID: L12C4B24-27-TN
% Mức: TH
\dangbt{Nhận biết phản ứng nhiệt hạch và điều kiện xảy ra}
\begin{ex}
% ID: L12C4B24-27-TN
% Mức: TH
Chọn phát biểu không trái với kiến thức về nhận biết phản ứng nhiệt hạch và điều kiện xảy ra?
\choice
{Nhiệt hạch dễ xảy ra ở nhiệt độ phòng.}
{Có thể bỏ qua điều kiện áp dụng và đơn vị trong mọi trường hợp.}
{\True Nhiệt hạch là sự kết hợp các hạt nhân nhẹ, cần nhiệt độ rất cao để thắng lực đẩy Coulomb và có thể tỏa năng lượng lớn.}
{Kết quả của bài toán không phụ thuộc dữ kiện đã cho.}
\loigiai{
Nhiệt hạch là sự kết hợp các hạt nhân nhẹ, cần nhiệt độ rất cao để thắng lực đẩy Coulomb và có thể tỏa năng lượng lớn. Vì vậy phương án $C$ đúng; các phương án còn lại trái với kiến thức hoặc điều kiện áp dụng.
}
\end{ex}

% ===== Câu 109 =====
% ID: L12C4B24-28-DS
% Mức: VD
\dangbt{Điện hạt nhân, nhiên liệu hạt nhân và xử lí chất thải}
\begin{ex}
% ID: L12C4B24-28-DS
% Mức: VD
Xét điện hạt nhân, nhiên liệu hạt nhân và xử lí chất thải (tình huống 3). Hãy xác định tính đúng, sai của bốn phát biểu sau.
\choiceTF
{\True Nhà máy điện hạt nhân biến năng lượng hạt nhân thành nhiệt, cơ năng rồi điện năng; nhiên liệu và chất thải phải được quản lí nghiêm ngặt.}
{Chất thải phóng xạ có thể xử lí như rác sinh hoạt thông thường.}
{\True Khi giải cần đọc đúng dữ kiện, dùng hệ đơn vị thống nhất và kiểm tra điều kiện áp dụng.}
{Có thể bỏ qua dữ kiện, đơn vị và ý nghĩa vật lí của kết quả.}
\loigiai{
\begin{itemchoice}
\item (Đúng) Nhà máy điện hạt nhân biến năng lượng hạt nhân thành nhiệt, cơ năng rồi điện năng; nhiên liệu và chất thải phải được quản lí nghiêm ngặt. (Vì): Nhà máy điện hạt nhân biến năng lượng hạt nhân thành nhiệt, cơ năng rồi điện năng; nhiên liệu và chất thải phải được quản lí nghiêm ngặt. b) Sai: Chất thải phóng xạ có thể xử lí như rác sinh hoạt thông thường. c) Đúng: Khi giải cần đọc đúng dữ kiện, dùng hệ đơn vị thống nhất và kiểm tra điều kiện áp dụng. d) Sai: Có thể bỏ qua dữ kiện, đơn vị và ý nghĩa vật lí của kết quả. Kiến thức cốt lõi: Nhà máy điện hạt nhân biến năng lượng hạt nhân thành nhiệt, cơ năng rồi điện năng; nhiên liệu và chất thải phải được quản lí nghiêm ngặt
\item (Sai) Chất thải phóng xạ có thể xử lí như rác sinh hoạt thông thường.
\item (Đúng) Khi giải cần đọc đúng dữ kiện, dùng hệ đơn vị thống nhất và kiểm tra điều kiện áp dụng.
\item (Sai) Có thể bỏ qua dữ kiện, đơn vị và ý nghĩa vật lí của kết quả.
\end{itemchoice}
}
\end{ex}

% ===== Câu 110 =====
% ID: L12C4B24-28-TL
% Mức: VD
\dangbt{Cấu tạo và nguyên lí hoạt động của lò phản ứng hạt nhân}
\begin{bt}
% ID: L12C4B24-28-TL
% Mức: VD
Đề xuất các bước giải một bài toán thuộc dạng cấu tạo và nguyên lí hoạt động của lò phản ứng hạt nhân.
\loigiai{
Cần nêu được: Lò phản ứng dùng nhiên liệu phân hạch, chất làm chậm, thanh điều khiển, chất tải nhiệt và hệ che chắn để duy trì phản ứng có kiểm soát. Khi vận dụng phải xác định đúng dữ kiện, điều kiện áp dụng, hệ đơn vị và kiểm tra tính hợp lí của kết quả. Nhận định “Thanh điều khiển có nhiệm vụ làm tăng không giới hạn số neutron trong lò.” sai vì trái với kiến thức trên.
}
\end{bt}

% ===== Câu 111 =====
% ID: L12C4B24-28-TLN
% Mức: VD
\begin{ex}
% ID: L12C4B24-28-TLN
% Mức: VD
Trong bốn phát biểu sau về cấu tạo và nguyên lí hoạt động của lò phản ứng hạt nhân, có bao nhiêu phát biểu đúng? Chỉ ghi một số nguyên. (1) Lò phản ứng dùng nhiên liệu phân hạch, chất làm chậm, thanh điều khiển, chất tải nhiệt và hệ che chắn để duy trì phản ứng có kiểm soát. (2) Cần đối chiếu kết quả với điều kiện bài toán. (3) Thanh điều khiển có nhiệm vụ làm tăng không giới hạn số neutron trong lò. (4) Mọi trường hợp đều cho cùng một kết quả.
\shortans{2}
\loigiai{
Đối chiếu từng phát biểu với kiến thức cốt lõi: Lò phản ứng dùng nhiên liệu phân hạch, chất làm chậm, thanh điều khiển, chất tải nhiệt và hệ che chắn để duy trì phản ứng có kiểm soát. Có 2 phát biểu đúng.
}
\end{ex}

% ===== Câu 112 =====
% ID: L12C4B24-28-TN
% Mức: VD
\dangbt{Nhận biết phản ứng nhiệt hạch và điều kiện xảy ra}
\begin{ex}
% ID: L12C4B24-28-TN
% Mức: VD
Cơ sở Vật lí đúng để giải bài là gì về nhận biết phản ứng nhiệt hạch và điều kiện xảy ra?
\choice
{Nhiệt hạch dễ xảy ra ở nhiệt độ phòng.}
{Có thể bỏ qua điều kiện áp dụng và đơn vị trong mọi trường hợp.}
{Kết quả của bài toán không phụ thuộc dữ kiện đã cho.}
{\True Nhiệt hạch là sự kết hợp các hạt nhân nhẹ, cần nhiệt độ rất cao để thắng lực đẩy Coulomb và có thể tỏa năng lượng lớn.}
\loigiai{
Nhiệt hạch là sự kết hợp các hạt nhân nhẹ, cần nhiệt độ rất cao để thắng lực đẩy Coulomb và có thể tỏa năng lượng lớn. Vì vậy phương án $D$ đúng; các phương án còn lại trái với kiến thức hoặc điều kiện áp dụng.
}
\end{ex}

% ===== Câu 113 =====
% ID: L12C4B24-29-DS
% Mức: VD
\dangbt{Điện hạt nhân, nhiên liệu hạt nhân và xử lí chất thải}
\begin{ex}
% ID: L12C4B24-29-DS
% Mức: VD
Xét điện hạt nhân, nhiên liệu hạt nhân và xử lí chất thải (tình huống 4). Hãy xác định tính đúng, sai của bốn phát biểu sau.
\choiceTF
{Chất thải phóng xạ có thể xử lí như rác sinh hoạt thông thường.}
{\True Nhà máy điện hạt nhân biến năng lượng hạt nhân thành nhiệt, cơ năng rồi điện năng; nhiên liệu và chất thải phải được quản lí nghiêm ngặt.}
{Có thể bỏ qua dữ kiện, đơn vị và ý nghĩa vật lí của kết quả.}
{\True Khi giải cần đọc đúng dữ kiện, dùng hệ đơn vị thống nhất và kiểm tra điều kiện áp dụng.}
\loigiai{
\begin{itemchoice}
\item (Sai) Chất thải phóng xạ có thể xử lí như rác sinh hoạt thông thường. (Vì): Chất thải phóng xạ có thể xử lí như rác sinh hoạt thông thường. b) Đúng: Nhà máy điện hạt nhân biến năng lượng hạt nhân thành nhiệt, cơ năng rồi điện năng; nhiên liệu và chất thải phải được quản lí nghiêm ngặt. c) Sai: Có thể bỏ qua dữ kiện, đơn vị và ý nghĩa vật lí của kết quả. d) Đúng: Khi giải cần đọc đúng dữ kiện, dùng hệ đơn vị thống nhất và kiểm tra điều kiện áp dụng. Kiến thức cốt lõi: Nhà máy điện hạt nhân biến năng lượng hạt nhân thành nhiệt, cơ năng rồi điện năng; nhiên liệu và chất thải phải được quản lí nghiêm ngặt
\item (Đúng) Nhà máy điện hạt nhân biến năng lượng hạt nhân thành nhiệt, cơ năng rồi điện năng; nhiên liệu và chất thải phải được quản lí nghiêm ngặt.
\item (Sai) Có thể bỏ qua dữ kiện, đơn vị và ý nghĩa vật lí của kết quả.
\item (Đúng) Khi giải cần đọc đúng dữ kiện, dùng hệ đơn vị thống nhất và kiểm tra điều kiện áp dụng.
\end{itemchoice}
}
\end{ex}

% ===== Câu 114 =====
% ID: L12C4B24-29-TL
% Mức: VD
\dangbt{Cấu tạo và nguyên lí hoạt động của lò phản ứng hạt nhân}
\begin{bt}
% ID: L12C4B24-29-TL
% Mức: VD
Nêu một tình huống thực tiễn liên quan đến cấu tạo và nguyên lí hoạt động của lò phản ứng hạt nhân và giải thích bằng kiến thức Vật lí.
\loigiai{
Cần nêu được: Lò phản ứng dùng nhiên liệu phân hạch, chất làm chậm, thanh điều khiển, chất tải nhiệt và hệ che chắn để duy trì phản ứng có kiểm soát. Khi vận dụng phải xác định đúng dữ kiện, điều kiện áp dụng, hệ đơn vị và kiểm tra tính hợp lí của kết quả. Nhận định “Thanh điều khiển có nhiệm vụ làm tăng không giới hạn số neutron trong lò.” sai vì trái với kiến thức trên.
}
\end{bt}

% ===== Câu 115 =====
% ID: L12C4B24-29-TLN
% Mức: VD
\begin{ex}
% ID: L12C4B24-29-TLN
% Mức: VD
Trong bốn phát biểu sau về cấu tạo và nguyên lí hoạt động của lò phản ứng hạt nhân, có bao nhiêu phát biểu đúng? Chỉ ghi một số nguyên. (1) Thanh điều khiển có nhiệm vụ làm tăng không giới hạn số neutron trong lò. (2) Cần đọc đúng dữ kiện và giả thiết. (3) Lò phản ứng dùng nhiên liệu phân hạch, chất làm chậm, thanh điều khiển, chất tải nhiệt và hệ che chắn để duy trì phản ứng có kiểm soát. (4) Có thể thay số trước khi chọn công thức.
\shortans{2}
\loigiai{
Đối chiếu từng phát biểu với kiến thức cốt lõi: Lò phản ứng dùng nhiên liệu phân hạch, chất làm chậm, thanh điều khiển, chất tải nhiệt và hệ che chắn để duy trì phản ứng có kiểm soát. Có 2 phát biểu đúng.
}
\end{ex}

% ===== Câu 116 =====
% ID: L12C4B24-29-TN
% Mức: VD
\dangbt{Nhận biết phản ứng nhiệt hạch và điều kiện xảy ra}
\begin{ex}
% ID: L12C4B24-29-TN
% Mức: VD
Trong một bài thực hành, nhận xét nào đúng về nhận biết phản ứng nhiệt hạch và điều kiện xảy ra?
\choice
{\True Nhiệt hạch là sự kết hợp các hạt nhân nhẹ, cần nhiệt độ rất cao để thắng lực đẩy Coulomb và có thể tỏa năng lượng lớn.}
{Nhiệt hạch dễ xảy ra ở nhiệt độ phòng.}
{Có thể bỏ qua điều kiện áp dụng và đơn vị trong mọi trường hợp.}
{Kết quả của bài toán không phụ thuộc dữ kiện đã cho.}
\loigiai{
Nhiệt hạch là sự kết hợp các hạt nhân nhẹ, cần nhiệt độ rất cao để thắng lực đẩy Coulomb và có thể tỏa năng lượng lớn. Vì vậy phương án $A$ đúng; các phương án còn lại trái với kiến thức hoặc điều kiện áp dụng.
}
\end{ex}

% ===== Câu 117 =====
% ID: L12C4B24-30-DS
% Mức: VD
\dangbt{Điện hạt nhân, nhiên liệu hạt nhân và xử lí chất thải}
\begin{ex}
% ID: L12C4B24-30-DS
% Mức: VD
Xét điện hạt nhân, nhiên liệu hạt nhân và xử lí chất thải (tình huống 5). Hãy xác định tính đúng, sai của bốn phát biểu sau.
\choiceTF
{\True Nhà máy điện hạt nhân biến năng lượng hạt nhân thành nhiệt, cơ năng rồi điện năng; nhiên liệu và chất thải phải được quản lí nghiêm ngặt.}
{Có thể bỏ qua dữ kiện, đơn vị và ý nghĩa vật lí của kết quả.}
{Chất thải phóng xạ có thể xử lí như rác sinh hoạt thông thường.}
{Có thể bỏ qua dữ kiện, đơn vị và ý nghĩa vật lí của kết quả.}
\loigiai{
\begin{itemchoice}
\item (Đúng) Nhà máy điện hạt nhân biến năng lượng hạt nhân thành nhiệt, cơ năng rồi điện năng; nhiên liệu và chất thải phải được quản lí nghiêm ngặt. (Vì): Nhà máy điện hạt nhân biến năng lượng hạt nhân thành nhiệt, cơ năng rồi điện năng; nhiên liệu và chất thải phải được quản lí nghiêm ngặt. b) Sai: Có thể bỏ qua dữ kiện, đơn vị và ý nghĩa vật lí của kết quả. c) Sai: Chất thải phóng xạ có thể xử lí như rác sinh hoạt thông thường. d) Sai: Có thể bỏ qua dữ kiện, đơn vị và ý nghĩa vật lí của kết quả. Kiến thức cốt lõi: Nhà máy điện hạt nhân biến năng lượng hạt nhân thành nhiệt, cơ năng rồi điện năng; nhiên liệu và chất thải phải được quản lí nghiêm ngặt
\item (Sai) Có thể bỏ qua dữ kiện, đơn vị và ý nghĩa vật lí của kết quả.
\item (Sai) Chất thải phóng xạ có thể xử lí như rác sinh hoạt thông thường.
\item (Sai) Có thể bỏ qua dữ kiện, đơn vị và ý nghĩa vật lí của kết quả.
\end{itemchoice}
}
\end{ex}

% ===== Câu 118 =====
% ID: L12C4B24-30-TL
% Mức: VDC
\dangbt{Cấu tạo và nguyên lí hoạt động của lò phản ứng hạt nhân}
\begin{bt}
% ID: L12C4B24-30-TL
% Mức: VDC
So sánh nhận định đúng “Lò phản ứng dùng nhiên liệu phân hạch, chất làm chậm, thanh điều khiển, chất tải nhiệt và hệ che chắn để duy trì phản ứng có kiểm soát.” với nhận định sai “Thanh điều khiển có nhiệm vụ làm tăng không giới hạn số neutron trong lò.”; chỉ rõ căn cứ Vật lí.
\loigiai{
Cần nêu được: Lò phản ứng dùng nhiên liệu phân hạch, chất làm chậm, thanh điều khiển, chất tải nhiệt và hệ che chắn để duy trì phản ứng có kiểm soát. Khi vận dụng phải xác định đúng dữ kiện, điều kiện áp dụng, hệ đơn vị và kiểm tra tính hợp lí của kết quả. Nhận định “Thanh điều khiển có nhiệm vụ làm tăng không giới hạn số neutron trong lò.” sai vì trái với kiến thức trên.
}
\end{bt}

% ===== Câu 119 =====
% ID: L12C4B24-30-TLN
% Mức: VDC
\begin{ex}
% ID: L12C4B24-30-TLN
% Mức: VDC
Trong bốn phát biểu sau về cấu tạo và nguyên lí hoạt động của lò phản ứng hạt nhân, có bao nhiêu phát biểu đúng? Chỉ ghi một số nguyên. (1) Kết quả cần có ý nghĩa vật lí. (2) Thanh điều khiển có nhiệm vụ làm tăng không giới hạn số neutron trong lò. (3) Phải dùng đúng định luật cho tình huống. (4) Lò phản ứng dùng nhiên liệu phân hạch, chất làm chậm, thanh điều khiển, chất tải nhiệt và hệ che chắn để duy trì phản ứng có kiểm soát.
\shortans{3}
\loigiai{
Đối chiếu từng phát biểu với kiến thức cốt lõi: Lò phản ứng dùng nhiên liệu phân hạch, chất làm chậm, thanh điều khiển, chất tải nhiệt và hệ che chắn để duy trì phản ứng có kiểm soát. Có 3 phát biểu đúng.
}
\end{ex}

% ===== Câu 120 =====
% ID: L12C4B24-30-TN
% Mức: VD
\dangbt{Nhận biết phản ứng nhiệt hạch và điều kiện xảy ra}
\begin{ex}
% ID: L12C4B24-30-TN
% Mức: VD
Kết luận đúng khi vận dụng là về nhận biết phản ứng nhiệt hạch và điều kiện xảy ra?
\choice
{Nhiệt hạch dễ xảy ra ở nhiệt độ phòng.}
{\True Nhiệt hạch là sự kết hợp các hạt nhân nhẹ, cần nhiệt độ rất cao để thắng lực đẩy Coulomb và có thể tỏa năng lượng lớn.}
{Có thể bỏ qua điều kiện áp dụng và đơn vị trong mọi trường hợp.}
{Kết quả của bài toán không phụ thuộc dữ kiện đã cho.}
\loigiai{
Nhiệt hạch là sự kết hợp các hạt nhân nhẹ, cần nhiệt độ rất cao để thắng lực đẩy Coulomb và có thể tỏa năng lượng lớn. Vì vậy phương án $B$ đúng; các phương án còn lại trái với kiến thức hoặc điều kiện áp dụng.
}
\end{ex}

% ===== Câu 121 =====
% ID: L12C4B24-31-DS
% Mức: TH
\dangbt{Ứng dụng công nghệ hạt nhân, an toàn và phát triển bền vững}
\begin{ex}
% ID: L12C4B24-31-DS
% Mức: TH
Xét ứng dụng công nghệ hạt nhân, an toàn và phát triển bền vững (tình huống 1). Hãy xác định tính đúng, sai của bốn phát biểu sau.
\choiceTF
{\True Công nghệ hạt nhân có ứng dụng trong năng lượng, y học, nông nghiệp và công nghiệp nhưng phải đánh giá rủi ro, an toàn và chất thải.}
{Ứng dụng hạt nhân hoàn toàn không tạo ra rủi ro môi trường.}
{\True Khi giải cần đọc đúng dữ kiện, dùng hệ đơn vị thống nhất và kiểm tra điều kiện áp dụng.}
{Có thể bỏ qua dữ kiện, đơn vị và ý nghĩa vật lí của kết quả.}
\loigiai{
\begin{itemchoice}
\item (Đúng) Công nghệ hạt nhân có ứng dụng trong năng lượng, y học, nông nghiệp và công nghiệp nhưng phải đánh giá rủi ro, an toàn và chất thải. (Vì): Công nghệ hạt nhân có ứng dụng trong năng lượng, y học, nông nghiệp và công nghiệp nhưng phải đánh giá rủi ro, an toàn và chất thải. b) Sai: Ứng dụng hạt nhân hoàn toàn không tạo ra rủi ro môi trường. c) Đúng: Khi giải cần đọc đúng dữ kiện, dùng hệ đơn vị thống nhất và kiểm tra điều kiện áp dụng. d) Sai: Có thể bỏ qua dữ kiện, đơn vị và ý nghĩa vật lí của kết quả. Kiến thức cốt lõi: Công nghệ hạt nhân có ứng dụng trong năng lượng, y học, nông nghiệp và công nghiệp nhưng phải đánh giá rủi ro, an toàn và chất thải
\item (Sai) Ứng dụng hạt nhân hoàn toàn không tạo ra rủi ro môi trường.
\item (Đúng) Khi giải cần đọc đúng dữ kiện, dùng hệ đơn vị thống nhất và kiểm tra điều kiện áp dụng.
\item (Sai) Có thể bỏ qua dữ kiện, đơn vị và ý nghĩa vật lí của kết quả.
\end{itemchoice}
}
\end{ex}

% ===== Câu 122 =====
% ID: L12C4B24-31-TL
% Mức: TH
\dangbt{Điện hạt nhân, nhiên liệu hạt nhân và xử lí chất thải}
\begin{bt}
% ID: L12C4B24-31-TL
% Mức: TH
Trình bày kiến thức cốt lõi của điện hạt nhân, nhiên liệu hạt nhân và xử lí chất thải và nêu điều kiện áp dụng.
\loigiai{
Cần nêu được: Nhà máy điện hạt nhân biến năng lượng hạt nhân thành nhiệt, cơ năng rồi điện năng; nhiên liệu và chất thải phải được quản lí nghiêm ngặt. Khi vận dụng phải xác định đúng dữ kiện, điều kiện áp dụng, hệ đơn vị và kiểm tra tính hợp lí của kết quả. Nhận định “Chất thải phóng xạ có thể xử lí như rác sinh hoạt thông thường.” sai vì trái với kiến thức trên.
}
\end{bt}

% ===== Câu 123 =====
% ID: L12C4B24-31-TLN
% Mức: TH
\begin{ex}
% ID: L12C4B24-31-TLN
% Mức: TH
Trong bốn phát biểu sau về điện hạt nhân, nhiên liệu hạt nhân và xử lí chất thải, có bao nhiêu phát biểu đúng? Chỉ ghi một số nguyên. (1) Nhà máy điện hạt nhân biến năng lượng hạt nhân thành nhiệt, cơ năng rồi điện năng; nhiên liệu và chất thải phải được quản lí nghiêm ngặt. (2) Chất thải phóng xạ có thể xử lí như rác sinh hoạt thông thường. (3) Cần kiểm tra điều kiện áp dụng trước khi tính toán. (4) Có thể bỏ qua đơn vị của kết quả.
\shortans{2}
\loigiai{
Đối chiếu từng phát biểu với kiến thức cốt lõi: Nhà máy điện hạt nhân biến năng lượng hạt nhân thành nhiệt, cơ năng rồi điện năng; nhiên liệu và chất thải phải được quản lí nghiêm ngặt. Có 2 phát biểu đúng.
}
\end{ex}

% ===== Câu 124 =====
% ID: L12C4B24-31-TN
% Mức: NB
\dangbt{Tính năng lượng trong phân hạch và nhiệt hạch}
\begin{ex}
% ID: L12C4B24-31-TN
% Mức: NB
Phát biểu nào sau đây đúng về tính năng lượng trong phân hạch và nhiệt hạch?
\choice
{\True Năng lượng toàn phần bằng năng lượng mỗi phản ứng nhân số phản ứng; số phản ứng được suy từ số hạt tham gia.}
{Năng lượng toàn phần không phụ thuộc số hạt đã phản ứng.}
{Có thể bỏ qua điều kiện áp dụng và đơn vị trong mọi trường hợp.}
{Kết quả của bài toán không phụ thuộc dữ kiện đã cho.}
\loigiai{
Năng lượng toàn phần bằng năng lượng mỗi phản ứng nhân số phản ứng; số phản ứng được suy từ số hạt tham gia. Vì vậy phương án $A$ đúng; các phương án còn lại trái với kiến thức hoặc điều kiện áp dụng.
}
\end{ex}

% ===== Câu 125 =====
% ID: L12C4B24-32-DS
% Mức: TH
\dangbt{Ứng dụng công nghệ hạt nhân, an toàn và phát triển bền vững}
\begin{ex}
% ID: L12C4B24-32-DS
% Mức: TH
Xét ứng dụng công nghệ hạt nhân, an toàn và phát triển bền vững (tình huống 2). Hãy xác định tính đúng, sai của bốn phát biểu sau.
\choiceTF
{Ứng dụng hạt nhân hoàn toàn không tạo ra rủi ro môi trường.}
{\True Công nghệ hạt nhân có ứng dụng trong năng lượng, y học, nông nghiệp và công nghiệp nhưng phải đánh giá rủi ro, an toàn và chất thải.}
{Có thể bỏ qua dữ kiện, đơn vị và ý nghĩa vật lí của kết quả.}
{\True Khi giải cần đọc đúng dữ kiện, dùng hệ đơn vị thống nhất và kiểm tra điều kiện áp dụng.}
\loigiai{
\begin{itemchoice}
\item (Sai) Ứng dụng hạt nhân hoàn toàn không tạo ra rủi ro môi trường. (Vì): Ứng dụng hạt nhân hoàn toàn không tạo ra rủi ro môi trường. b) Đúng: Công nghệ hạt nhân có ứng dụng trong năng lượng, y học, nông nghiệp và công nghiệp nhưng phải đánh giá rủi ro, an toàn và chất thải. c) Sai: Có thể bỏ qua dữ kiện, đơn vị và ý nghĩa vật lí của kết quả. d) Đúng: Khi giải cần đọc đúng dữ kiện, dùng hệ đơn vị thống nhất và kiểm tra điều kiện áp dụng. Kiến thức cốt lõi: Công nghệ hạt nhân có ứng dụng trong năng lượng, y học, nông nghiệp và công nghiệp nhưng phải đánh giá rủi ro, an toàn và chất thải
\item (Đúng) Công nghệ hạt nhân có ứng dụng trong năng lượng, y học, nông nghiệp và công nghiệp nhưng phải đánh giá rủi ro, an toàn và chất thải.
\item (Sai) Có thể bỏ qua dữ kiện, đơn vị và ý nghĩa vật lí của kết quả.
\item (Đúng) Khi giải cần đọc đúng dữ kiện, dùng hệ đơn vị thống nhất và kiểm tra điều kiện áp dụng.
\end{itemchoice}
}
\end{ex}

% ===== Câu 126 =====
% ID: L12C4B24-32-TL
% Mức: TH
\dangbt{Điện hạt nhân, nhiên liệu hạt nhân và xử lí chất thải}
\begin{bt}
% ID: L12C4B24-32-TL
% Mức: TH
Giải thích vì sao nhận định sau đúng: “Nhà máy điện hạt nhân biến năng lượng hạt nhân thành nhiệt, cơ năng rồi điện năng; nhiên liệu và chất thải phải được quản lí nghiêm ngặt.”
\loigiai{
Cần nêu được: Nhà máy điện hạt nhân biến năng lượng hạt nhân thành nhiệt, cơ năng rồi điện năng; nhiên liệu và chất thải phải được quản lí nghiêm ngặt. Khi vận dụng phải xác định đúng dữ kiện, điều kiện áp dụng, hệ đơn vị và kiểm tra tính hợp lí của kết quả. Nhận định “Chất thải phóng xạ có thể xử lí như rác sinh hoạt thông thường.” sai vì trái với kiến thức trên.
}
\end{bt}

% ===== Câu 127 =====
% ID: L12C4B24-32-TLN
% Mức: TH
\begin{ex}
% ID: L12C4B24-32-TLN
% Mức: TH
Trong bốn phát biểu sau về điện hạt nhân, nhiên liệu hạt nhân và xử lí chất thải, có bao nhiêu phát biểu đúng? Chỉ ghi một số nguyên. (1) Chất thải phóng xạ có thể xử lí như rác sinh hoạt thông thường. (2) Nhà máy điện hạt nhân biến năng lượng hạt nhân thành nhiệt, cơ năng rồi điện năng; nhiên liệu và chất thải phải được quản lí nghiêm ngặt. (3) Kết quả phải phù hợp ý nghĩa vật lí. (4) Mọi đại lượng đều có cùng bản chất.
\shortans{2}
\loigiai{
Đối chiếu từng phát biểu với kiến thức cốt lõi: Nhà máy điện hạt nhân biến năng lượng hạt nhân thành nhiệt, cơ năng rồi điện năng; nhiên liệu và chất thải phải được quản lí nghiêm ngặt. Có 2 phát biểu đúng.
}
\end{ex}

% ===== Câu 128 =====
% ID: L12C4B24-32-TN
% Mức: NB
\dangbt{Tính năng lượng trong phân hạch và nhiệt hạch}
\begin{ex}
% ID: L12C4B24-32-TN
% Mức: NB
Nhận định nào phù hợp với kiến thức Vật lí về tính năng lượng trong phân hạch và nhiệt hạch?
\choice
{Năng lượng toàn phần không phụ thuộc số hạt đã phản ứng.}
{\True Năng lượng toàn phần bằng năng lượng mỗi phản ứng nhân số phản ứng; số phản ứng được suy từ số hạt tham gia.}
{Có thể bỏ qua điều kiện áp dụng và đơn vị trong mọi trường hợp.}
{Kết quả của bài toán không phụ thuộc dữ kiện đã cho.}
\loigiai{
Năng lượng toàn phần bằng năng lượng mỗi phản ứng nhân số phản ứng; số phản ứng được suy từ số hạt tham gia. Vì vậy phương án $B$ đúng; các phương án còn lại trái với kiến thức hoặc điều kiện áp dụng.
}
\end{ex}

% ===== Câu 129 =====
% ID: L12C4B24-33-DS
% Mức: VD
\dangbt{Ứng dụng công nghệ hạt nhân, an toàn và phát triển bền vững}
\begin{ex}
% ID: L12C4B24-33-DS
% Mức: VD
Xét ứng dụng công nghệ hạt nhân, an toàn và phát triển bền vững (tình huống 3). Hãy xác định tính đúng, sai của bốn phát biểu sau.
\choiceTF
{\True Công nghệ hạt nhân có ứng dụng trong năng lượng, y học, nông nghiệp và công nghiệp nhưng phải đánh giá rủi ro, an toàn và chất thải.}
{Ứng dụng hạt nhân hoàn toàn không tạo ra rủi ro môi trường.}
{\True Khi giải cần đọc đúng dữ kiện, dùng hệ đơn vị thống nhất và kiểm tra điều kiện áp dụng.}
{Có thể bỏ qua dữ kiện, đơn vị và ý nghĩa vật lí của kết quả.}
\loigiai{
\begin{itemchoice}
\item (Đúng) Công nghệ hạt nhân có ứng dụng trong năng lượng, y học, nông nghiệp và công nghiệp nhưng phải đánh giá rủi ro, an toàn và chất thải. (Vì): Công nghệ hạt nhân có ứng dụng trong năng lượng, y học, nông nghiệp và công nghiệp nhưng phải đánh giá rủi ro, an toàn và chất thải. b) Sai: Ứng dụng hạt nhân hoàn toàn không tạo ra rủi ro môi trường. c) Đúng: Khi giải cần đọc đúng dữ kiện, dùng hệ đơn vị thống nhất và kiểm tra điều kiện áp dụng. d) Sai: Có thể bỏ qua dữ kiện, đơn vị và ý nghĩa vật lí của kết quả. Kiến thức cốt lõi: Công nghệ hạt nhân có ứng dụng trong năng lượng, y học, nông nghiệp và công nghiệp nhưng phải đánh giá rủi ro, an toàn và chất thải
\item (Sai) Ứng dụng hạt nhân hoàn toàn không tạo ra rủi ro môi trường.
\item (Đúng) Khi giải cần đọc đúng dữ kiện, dùng hệ đơn vị thống nhất và kiểm tra điều kiện áp dụng.
\item (Sai) Có thể bỏ qua dữ kiện, đơn vị và ý nghĩa vật lí của kết quả.
\end{itemchoice}
}
\end{ex}

% ===== Câu 130 =====
% ID: L12C4B24-33-TL
% Mức: VD
\dangbt{Điện hạt nhân, nhiên liệu hạt nhân và xử lí chất thải}
\begin{bt}
% ID: L12C4B24-33-TL
% Mức: VD
Phân tích sai lầm trong nhận định: “Chất thải phóng xạ có thể xử lí như rác sinh hoạt thông thường.”
\loigiai{
Cần nêu được: Nhà máy điện hạt nhân biến năng lượng hạt nhân thành nhiệt, cơ năng rồi điện năng; nhiên liệu và chất thải phải được quản lí nghiêm ngặt. Khi vận dụng phải xác định đúng dữ kiện, điều kiện áp dụng, hệ đơn vị và kiểm tra tính hợp lí của kết quả. Nhận định “Chất thải phóng xạ có thể xử lí như rác sinh hoạt thông thường.” sai vì trái với kiến thức trên.
}
\end{bt}

% ===== Câu 131 =====
% ID: L12C4B24-33-TLN
% Mức: TH
\begin{ex}
% ID: L12C4B24-33-TLN
% Mức: TH
Trong bốn phát biểu sau về điện hạt nhân, nhiên liệu hạt nhân và xử lí chất thải, có bao nhiêu phát biểu đúng? Chỉ ghi một số nguyên. (1) Cần đổi các đại lượng về hệ đơn vị thống nhất. (2) Nhà máy điện hạt nhân biến năng lượng hạt nhân thành nhiệt, cơ năng rồi điện năng; nhiên liệu và chất thải phải được quản lí nghiêm ngặt. (3) Chất thải phóng xạ có thể xử lí như rác sinh hoạt thông thường. (4) Không cần kiểm tra dữ kiện.
\shortans{2}
\loigiai{
Đối chiếu từng phát biểu với kiến thức cốt lõi: Nhà máy điện hạt nhân biến năng lượng hạt nhân thành nhiệt, cơ năng rồi điện năng; nhiên liệu và chất thải phải được quản lí nghiêm ngặt. Có 2 phát biểu đúng.
}
\end{ex}

% ===== Câu 132 =====
% ID: L12C4B24-33-TN
% Mức: NB
\dangbt{Tính năng lượng trong phân hạch và nhiệt hạch}
\begin{ex}
% ID: L12C4B24-33-TN
% Mức: NB
Chọn kết luận chính xác về tính năng lượng trong phân hạch và nhiệt hạch?
\choice
{Năng lượng toàn phần không phụ thuộc số hạt đã phản ứng.}
{Có thể bỏ qua điều kiện áp dụng và đơn vị trong mọi trường hợp.}
{\True Năng lượng toàn phần bằng năng lượng mỗi phản ứng nhân số phản ứng; số phản ứng được suy từ số hạt tham gia.}
{Kết quả của bài toán không phụ thuộc dữ kiện đã cho.}
\loigiai{
Năng lượng toàn phần bằng năng lượng mỗi phản ứng nhân số phản ứng; số phản ứng được suy từ số hạt tham gia. Vì vậy phương án $C$ đúng; các phương án còn lại trái với kiến thức hoặc điều kiện áp dụng.
}
\end{ex}

% ===== Câu 133 =====
% ID: L12C4B24-34-DS
% Mức: VD
\dangbt{Ứng dụng công nghệ hạt nhân, an toàn và phát triển bền vững}
\begin{ex}
% ID: L12C4B24-34-DS
% Mức: VD
Xét ứng dụng công nghệ hạt nhân, an toàn và phát triển bền vững (tình huống 4). Hãy xác định tính đúng, sai của bốn phát biểu sau.
\choiceTF
{Ứng dụng hạt nhân hoàn toàn không tạo ra rủi ro môi trường.}
{\True Công nghệ hạt nhân có ứng dụng trong năng lượng, y học, nông nghiệp và công nghiệp nhưng phải đánh giá rủi ro, an toàn và chất thải.}
{Có thể bỏ qua dữ kiện, đơn vị và ý nghĩa vật lí của kết quả.}
{\True Khi giải cần đọc đúng dữ kiện, dùng hệ đơn vị thống nhất và kiểm tra điều kiện áp dụng.}
\loigiai{
\begin{itemchoice}
\item (Sai) Ứng dụng hạt nhân hoàn toàn không tạo ra rủi ro môi trường. (Vì): Ứng dụng hạt nhân hoàn toàn không tạo ra rủi ro môi trường. b) Đúng: Công nghệ hạt nhân có ứng dụng trong năng lượng, y học, nông nghiệp và công nghiệp nhưng phải đánh giá rủi ro, an toàn và chất thải. c) Sai: Có thể bỏ qua dữ kiện, đơn vị và ý nghĩa vật lí của kết quả. d) Đúng: Khi giải cần đọc đúng dữ kiện, dùng hệ đơn vị thống nhất và kiểm tra điều kiện áp dụng. Kiến thức cốt lõi: Công nghệ hạt nhân có ứng dụng trong năng lượng, y học, nông nghiệp và công nghiệp nhưng phải đánh giá rủi ro, an toàn và chất thải
\item (Đúng) Công nghệ hạt nhân có ứng dụng trong năng lượng, y học, nông nghiệp và công nghiệp nhưng phải đánh giá rủi ro, an toàn và chất thải.
\item (Sai) Có thể bỏ qua dữ kiện, đơn vị và ý nghĩa vật lí của kết quả.
\item (Đúng) Khi giải cần đọc đúng dữ kiện, dùng hệ đơn vị thống nhất và kiểm tra điều kiện áp dụng.
\end{itemchoice}
}
\end{ex}

% ===== Câu 134 =====
% ID: L12C4B24-34-TL
% Mức: VD
\dangbt{Điện hạt nhân, nhiên liệu hạt nhân và xử lí chất thải}
\begin{bt}
% ID: L12C4B24-34-TL
% Mức: VD
Đề xuất các bước giải một bài toán thuộc dạng điện hạt nhân, nhiên liệu hạt nhân và xử lí chất thải.
\loigiai{
Cần nêu được: Nhà máy điện hạt nhân biến năng lượng hạt nhân thành nhiệt, cơ năng rồi điện năng; nhiên liệu và chất thải phải được quản lí nghiêm ngặt. Khi vận dụng phải xác định đúng dữ kiện, điều kiện áp dụng, hệ đơn vị và kiểm tra tính hợp lí của kết quả. Nhận định “Chất thải phóng xạ có thể xử lí như rác sinh hoạt thông thường.” sai vì trái với kiến thức trên.
}
\end{bt}

% ===== Câu 135 =====
% ID: L12C4B24-34-TLN
% Mức: VD
\begin{ex}
% ID: L12C4B24-34-TLN
% Mức: VD
Trong bốn phát biểu sau về điện hạt nhân, nhiên liệu hạt nhân và xử lí chất thải, có bao nhiêu phát biểu đúng? Chỉ ghi một số nguyên. (1) Nhà máy điện hạt nhân biến năng lượng hạt nhân thành nhiệt, cơ năng rồi điện năng; nhiên liệu và chất thải phải được quản lí nghiêm ngặt. (2) Cần đối chiếu kết quả với điều kiện bài toán. (3) Chất thải phóng xạ có thể xử lí như rác sinh hoạt thông thường. (4) Mọi trường hợp đều cho cùng một kết quả.
\shortans{2}
\loigiai{
Đối chiếu từng phát biểu với kiến thức cốt lõi: Nhà máy điện hạt nhân biến năng lượng hạt nhân thành nhiệt, cơ năng rồi điện năng; nhiên liệu và chất thải phải được quản lí nghiêm ngặt. Có 2 phát biểu đúng.
}
\end{ex}

% ===== Câu 136 =====
% ID: L12C4B24-34-TN
% Mức: TH
\dangbt{Tính năng lượng trong phân hạch và nhiệt hạch}
\begin{ex}
% ID: L12C4B24-34-TN
% Mức: TH
Nội dung nào mô tả đúng nhất về tính năng lượng trong phân hạch và nhiệt hạch?
\choice
{Năng lượng toàn phần không phụ thuộc số hạt đã phản ứng.}
{Có thể bỏ qua điều kiện áp dụng và đơn vị trong mọi trường hợp.}
{Kết quả của bài toán không phụ thuộc dữ kiện đã cho.}
{\True Năng lượng toàn phần bằng năng lượng mỗi phản ứng nhân số phản ứng; số phản ứng được suy từ số hạt tham gia.}
\loigiai{
Năng lượng toàn phần bằng năng lượng mỗi phản ứng nhân số phản ứng; số phản ứng được suy từ số hạt tham gia. Vì vậy phương án $D$ đúng; các phương án còn lại trái với kiến thức hoặc điều kiện áp dụng.
}
\end{ex}

% ===== Câu 137 =====
% ID: L12C4B24-35-DS
% Mức: VD
\dangbt{Ứng dụng công nghệ hạt nhân, an toàn và phát triển bền vững}
\begin{ex}
% ID: L12C4B24-35-DS
% Mức: VD
Xét ứng dụng công nghệ hạt nhân, an toàn và phát triển bền vững (tình huống 5). Hãy xác định tính đúng, sai của bốn phát biểu sau.
\choiceTF
{\True Công nghệ hạt nhân có ứng dụng trong năng lượng, y học, nông nghiệp và công nghiệp nhưng phải đánh giá rủi ro, an toàn và chất thải.}
{Có thể bỏ qua dữ kiện, đơn vị và ý nghĩa vật lí của kết quả.}
{Ứng dụng hạt nhân hoàn toàn không tạo ra rủi ro môi trường.}
{Có thể bỏ qua dữ kiện, đơn vị và ý nghĩa vật lí của kết quả.}
\loigiai{
\begin{itemchoice}
\item (Đúng) Công nghệ hạt nhân có ứng dụng trong năng lượng, y học, nông nghiệp và công nghiệp nhưng phải đánh giá rủi ro, an toàn và chất thải. (Vì): Công nghệ hạt nhân có ứng dụng trong năng lượng, y học, nông nghiệp và công nghiệp nhưng phải đánh giá rủi ro, an toàn và chất thải. b) Sai: Có thể bỏ qua dữ kiện, đơn vị và ý nghĩa vật lí của kết quả. c) Sai: Ứng dụng hạt nhân hoàn toàn không tạo ra rủi ro môi trường. d) Sai: Có thể bỏ qua dữ kiện, đơn vị và ý nghĩa vật lí của kết quả. Kiến thức cốt lõi: Công nghệ hạt nhân có ứng dụng trong năng lượng, y học, nông nghiệp và công nghiệp nhưng phải đánh giá rủi ro, an toàn và chất thải
\item (Sai) Có thể bỏ qua dữ kiện, đơn vị và ý nghĩa vật lí của kết quả.
\item (Sai) Ứng dụng hạt nhân hoàn toàn không tạo ra rủi ro môi trường.
\item (Sai) Có thể bỏ qua dữ kiện, đơn vị và ý nghĩa vật lí của kết quả.
\end{itemchoice}
}
\end{ex}

% ===== Câu 138 =====
% ID: L12C4B24-35-TL
% Mức: VD
\dangbt{Điện hạt nhân, nhiên liệu hạt nhân và xử lí chất thải}
\begin{bt}
% ID: L12C4B24-35-TL
% Mức: VD
Nêu một tình huống thực tiễn liên quan đến điện hạt nhân, nhiên liệu hạt nhân và xử lí chất thải và giải thích bằng kiến thức Vật lí.
\loigiai{
Cần nêu được: Nhà máy điện hạt nhân biến năng lượng hạt nhân thành nhiệt, cơ năng rồi điện năng; nhiên liệu và chất thải phải được quản lí nghiêm ngặt. Khi vận dụng phải xác định đúng dữ kiện, điều kiện áp dụng, hệ đơn vị và kiểm tra tính hợp lí của kết quả. Nhận định “Chất thải phóng xạ có thể xử lí như rác sinh hoạt thông thường.” sai vì trái với kiến thức trên.
}
\end{bt}

% ===== Câu 139 =====
% ID: L12C4B24-35-TLN
% Mức: VD
\begin{ex}
% ID: L12C4B24-35-TLN
% Mức: VD
Trong bốn phát biểu sau về điện hạt nhân, nhiên liệu hạt nhân và xử lí chất thải, có bao nhiêu phát biểu đúng? Chỉ ghi một số nguyên. (1) Chất thải phóng xạ có thể xử lí như rác sinh hoạt thông thường. (2) Cần đọc đúng dữ kiện và giả thiết. (3) Nhà máy điện hạt nhân biến năng lượng hạt nhân thành nhiệt, cơ năng rồi điện năng; nhiên liệu và chất thải phải được quản lí nghiêm ngặt. (4) Có thể thay số trước khi chọn công thức.
\shortans{2}
\loigiai{
Đối chiếu từng phát biểu với kiến thức cốt lõi: Nhà máy điện hạt nhân biến năng lượng hạt nhân thành nhiệt, cơ năng rồi điện năng; nhiên liệu và chất thải phải được quản lí nghiêm ngặt. Có 2 phát biểu đúng.
}
\end{ex}

% ===== Câu 140 =====
% ID: L12C4B24-35-TN
% Mức: TH
\dangbt{Tính năng lượng trong phân hạch và nhiệt hạch}
\begin{ex}
% ID: L12C4B24-35-TN
% Mức: TH
Khi phân tích, phát biểu nào đúng về tính năng lượng trong phân hạch và nhiệt hạch?
\choice
{\True Năng lượng toàn phần bằng năng lượng mỗi phản ứng nhân số phản ứng; số phản ứng được suy từ số hạt tham gia.}
{Năng lượng toàn phần không phụ thuộc số hạt đã phản ứng.}
{Có thể bỏ qua điều kiện áp dụng và đơn vị trong mọi trường hợp.}
{Kết quả của bài toán không phụ thuộc dữ kiện đã cho.}
\loigiai{
Năng lượng toàn phần bằng năng lượng mỗi phản ứng nhân số phản ứng; số phản ứng được suy từ số hạt tham gia. Vì vậy phương án $A$ đúng; các phương án còn lại trái với kiến thức hoặc điều kiện áp dụng.
}
\end{ex}

% ===== Câu 141 =====
% ID: L12C4B24-36-TL
% Mức: VDC
\dangbt{Điện hạt nhân, nhiên liệu hạt nhân và xử lí chất thải}
\begin{bt}
% ID: L12C4B24-36-TL
% Mức: VDC
So sánh nhận định đúng “Nhà máy điện hạt nhân biến năng lượng hạt nhân thành nhiệt, cơ năng rồi điện năng; nhiên liệu và chất thải phải được quản lí nghiêm ngặt.” với nhận định sai “Chất thải phóng xạ có thể xử lí như rác sinh hoạt thông thường.”; chỉ rõ căn cứ Vật lí.
\loigiai{
Cần nêu được: Nhà máy điện hạt nhân biến năng lượng hạt nhân thành nhiệt, cơ năng rồi điện năng; nhiên liệu và chất thải phải được quản lí nghiêm ngặt. Khi vận dụng phải xác định đúng dữ kiện, điều kiện áp dụng, hệ đơn vị và kiểm tra tính hợp lí của kết quả. Nhận định “Chất thải phóng xạ có thể xử lí như rác sinh hoạt thông thường.” sai vì trái với kiến thức trên.
}
\end{bt}

% ===== Câu 142 =====
% ID: L12C4B24-36-TLN
% Mức: VDC
\begin{ex}
% ID: L12C4B24-36-TLN
% Mức: VDC
Trong bốn phát biểu sau về điện hạt nhân, nhiên liệu hạt nhân và xử lí chất thải, có bao nhiêu phát biểu đúng? Chỉ ghi một số nguyên. (1) Kết quả cần có ý nghĩa vật lí. (2) Chất thải phóng xạ có thể xử lí như rác sinh hoạt thông thường. (3) Phải dùng đúng định luật cho tình huống. (4) Nhà máy điện hạt nhân biến năng lượng hạt nhân thành nhiệt, cơ năng rồi điện năng; nhiên liệu và chất thải phải được quản lí nghiêm ngặt.
\shortans{3}
\loigiai{
Đối chiếu từng phát biểu với kiến thức cốt lõi: Nhà máy điện hạt nhân biến năng lượng hạt nhân thành nhiệt, cơ năng rồi điện năng; nhiên liệu và chất thải phải được quản lí nghiêm ngặt. Có 3 phát biểu đúng.
}
\end{ex}

% ===== Câu 143 =====
% ID: L12C4B24-36-TN
% Mức: TH
\dangbt{Tính năng lượng trong phân hạch và nhiệt hạch}
\begin{ex}
% ID: L12C4B24-36-TN
% Mức: TH
Kết luận nào của học sinh là đúng về tính năng lượng trong phân hạch và nhiệt hạch?
\choice
{Năng lượng toàn phần không phụ thuộc số hạt đã phản ứng.}
{\True Năng lượng toàn phần bằng năng lượng mỗi phản ứng nhân số phản ứng; số phản ứng được suy từ số hạt tham gia.}
{Có thể bỏ qua điều kiện áp dụng và đơn vị trong mọi trường hợp.}
{Kết quả của bài toán không phụ thuộc dữ kiện đã cho.}
\loigiai{
Năng lượng toàn phần bằng năng lượng mỗi phản ứng nhân số phản ứng; số phản ứng được suy từ số hạt tham gia. Vì vậy phương án $B$ đúng; các phương án còn lại trái với kiến thức hoặc điều kiện áp dụng.
}
\end{ex}

% ===== Câu 144 =====
% ID: L12C4B24-37-TL
% Mức: TH
\dangbt{Ứng dụng công nghệ hạt nhân, an toàn và phát triển bền vững}
\begin{bt}
% ID: L12C4B24-37-TL
% Mức: TH
Trình bày kiến thức cốt lõi của ứng dụng công nghệ hạt nhân, an toàn và phát triển bền vững và nêu điều kiện áp dụng.
\loigiai{
Cần nêu được: Công nghệ hạt nhân có ứng dụng trong năng lượng, y học, nông nghiệp và công nghiệp nhưng phải đánh giá rủi ro, an toàn và chất thải. Khi vận dụng phải xác định đúng dữ kiện, điều kiện áp dụng, hệ đơn vị và kiểm tra tính hợp lí của kết quả. Nhận định “Ứng dụng hạt nhân hoàn toàn không tạo ra rủi ro môi trường.” sai vì trái với kiến thức trên.
}
\end{bt}

% ===== Câu 145 =====
% ID: L12C4B24-37-TLN
% Mức: TH
\begin{ex}
% ID: L12C4B24-37-TLN
% Mức: TH
Trong bốn phát biểu sau về ứng dụng công nghệ hạt nhân, an toàn và phát triển bền vững, có bao nhiêu phát biểu đúng? Chỉ ghi một số nguyên. (1) Công nghệ hạt nhân có ứng dụng trong năng lượng, y học, nông nghiệp và công nghiệp nhưng phải đánh giá rủi ro, an toàn và chất thải. (2) Ứng dụng hạt nhân hoàn toàn không tạo ra rủi ro môi trường. (3) Cần kiểm tra điều kiện áp dụng trước khi tính toán. (4) Có thể bỏ qua đơn vị của kết quả.
\shortans{2}
\loigiai{
Đối chiếu từng phát biểu với kiến thức cốt lõi: Công nghệ hạt nhân có ứng dụng trong năng lượng, y học, nông nghiệp và công nghiệp nhưng phải đánh giá rủi ro, an toàn và chất thải. Có 2 phát biểu đúng.
}
\end{ex}

% ===== Câu 146 =====
% ID: L12C4B24-37-TN
% Mức: TH
\dangbt{Tính năng lượng trong phân hạch và nhiệt hạch}
\begin{ex}
% ID: L12C4B24-37-TN
% Mức: TH
Chọn phát biểu không trái với kiến thức về tính năng lượng trong phân hạch và nhiệt hạch?
\choice
{Năng lượng toàn phần không phụ thuộc số hạt đã phản ứng.}
{Có thể bỏ qua điều kiện áp dụng và đơn vị trong mọi trường hợp.}
{\True Năng lượng toàn phần bằng năng lượng mỗi phản ứng nhân số phản ứng; số phản ứng được suy từ số hạt tham gia.}
{Kết quả của bài toán không phụ thuộc dữ kiện đã cho.}
\loigiai{
Năng lượng toàn phần bằng năng lượng mỗi phản ứng nhân số phản ứng; số phản ứng được suy từ số hạt tham gia. Vì vậy phương án $C$ đúng; các phương án còn lại trái với kiến thức hoặc điều kiện áp dụng.
}
\end{ex}

% ===== Câu 147 =====
% ID: L12C4B24-38-TL
% Mức: TH
\dangbt{Ứng dụng công nghệ hạt nhân, an toàn và phát triển bền vững}
\begin{bt}
% ID: L12C4B24-38-TL
% Mức: TH
Giải thích vì sao nhận định sau đúng: “Công nghệ hạt nhân có ứng dụng trong năng lượng, y học, nông nghiệp và công nghiệp nhưng phải đánh giá rủi ro, an toàn và chất thải.”
\loigiai{
Cần nêu được: Công nghệ hạt nhân có ứng dụng trong năng lượng, y học, nông nghiệp và công nghiệp nhưng phải đánh giá rủi ro, an toàn và chất thải. Khi vận dụng phải xác định đúng dữ kiện, điều kiện áp dụng, hệ đơn vị và kiểm tra tính hợp lí của kết quả. Nhận định “Ứng dụng hạt nhân hoàn toàn không tạo ra rủi ro môi trường.” sai vì trái với kiến thức trên.
}
\end{bt}

% ===== Câu 148 =====
% ID: L12C4B24-38-TLN
% Mức: TH
\begin{ex}
% ID: L12C4B24-38-TLN
% Mức: TH
Trong bốn phát biểu sau về ứng dụng công nghệ hạt nhân, an toàn và phát triển bền vững, có bao nhiêu phát biểu đúng? Chỉ ghi một số nguyên. (1) Ứng dụng hạt nhân hoàn toàn không tạo ra rủi ro môi trường. (2) Công nghệ hạt nhân có ứng dụng trong năng lượng, y học, nông nghiệp và công nghiệp nhưng phải đánh giá rủi ro, an toàn và chất thải. (3) Kết quả phải phù hợp ý nghĩa vật lí. (4) Mọi đại lượng đều có cùng bản chất.
\shortans{2}
\loigiai{
Đối chiếu từng phát biểu với kiến thức cốt lõi: Công nghệ hạt nhân có ứng dụng trong năng lượng, y học, nông nghiệp và công nghiệp nhưng phải đánh giá rủi ro, an toàn và chất thải. Có 2 phát biểu đúng.
}
\end{ex}

% ===== Câu 149 =====
% ID: L12C4B24-38-TN
% Mức: VD
\dangbt{Tính năng lượng trong phân hạch và nhiệt hạch}
\begin{ex}
% ID: L12C4B24-38-TN
% Mức: VD
Cơ sở Vật lí đúng để giải bài là gì về tính năng lượng trong phân hạch và nhiệt hạch?
\choice
{Năng lượng toàn phần không phụ thuộc số hạt đã phản ứng.}
{Có thể bỏ qua điều kiện áp dụng và đơn vị trong mọi trường hợp.}
{Kết quả của bài toán không phụ thuộc dữ kiện đã cho.}
{\True Năng lượng toàn phần bằng năng lượng mỗi phản ứng nhân số phản ứng; số phản ứng được suy từ số hạt tham gia.}
\loigiai{
Năng lượng toàn phần bằng năng lượng mỗi phản ứng nhân số phản ứng; số phản ứng được suy từ số hạt tham gia. Vì vậy phương án $D$ đúng; các phương án còn lại trái với kiến thức hoặc điều kiện áp dụng.
}
\end{ex}

% ===== Câu 150 =====
% ID: L12C4B24-39-TL
% Mức: VD
\dangbt{Ứng dụng công nghệ hạt nhân, an toàn và phát triển bền vững}
\begin{bt}
% ID: L12C4B24-39-TL
% Mức: VD
Phân tích sai lầm trong nhận định: “Ứng dụng hạt nhân hoàn toàn không tạo ra rủi ro môi trường.”
\loigiai{
Cần nêu được: Công nghệ hạt nhân có ứng dụng trong năng lượng, y học, nông nghiệp và công nghiệp nhưng phải đánh giá rủi ro, an toàn và chất thải. Khi vận dụng phải xác định đúng dữ kiện, điều kiện áp dụng, hệ đơn vị và kiểm tra tính hợp lí của kết quả. Nhận định “Ứng dụng hạt nhân hoàn toàn không tạo ra rủi ro môi trường.” sai vì trái với kiến thức trên.
}
\end{bt}

% ===== Câu 151 =====
% ID: L12C4B24-39-TLN
% Mức: TH
\begin{ex}
% ID: L12C4B24-39-TLN
% Mức: TH
Trong bốn phát biểu sau về ứng dụng công nghệ hạt nhân, an toàn và phát triển bền vững, có bao nhiêu phát biểu đúng? Chỉ ghi một số nguyên. (1) Cần đổi các đại lượng về hệ đơn vị thống nhất. (2) Công nghệ hạt nhân có ứng dụng trong năng lượng, y học, nông nghiệp và công nghiệp nhưng phải đánh giá rủi ro, an toàn và chất thải. (3) Ứng dụng hạt nhân hoàn toàn không tạo ra rủi ro môi trường. (4) Không cần kiểm tra dữ kiện.
\shortans{2}
\loigiai{
Đối chiếu từng phát biểu với kiến thức cốt lõi: Công nghệ hạt nhân có ứng dụng trong năng lượng, y học, nông nghiệp và công nghiệp nhưng phải đánh giá rủi ro, an toàn và chất thải. Có 2 phát biểu đúng.
}
\end{ex}

% ===== Câu 152 =====
% ID: L12C4B24-39-TN
% Mức: VD
\dangbt{Tính năng lượng trong phân hạch và nhiệt hạch}
\begin{ex}
% ID: L12C4B24-39-TN
% Mức: VD
Trong một bài thực hành, nhận xét nào đúng về tính năng lượng trong phân hạch và nhiệt hạch?
\choice
{\True Năng lượng toàn phần bằng năng lượng mỗi phản ứng nhân số phản ứng; số phản ứng được suy từ số hạt tham gia.}
{Năng lượng toàn phần không phụ thuộc số hạt đã phản ứng.}
{Có thể bỏ qua điều kiện áp dụng và đơn vị trong mọi trường hợp.}
{Kết quả của bài toán không phụ thuộc dữ kiện đã cho.}
\loigiai{
Năng lượng toàn phần bằng năng lượng mỗi phản ứng nhân số phản ứng; số phản ứng được suy từ số hạt tham gia. Vì vậy phương án $A$ đúng; các phương án còn lại trái với kiến thức hoặc điều kiện áp dụng.
}
\end{ex}

% ===== Câu 153 =====
% ID: L12C4B24-40-TL
% Mức: VD
\dangbt{Ứng dụng công nghệ hạt nhân, an toàn và phát triển bền vững}
\begin{bt}
% ID: L12C4B24-40-TL
% Mức: VD
Đề xuất các bước giải một bài toán thuộc dạng ứng dụng công nghệ hạt nhân, an toàn và phát triển bền vững.
\loigiai{
Cần nêu được: Công nghệ hạt nhân có ứng dụng trong năng lượng, y học, nông nghiệp và công nghiệp nhưng phải đánh giá rủi ro, an toàn và chất thải. Khi vận dụng phải xác định đúng dữ kiện, điều kiện áp dụng, hệ đơn vị và kiểm tra tính hợp lí của kết quả. Nhận định “Ứng dụng hạt nhân hoàn toàn không tạo ra rủi ro môi trường.” sai vì trái với kiến thức trên.
}
\end{bt}

% ===== Câu 154 =====
% ID: L12C4B24-40-TLN
% Mức: VD
\begin{ex}
% ID: L12C4B24-40-TLN
% Mức: VD
Trong bốn phát biểu sau về ứng dụng công nghệ hạt nhân, an toàn và phát triển bền vững, có bao nhiêu phát biểu đúng? Chỉ ghi một số nguyên. (1) Công nghệ hạt nhân có ứng dụng trong năng lượng, y học, nông nghiệp và công nghiệp nhưng phải đánh giá rủi ro, an toàn và chất thải. (2) Cần đối chiếu kết quả với điều kiện bài toán. (3) Ứng dụng hạt nhân hoàn toàn không tạo ra rủi ro môi trường. (4) Mọi trường hợp đều cho cùng một kết quả.
\shortans{2}
\loigiai{
Đối chiếu từng phát biểu với kiến thức cốt lõi: Công nghệ hạt nhân có ứng dụng trong năng lượng, y học, nông nghiệp và công nghiệp nhưng phải đánh giá rủi ro, an toàn và chất thải. Có 2 phát biểu đúng.
}
\end{ex}

% ===== Câu 155 =====
% ID: L12C4B24-40-TN
% Mức: VD
\dangbt{Tính năng lượng trong phân hạch và nhiệt hạch}
\begin{ex}
% ID: L12C4B24-40-TN
% Mức: VD
Kết luận đúng khi vận dụng là về tính năng lượng trong phân hạch và nhiệt hạch?
\choice
{Năng lượng toàn phần không phụ thuộc số hạt đã phản ứng.}
{\True Năng lượng toàn phần bằng năng lượng mỗi phản ứng nhân số phản ứng; số phản ứng được suy từ số hạt tham gia.}
{Có thể bỏ qua điều kiện áp dụng và đơn vị trong mọi trường hợp.}
{Kết quả của bài toán không phụ thuộc dữ kiện đã cho.}
\loigiai{
Năng lượng toàn phần bằng năng lượng mỗi phản ứng nhân số phản ứng; số phản ứng được suy từ số hạt tham gia. Vì vậy phương án $B$ đúng; các phương án còn lại trái với kiến thức hoặc điều kiện áp dụng.
}
\end{ex}

% ===== Câu 156 =====
% ID: L12C4B24-41-TL
% Mức: VD
\dangbt{Ứng dụng công nghệ hạt nhân, an toàn và phát triển bền vững}
\begin{bt}
% ID: L12C4B24-41-TL
% Mức: VD
Nêu một tình huống thực tiễn liên quan đến ứng dụng công nghệ hạt nhân, an toàn và phát triển bền vững và giải thích bằng kiến thức Vật lí.
\loigiai{
Cần nêu được: Công nghệ hạt nhân có ứng dụng trong năng lượng, y học, nông nghiệp và công nghiệp nhưng phải đánh giá rủi ro, an toàn và chất thải. Khi vận dụng phải xác định đúng dữ kiện, điều kiện áp dụng, hệ đơn vị và kiểm tra tính hợp lí của kết quả. Nhận định “Ứng dụng hạt nhân hoàn toàn không tạo ra rủi ro môi trường.” sai vì trái với kiến thức trên.
}
\end{bt}

% ===== Câu 157 =====
% ID: L12C4B24-41-TLN
% Mức: VD
\begin{ex}
% ID: L12C4B24-41-TLN
% Mức: VD
Trong bốn phát biểu sau về ứng dụng công nghệ hạt nhân, an toàn và phát triển bền vững, có bao nhiêu phát biểu đúng? Chỉ ghi một số nguyên. (1) Ứng dụng hạt nhân hoàn toàn không tạo ra rủi ro môi trường. (2) Cần đọc đúng dữ kiện và giả thiết. (3) Công nghệ hạt nhân có ứng dụng trong năng lượng, y học, nông nghiệp và công nghiệp nhưng phải đánh giá rủi ro, an toàn và chất thải. (4) Có thể thay số trước khi chọn công thức.
\shortans{2}
\loigiai{
Đối chiếu từng phát biểu với kiến thức cốt lõi: Công nghệ hạt nhân có ứng dụng trong năng lượng, y học, nông nghiệp và công nghiệp nhưng phải đánh giá rủi ro, an toàn và chất thải. Có 2 phát biểu đúng.
}
\end{ex}

% ===== Câu 158 =====
% ID: L12C4B24-41-TN
% Mức: NB
\dangbt{Cấu tạo và nguyên lí hoạt động của lò phản ứng hạt nhân}
\begin{ex}
% ID: L12C4B24-41-TN
% Mức: NB
Phát biểu nào sau đây đúng về cấu tạo và nguyên lí hoạt động của lò phản ứng hạt nhân?
\choice
{\True Lò phản ứng dùng nhiên liệu phân hạch, chất làm chậm, thanh điều khiển, chất tải nhiệt và hệ che chắn để duy trì phản ứng có kiểm soát.}
{Thanh điều khiển có nhiệm vụ làm tăng không giới hạn số neutron trong lò.}
{Có thể bỏ qua điều kiện áp dụng và đơn vị trong mọi trường hợp.}
{Kết quả của bài toán không phụ thuộc dữ kiện đã cho.}
\loigiai{
Lò phản ứng dùng nhiên liệu phân hạch, chất làm chậm, thanh điều khiển, chất tải nhiệt và hệ che chắn để duy trì phản ứng có kiểm soát. Vì vậy phương án $A$ đúng; các phương án còn lại trái với kiến thức hoặc điều kiện áp dụng.
}
\end{ex}

% ===== Câu 159 =====
% ID: L12C4B24-42-TL
% Mức: VDC
\dangbt{Ứng dụng công nghệ hạt nhân, an toàn và phát triển bền vững}
\begin{bt}
% ID: L12C4B24-42-TL
% Mức: VDC
So sánh nhận định đúng “Công nghệ hạt nhân có ứng dụng trong năng lượng, y học, nông nghiệp và công nghiệp nhưng phải đánh giá rủi ro, an toàn và chất thải.” với nhận định sai “Ứng dụng hạt nhân hoàn toàn không tạo ra rủi ro môi trường.”; chỉ rõ căn cứ Vật lí.
\loigiai{
Cần nêu được: Công nghệ hạt nhân có ứng dụng trong năng lượng, y học, nông nghiệp và công nghiệp nhưng phải đánh giá rủi ro, an toàn và chất thải. Khi vận dụng phải xác định đúng dữ kiện, điều kiện áp dụng, hệ đơn vị và kiểm tra tính hợp lí của kết quả. Nhận định “Ứng dụng hạt nhân hoàn toàn không tạo ra rủi ro môi trường.” sai vì trái với kiến thức trên.
}
\end{bt}

% ===== Câu 160 =====
% ID: L12C4B24-42-TLN
% Mức: VDC
\begin{ex}
% ID: L12C4B24-42-TLN
% Mức: VDC
Trong bốn phát biểu sau về ứng dụng công nghệ hạt nhân, an toàn và phát triển bền vững, có bao nhiêu phát biểu đúng? Chỉ ghi một số nguyên. (1) Kết quả cần có ý nghĩa vật lí. (2) Ứng dụng hạt nhân hoàn toàn không tạo ra rủi ro môi trường. (3) Phải dùng đúng định luật cho tình huống. (4) Công nghệ hạt nhân có ứng dụng trong năng lượng, y học, nông nghiệp và công nghiệp nhưng phải đánh giá rủi ro, an toàn và chất thải.
\shortans{3}
\loigiai{
Đối chiếu từng phát biểu với kiến thức cốt lõi: Công nghệ hạt nhân có ứng dụng trong năng lượng, y học, nông nghiệp và công nghiệp nhưng phải đánh giá rủi ro, an toàn và chất thải. Có 3 phát biểu đúng.
}
\end{ex}

% ===== Câu 161 =====
% ID: L12C4B24-42-TN
% Mức: NB
\dangbt{Cấu tạo và nguyên lí hoạt động của lò phản ứng hạt nhân}
\begin{ex}
% ID: L12C4B24-42-TN
% Mức: NB
Nhận định nào phù hợp với kiến thức Vật lí về cấu tạo và nguyên lí hoạt động của lò phản ứng hạt nhân?
\choice
{Thanh điều khiển có nhiệm vụ làm tăng không giới hạn số neutron trong lò.}
{\True Lò phản ứng dùng nhiên liệu phân hạch, chất làm chậm, thanh điều khiển, chất tải nhiệt và hệ che chắn để duy trì phản ứng có kiểm soát.}
{Có thể bỏ qua điều kiện áp dụng và đơn vị trong mọi trường hợp.}
{Kết quả của bài toán không phụ thuộc dữ kiện đã cho.}
\loigiai{
Lò phản ứng dùng nhiên liệu phân hạch, chất làm chậm, thanh điều khiển, chất tải nhiệt và hệ che chắn để duy trì phản ứng có kiểm soát. Vì vậy phương án $B$ đúng; các phương án còn lại trái với kiến thức hoặc điều kiện áp dụng.
}
\end{ex}

% ===== Câu 162 =====
% ID: L12C4B24-43-TN
% Mức: NB
\begin{ex}
% ID: L12C4B24-43-TN
% Mức: NB
Chọn kết luận chính xác về cấu tạo và nguyên lí hoạt động của lò phản ứng hạt nhân?
\choice
{Thanh điều khiển có nhiệm vụ làm tăng không giới hạn số neutron trong lò.}
{Có thể bỏ qua điều kiện áp dụng và đơn vị trong mọi trường hợp.}
{\True Lò phản ứng dùng nhiên liệu phân hạch, chất làm chậm, thanh điều khiển, chất tải nhiệt và hệ che chắn để duy trì phản ứng có kiểm soát.}
{Kết quả của bài toán không phụ thuộc dữ kiện đã cho.}
\loigiai{
Lò phản ứng dùng nhiên liệu phân hạch, chất làm chậm, thanh điều khiển, chất tải nhiệt và hệ che chắn để duy trì phản ứng có kiểm soát. Vì vậy phương án $C$ đúng; các phương án còn lại trái với kiến thức hoặc điều kiện áp dụng.
}
\end{ex}

% ===== Câu 163 =====
% ID: L12C4B24-44-TN
% Mức: TH
\begin{ex}
% ID: L12C4B24-44-TN
% Mức: TH
Nội dung nào mô tả đúng nhất về cấu tạo và nguyên lí hoạt động của lò phản ứng hạt nhân?
\choice
{Thanh điều khiển có nhiệm vụ làm tăng không giới hạn số neutron trong lò.}
{Có thể bỏ qua điều kiện áp dụng và đơn vị trong mọi trường hợp.}
{Kết quả của bài toán không phụ thuộc dữ kiện đã cho.}
{\True Lò phản ứng dùng nhiên liệu phân hạch, chất làm chậm, thanh điều khiển, chất tải nhiệt và hệ che chắn để duy trì phản ứng có kiểm soát.}
\loigiai{
Lò phản ứng dùng nhiên liệu phân hạch, chất làm chậm, thanh điều khiển, chất tải nhiệt và hệ che chắn để duy trì phản ứng có kiểm soát. Vì vậy phương án $D$ đúng; các phương án còn lại trái với kiến thức hoặc điều kiện áp dụng.
}
\end{ex}

% ===== Câu 164 =====
% ID: L12C4B24-45-TN
% Mức: TH
\begin{ex}
% ID: L12C4B24-45-TN
% Mức: TH
Khi phân tích, phát biểu nào đúng về cấu tạo và nguyên lí hoạt động của lò phản ứng hạt nhân?
\choice
{\True Lò phản ứng dùng nhiên liệu phân hạch, chất làm chậm, thanh điều khiển, chất tải nhiệt và hệ che chắn để duy trì phản ứng có kiểm soát.}
{Thanh điều khiển có nhiệm vụ làm tăng không giới hạn số neutron trong lò.}
{Có thể bỏ qua điều kiện áp dụng và đơn vị trong mọi trường hợp.}
{Kết quả của bài toán không phụ thuộc dữ kiện đã cho.}
\loigiai{
Lò phản ứng dùng nhiên liệu phân hạch, chất làm chậm, thanh điều khiển, chất tải nhiệt và hệ che chắn để duy trì phản ứng có kiểm soát. Vì vậy phương án $A$ đúng; các phương án còn lại trái với kiến thức hoặc điều kiện áp dụng.
}
\end{ex}

% ===== Câu 165 =====
% ID: L12C4B24-46-TN
% Mức: TH
\begin{ex}
% ID: L12C4B24-46-TN
% Mức: TH
Kết luận nào của học sinh là đúng về cấu tạo và nguyên lí hoạt động của lò phản ứng hạt nhân?
\choice
{Thanh điều khiển có nhiệm vụ làm tăng không giới hạn số neutron trong lò.}
{\True Lò phản ứng dùng nhiên liệu phân hạch, chất làm chậm, thanh điều khiển, chất tải nhiệt và hệ che chắn để duy trì phản ứng có kiểm soát.}
{Có thể bỏ qua điều kiện áp dụng và đơn vị trong mọi trường hợp.}
{Kết quả của bài toán không phụ thuộc dữ kiện đã cho.}
\loigiai{
Lò phản ứng dùng nhiên liệu phân hạch, chất làm chậm, thanh điều khiển, chất tải nhiệt và hệ che chắn để duy trì phản ứng có kiểm soát. Vì vậy phương án $B$ đúng; các phương án còn lại trái với kiến thức hoặc điều kiện áp dụng.
}
\end{ex}

% ===== Câu 166 =====
% ID: L12C4B24-47-TN
% Mức: TH
\begin{ex}
% ID: L12C4B24-47-TN
% Mức: TH
Chọn phát biểu không trái với kiến thức về cấu tạo và nguyên lí hoạt động của lò phản ứng hạt nhân?
\choice
{Thanh điều khiển có nhiệm vụ làm tăng không giới hạn số neutron trong lò.}
{Có thể bỏ qua điều kiện áp dụng và đơn vị trong mọi trường hợp.}
{\True Lò phản ứng dùng nhiên liệu phân hạch, chất làm chậm, thanh điều khiển, chất tải nhiệt và hệ che chắn để duy trì phản ứng có kiểm soát.}
{Kết quả của bài toán không phụ thuộc dữ kiện đã cho.}
\loigiai{
Lò phản ứng dùng nhiên liệu phân hạch, chất làm chậm, thanh điều khiển, chất tải nhiệt và hệ che chắn để duy trì phản ứng có kiểm soát. Vì vậy phương án $C$ đúng; các phương án còn lại trái với kiến thức hoặc điều kiện áp dụng.
}
\end{ex}

% ===== Câu 167 =====
% ID: L12C4B24-48-TN
% Mức: VD
\begin{ex}
% ID: L12C4B24-48-TN
% Mức: VD
Cơ sở Vật lí đúng để giải bài là gì về cấu tạo và nguyên lí hoạt động của lò phản ứng hạt nhân?
\choice
{Thanh điều khiển có nhiệm vụ làm tăng không giới hạn số neutron trong lò.}
{Có thể bỏ qua điều kiện áp dụng và đơn vị trong mọi trường hợp.}
{Kết quả của bài toán không phụ thuộc dữ kiện đã cho.}
{\True Lò phản ứng dùng nhiên liệu phân hạch, chất làm chậm, thanh điều khiển, chất tải nhiệt và hệ che chắn để duy trì phản ứng có kiểm soát.}
\loigiai{
Lò phản ứng dùng nhiên liệu phân hạch, chất làm chậm, thanh điều khiển, chất tải nhiệt và hệ che chắn để duy trì phản ứng có kiểm soát. Vì vậy phương án $D$ đúng; các phương án còn lại trái với kiến thức hoặc điều kiện áp dụng.
}
\end{ex}

% ===== Câu 168 =====
% ID: L12C4B24-49-TN
% Mức: VD
\begin{ex}
% ID: L12C4B24-49-TN
% Mức: VD
Trong một bài thực hành, nhận xét nào đúng về cấu tạo và nguyên lí hoạt động của lò phản ứng hạt nhân?
\choice
{\True Lò phản ứng dùng nhiên liệu phân hạch, chất làm chậm, thanh điều khiển, chất tải nhiệt và hệ che chắn để duy trì phản ứng có kiểm soát.}
{Thanh điều khiển có nhiệm vụ làm tăng không giới hạn số neutron trong lò.}
{Có thể bỏ qua điều kiện áp dụng và đơn vị trong mọi trường hợp.}
{Kết quả của bài toán không phụ thuộc dữ kiện đã cho.}
\loigiai{
Lò phản ứng dùng nhiên liệu phân hạch, chất làm chậm, thanh điều khiển, chất tải nhiệt và hệ che chắn để duy trì phản ứng có kiểm soát. Vì vậy phương án $A$ đúng; các phương án còn lại trái với kiến thức hoặc điều kiện áp dụng.
}
\end{ex}

% ===== Câu 169 =====
% ID: L12C4B24-50-TN
% Mức: VD
\begin{ex}
% ID: L12C4B24-50-TN
% Mức: VD
Kết luận đúng khi vận dụng là về cấu tạo và nguyên lí hoạt động của lò phản ứng hạt nhân?
\choice
{Thanh điều khiển có nhiệm vụ làm tăng không giới hạn số neutron trong lò.}
{\True Lò phản ứng dùng nhiên liệu phân hạch, chất làm chậm, thanh điều khiển, chất tải nhiệt và hệ che chắn để duy trì phản ứng có kiểm soát.}
{Có thể bỏ qua điều kiện áp dụng và đơn vị trong mọi trường hợp.}
{Kết quả của bài toán không phụ thuộc dữ kiện đã cho.}
\loigiai{
Lò phản ứng dùng nhiên liệu phân hạch, chất làm chậm, thanh điều khiển, chất tải nhiệt và hệ che chắn để duy trì phản ứng có kiểm soát. Vì vậy phương án $B$ đúng; các phương án còn lại trái với kiến thức hoặc điều kiện áp dụng.
}
\end{ex}

% ===== Câu 170 =====
% ID: L12C4B24-51-TN
% Mức: NB
\dangbt{Điện hạt nhân, nhiên liệu hạt nhân và xử lí chất thải}
\begin{ex}
% ID: L12C4B24-51-TN
% Mức: NB
Phát biểu nào sau đây đúng về điện hạt nhân, nhiên liệu hạt nhân và xử lí chất thải?
\choice
{\True Nhà máy điện hạt nhân biến năng lượng hạt nhân thành nhiệt, cơ năng rồi điện năng; nhiên liệu và chất thải phải được quản lí nghiêm ngặt.}
{Chất thải phóng xạ có thể xử lí như rác sinh hoạt thông thường.}
{Có thể bỏ qua điều kiện áp dụng và đơn vị trong mọi trường hợp.}
{Kết quả của bài toán không phụ thuộc dữ kiện đã cho.}
\loigiai{
Nhà máy điện hạt nhân biến năng lượng hạt nhân thành nhiệt, cơ năng rồi điện năng; nhiên liệu và chất thải phải được quản lí nghiêm ngặt. Vì vậy phương án $A$ đúng; các phương án còn lại trái với kiến thức hoặc điều kiện áp dụng.
}
\end{ex}

% ===== Câu 171 =====
% ID: L12C4B24-52-TN
% Mức: NB
\begin{ex}
% ID: L12C4B24-52-TN
% Mức: NB
Nhận định nào phù hợp với kiến thức Vật lí về điện hạt nhân, nhiên liệu hạt nhân và xử lí chất thải?
\choice
{Chất thải phóng xạ có thể xử lí như rác sinh hoạt thông thường.}
{\True Nhà máy điện hạt nhân biến năng lượng hạt nhân thành nhiệt, cơ năng rồi điện năng; nhiên liệu và chất thải phải được quản lí nghiêm ngặt.}
{Có thể bỏ qua điều kiện áp dụng và đơn vị trong mọi trường hợp.}
{Kết quả của bài toán không phụ thuộc dữ kiện đã cho.}
\loigiai{
Nhà máy điện hạt nhân biến năng lượng hạt nhân thành nhiệt, cơ năng rồi điện năng; nhiên liệu và chất thải phải được quản lí nghiêm ngặt. Vì vậy phương án $B$ đúng; các phương án còn lại trái với kiến thức hoặc điều kiện áp dụng.
}
\end{ex}

% ===== Câu 172 =====
% ID: L12C4B24-53-TN
% Mức: NB
\begin{ex}
% ID: L12C4B24-53-TN
% Mức: NB
Chọn kết luận chính xác về điện hạt nhân, nhiên liệu hạt nhân và xử lí chất thải?
\choice
{Chất thải phóng xạ có thể xử lí như rác sinh hoạt thông thường.}
{Có thể bỏ qua điều kiện áp dụng và đơn vị trong mọi trường hợp.}
{\True Nhà máy điện hạt nhân biến năng lượng hạt nhân thành nhiệt, cơ năng rồi điện năng; nhiên liệu và chất thải phải được quản lí nghiêm ngặt.}
{Kết quả của bài toán không phụ thuộc dữ kiện đã cho.}
\loigiai{
Nhà máy điện hạt nhân biến năng lượng hạt nhân thành nhiệt, cơ năng rồi điện năng; nhiên liệu và chất thải phải được quản lí nghiêm ngặt. Vì vậy phương án $C$ đúng; các phương án còn lại trái với kiến thức hoặc điều kiện áp dụng.
}
\end{ex}

% ===== Câu 173 =====
% ID: L12C4B24-54-TN
% Mức: TH
\begin{ex}
% ID: L12C4B24-54-TN
% Mức: TH
Nội dung nào mô tả đúng nhất về điện hạt nhân, nhiên liệu hạt nhân và xử lí chất thải?
\choice
{Chất thải phóng xạ có thể xử lí như rác sinh hoạt thông thường.}
{Có thể bỏ qua điều kiện áp dụng và đơn vị trong mọi trường hợp.}
{Kết quả của bài toán không phụ thuộc dữ kiện đã cho.}
{\True Nhà máy điện hạt nhân biến năng lượng hạt nhân thành nhiệt, cơ năng rồi điện năng; nhiên liệu và chất thải phải được quản lí nghiêm ngặt.}
\loigiai{
Nhà máy điện hạt nhân biến năng lượng hạt nhân thành nhiệt, cơ năng rồi điện năng; nhiên liệu và chất thải phải được quản lí nghiêm ngặt. Vì vậy phương án $D$ đúng; các phương án còn lại trái với kiến thức hoặc điều kiện áp dụng.
}
\end{ex}

% ===== Câu 174 =====
% ID: L12C4B24-55-TN
% Mức: TH
\begin{ex}
% ID: L12C4B24-55-TN
% Mức: TH
Khi phân tích, phát biểu nào đúng về điện hạt nhân, nhiên liệu hạt nhân và xử lí chất thải?
\choice
{\True Nhà máy điện hạt nhân biến năng lượng hạt nhân thành nhiệt, cơ năng rồi điện năng; nhiên liệu và chất thải phải được quản lí nghiêm ngặt.}
{Chất thải phóng xạ có thể xử lí như rác sinh hoạt thông thường.}
{Có thể bỏ qua điều kiện áp dụng và đơn vị trong mọi trường hợp.}
{Kết quả của bài toán không phụ thuộc dữ kiện đã cho.}
\loigiai{
Nhà máy điện hạt nhân biến năng lượng hạt nhân thành nhiệt, cơ năng rồi điện năng; nhiên liệu và chất thải phải được quản lí nghiêm ngặt. Vì vậy phương án $A$ đúng; các phương án còn lại trái với kiến thức hoặc điều kiện áp dụng.
}
\end{ex}

% ===== Câu 175 =====
% ID: L12C4B24-56-TN
% Mức: TH
\begin{ex}
% ID: L12C4B24-56-TN
% Mức: TH
Kết luận nào của học sinh là đúng về điện hạt nhân, nhiên liệu hạt nhân và xử lí chất thải?
\choice
{Chất thải phóng xạ có thể xử lí như rác sinh hoạt thông thường.}
{\True Nhà máy điện hạt nhân biến năng lượng hạt nhân thành nhiệt, cơ năng rồi điện năng; nhiên liệu và chất thải phải được quản lí nghiêm ngặt.}
{Có thể bỏ qua điều kiện áp dụng và đơn vị trong mọi trường hợp.}
{Kết quả của bài toán không phụ thuộc dữ kiện đã cho.}
\loigiai{
Nhà máy điện hạt nhân biến năng lượng hạt nhân thành nhiệt, cơ năng rồi điện năng; nhiên liệu và chất thải phải được quản lí nghiêm ngặt. Vì vậy phương án $B$ đúng; các phương án còn lại trái với kiến thức hoặc điều kiện áp dụng.
}
\end{ex}

% ===== Câu 176 =====
% ID: L12C4B24-57-TN
% Mức: TH
\begin{ex}
% ID: L12C4B24-57-TN
% Mức: TH
Chọn phát biểu không trái với kiến thức về điện hạt nhân, nhiên liệu hạt nhân và xử lí chất thải?
\choice
{Chất thải phóng xạ có thể xử lí như rác sinh hoạt thông thường.}
{Có thể bỏ qua điều kiện áp dụng và đơn vị trong mọi trường hợp.}
{\True Nhà máy điện hạt nhân biến năng lượng hạt nhân thành nhiệt, cơ năng rồi điện năng; nhiên liệu và chất thải phải được quản lí nghiêm ngặt.}
{Kết quả của bài toán không phụ thuộc dữ kiện đã cho.}
\loigiai{
Nhà máy điện hạt nhân biến năng lượng hạt nhân thành nhiệt, cơ năng rồi điện năng; nhiên liệu và chất thải phải được quản lí nghiêm ngặt. Vì vậy phương án $C$ đúng; các phương án còn lại trái với kiến thức hoặc điều kiện áp dụng.
}
\end{ex}

% ===== Câu 177 =====
% ID: L12C4B24-58-TN
% Mức: VD
\begin{ex}
% ID: L12C4B24-58-TN
% Mức: VD
Cơ sở Vật lí đúng để giải bài là gì về điện hạt nhân, nhiên liệu hạt nhân và xử lí chất thải?
\choice
{Chất thải phóng xạ có thể xử lí như rác sinh hoạt thông thường.}
{Có thể bỏ qua điều kiện áp dụng và đơn vị trong mọi trường hợp.}
{Kết quả của bài toán không phụ thuộc dữ kiện đã cho.}
{\True Nhà máy điện hạt nhân biến năng lượng hạt nhân thành nhiệt, cơ năng rồi điện năng; nhiên liệu và chất thải phải được quản lí nghiêm ngặt.}
\loigiai{
Nhà máy điện hạt nhân biến năng lượng hạt nhân thành nhiệt, cơ năng rồi điện năng; nhiên liệu và chất thải phải được quản lí nghiêm ngặt. Vì vậy phương án $D$ đúng; các phương án còn lại trái với kiến thức hoặc điều kiện áp dụng.
}
\end{ex}

% ===== Câu 178 =====
% ID: L12C4B24-59-TN
% Mức: VD
\begin{ex}
% ID: L12C4B24-59-TN
% Mức: VD
Trong một bài thực hành, nhận xét nào đúng về điện hạt nhân, nhiên liệu hạt nhân và xử lí chất thải?
\choice
{\True Nhà máy điện hạt nhân biến năng lượng hạt nhân thành nhiệt, cơ năng rồi điện năng; nhiên liệu và chất thải phải được quản lí nghiêm ngặt.}
{Chất thải phóng xạ có thể xử lí như rác sinh hoạt thông thường.}
{Có thể bỏ qua điều kiện áp dụng và đơn vị trong mọi trường hợp.}
{Kết quả của bài toán không phụ thuộc dữ kiện đã cho.}
\loigiai{
Nhà máy điện hạt nhân biến năng lượng hạt nhân thành nhiệt, cơ năng rồi điện năng; nhiên liệu và chất thải phải được quản lí nghiêm ngặt. Vì vậy phương án $A$ đúng; các phương án còn lại trái với kiến thức hoặc điều kiện áp dụng.
}
\end{ex}

% ===== Câu 179 =====
% ID: L12C4B24-60-TN
% Mức: VD
\begin{ex}
% ID: L12C4B24-60-TN
% Mức: VD
Kết luận đúng khi vận dụng là về điện hạt nhân, nhiên liệu hạt nhân và xử lí chất thải?
\choice
{Chất thải phóng xạ có thể xử lí như rác sinh hoạt thông thường.}
{\True Nhà máy điện hạt nhân biến năng lượng hạt nhân thành nhiệt, cơ năng rồi điện năng; nhiên liệu và chất thải phải được quản lí nghiêm ngặt.}
{Có thể bỏ qua điều kiện áp dụng và đơn vị trong mọi trường hợp.}
{Kết quả của bài toán không phụ thuộc dữ kiện đã cho.}
\loigiai{
Nhà máy điện hạt nhân biến năng lượng hạt nhân thành nhiệt, cơ năng rồi điện năng; nhiên liệu và chất thải phải được quản lí nghiêm ngặt. Vì vậy phương án $B$ đúng; các phương án còn lại trái với kiến thức hoặc điều kiện áp dụng.
}
\end{ex}

% ===== Câu 180 =====
% ID: L12C4B24-61-TN
% Mức: NB
\dangbt{Ứng dụng công nghệ hạt nhân, an toàn và phát triển bền vững}
\begin{ex}
% ID: L12C4B24-61-TN
% Mức: NB
Phát biểu nào sau đây đúng về ứng dụng công nghệ hạt nhân, an toàn và phát triển bền vững?
\choice
{\True Công nghệ hạt nhân có ứng dụng trong năng lượng, y học, nông nghiệp và công nghiệp nhưng phải đánh giá rủi ro, an toàn và chất thải.}
{Ứng dụng hạt nhân hoàn toàn không tạo ra rủi ro môi trường.}
{Có thể bỏ qua điều kiện áp dụng và đơn vị trong mọi trường hợp.}
{Kết quả của bài toán không phụ thuộc dữ kiện đã cho.}
\loigiai{
Công nghệ hạt nhân có ứng dụng trong năng lượng, y học, nông nghiệp và công nghiệp nhưng phải đánh giá rủi ro, an toàn và chất thải. Vì vậy phương án $A$ đúng; các phương án còn lại trái với kiến thức hoặc điều kiện áp dụng.
}
\end{ex}

% ===== Câu 181 =====
% ID: L12C4B24-62-TN
% Mức: NB
\begin{ex}
% ID: L12C4B24-62-TN
% Mức: NB
Nhận định nào phù hợp với kiến thức Vật lí về ứng dụng công nghệ hạt nhân, an toàn và phát triển bền vững?
\choice
{Ứng dụng hạt nhân hoàn toàn không tạo ra rủi ro môi trường.}
{\True Công nghệ hạt nhân có ứng dụng trong năng lượng, y học, nông nghiệp và công nghiệp nhưng phải đánh giá rủi ro, an toàn và chất thải.}
{Có thể bỏ qua điều kiện áp dụng và đơn vị trong mọi trường hợp.}
{Kết quả của bài toán không phụ thuộc dữ kiện đã cho.}
\loigiai{
Công nghệ hạt nhân có ứng dụng trong năng lượng, y học, nông nghiệp và công nghiệp nhưng phải đánh giá rủi ro, an toàn và chất thải. Vì vậy phương án $B$ đúng; các phương án còn lại trái với kiến thức hoặc điều kiện áp dụng.
}
\end{ex}

% ===== Câu 182 =====
% ID: L12C4B24-63-TN
% Mức: NB
\begin{ex}
% ID: L12C4B24-63-TN
% Mức: NB
Chọn kết luận chính xác về ứng dụng công nghệ hạt nhân, an toàn và phát triển bền vững?
\choice
{Ứng dụng hạt nhân hoàn toàn không tạo ra rủi ro môi trường.}
{Có thể bỏ qua điều kiện áp dụng và đơn vị trong mọi trường hợp.}
{\True Công nghệ hạt nhân có ứng dụng trong năng lượng, y học, nông nghiệp và công nghiệp nhưng phải đánh giá rủi ro, an toàn và chất thải.}
{Kết quả của bài toán không phụ thuộc dữ kiện đã cho.}
\loigiai{
Công nghệ hạt nhân có ứng dụng trong năng lượng, y học, nông nghiệp và công nghiệp nhưng phải đánh giá rủi ro, an toàn và chất thải. Vì vậy phương án $C$ đúng; các phương án còn lại trái với kiến thức hoặc điều kiện áp dụng.
}
\end{ex}

% ===== Câu 183 =====
% ID: L12C4B24-64-TN
% Mức: TH
\begin{ex}
% ID: L12C4B24-64-TN
% Mức: TH
Nội dung nào mô tả đúng nhất về ứng dụng công nghệ hạt nhân, an toàn và phát triển bền vững?
\choice
{Ứng dụng hạt nhân hoàn toàn không tạo ra rủi ro môi trường.}
{Có thể bỏ qua điều kiện áp dụng và đơn vị trong mọi trường hợp.}
{Kết quả của bài toán không phụ thuộc dữ kiện đã cho.}
{\True Công nghệ hạt nhân có ứng dụng trong năng lượng, y học, nông nghiệp và công nghiệp nhưng phải đánh giá rủi ro, an toàn và chất thải.}
\loigiai{
Công nghệ hạt nhân có ứng dụng trong năng lượng, y học, nông nghiệp và công nghiệp nhưng phải đánh giá rủi ro, an toàn và chất thải. Vì vậy phương án $D$ đúng; các phương án còn lại trái với kiến thức hoặc điều kiện áp dụng.
}
\end{ex}

% ===== Câu 184 =====
% ID: L12C4B24-65-TN
% Mức: TH
\begin{ex}
% ID: L12C4B24-65-TN
% Mức: TH
Khi phân tích, phát biểu nào đúng về ứng dụng công nghệ hạt nhân, an toàn và phát triển bền vững?
\choice
{\True Công nghệ hạt nhân có ứng dụng trong năng lượng, y học, nông nghiệp và công nghiệp nhưng phải đánh giá rủi ro, an toàn và chất thải.}
{Ứng dụng hạt nhân hoàn toàn không tạo ra rủi ro môi trường.}
{Có thể bỏ qua điều kiện áp dụng và đơn vị trong mọi trường hợp.}
{Kết quả của bài toán không phụ thuộc dữ kiện đã cho.}
\loigiai{
Công nghệ hạt nhân có ứng dụng trong năng lượng, y học, nông nghiệp và công nghiệp nhưng phải đánh giá rủi ro, an toàn và chất thải. Vì vậy phương án $A$ đúng; các phương án còn lại trái với kiến thức hoặc điều kiện áp dụng.
}
\end{ex}

% ===== Câu 185 =====
% ID: L12C4B24-66-TN
% Mức: TH
\begin{ex}
% ID: L12C4B24-66-TN
% Mức: TH
Kết luận nào của học sinh là đúng về ứng dụng công nghệ hạt nhân, an toàn và phát triển bền vững?
\choice
{Ứng dụng hạt nhân hoàn toàn không tạo ra rủi ro môi trường.}
{\True Công nghệ hạt nhân có ứng dụng trong năng lượng, y học, nông nghiệp và công nghiệp nhưng phải đánh giá rủi ro, an toàn và chất thải.}
{Có thể bỏ qua điều kiện áp dụng và đơn vị trong mọi trường hợp.}
{Kết quả của bài toán không phụ thuộc dữ kiện đã cho.}
\loigiai{
Công nghệ hạt nhân có ứng dụng trong năng lượng, y học, nông nghiệp và công nghiệp nhưng phải đánh giá rủi ro, an toàn và chất thải. Vì vậy phương án $B$ đúng; các phương án còn lại trái với kiến thức hoặc điều kiện áp dụng.
}
\end{ex}

% ===== Câu 186 =====
% ID: L12C4B24-67-TN
% Mức: TH
\begin{ex}
% ID: L12C4B24-67-TN
% Mức: TH
Chọn phát biểu không trái với kiến thức về ứng dụng công nghệ hạt nhân, an toàn và phát triển bền vững?
\choice
{Ứng dụng hạt nhân hoàn toàn không tạo ra rủi ro môi trường.}
{Có thể bỏ qua điều kiện áp dụng và đơn vị trong mọi trường hợp.}
{\True Công nghệ hạt nhân có ứng dụng trong năng lượng, y học, nông nghiệp và công nghiệp nhưng phải đánh giá rủi ro, an toàn và chất thải.}
{Kết quả của bài toán không phụ thuộc dữ kiện đã cho.}
\loigiai{
Công nghệ hạt nhân có ứng dụng trong năng lượng, y học, nông nghiệp và công nghiệp nhưng phải đánh giá rủi ro, an toàn và chất thải. Vì vậy phương án $C$ đúng; các phương án còn lại trái với kiến thức hoặc điều kiện áp dụng.
}
\end{ex}

% ===== Câu 187 =====
% ID: L12C4B24-68-TN
% Mức: VD
\begin{ex}
% ID: L12C4B24-68-TN
% Mức: VD
Cơ sở Vật lí đúng để giải bài là gì về ứng dụng công nghệ hạt nhân, an toàn và phát triển bền vững?
\choice
{Ứng dụng hạt nhân hoàn toàn không tạo ra rủi ro môi trường.}
{Có thể bỏ qua điều kiện áp dụng và đơn vị trong mọi trường hợp.}
{Kết quả của bài toán không phụ thuộc dữ kiện đã cho.}
{\True Công nghệ hạt nhân có ứng dụng trong năng lượng, y học, nông nghiệp và công nghiệp nhưng phải đánh giá rủi ro, an toàn và chất thải.}
\loigiai{
Công nghệ hạt nhân có ứng dụng trong năng lượng, y học, nông nghiệp và công nghiệp nhưng phải đánh giá rủi ro, an toàn và chất thải. Vì vậy phương án $D$ đúng; các phương án còn lại trái với kiến thức hoặc điều kiện áp dụng.
}
\end{ex}

% ===== Câu 188 =====
% ID: L12C4B24-69-TN
% Mức: VD
\begin{ex}
% ID: L12C4B24-69-TN
% Mức: VD
Trong một bài thực hành, nhận xét nào đúng về ứng dụng công nghệ hạt nhân, an toàn và phát triển bền vững?
\choice
{\True Công nghệ hạt nhân có ứng dụng trong năng lượng, y học, nông nghiệp và công nghiệp nhưng phải đánh giá rủi ro, an toàn và chất thải.}
{Ứng dụng hạt nhân hoàn toàn không tạo ra rủi ro môi trường.}
{Có thể bỏ qua điều kiện áp dụng và đơn vị trong mọi trường hợp.}
{Kết quả của bài toán không phụ thuộc dữ kiện đã cho.}
\loigiai{
Công nghệ hạt nhân có ứng dụng trong năng lượng, y học, nông nghiệp và công nghiệp nhưng phải đánh giá rủi ro, an toàn và chất thải. Vì vậy phương án $A$ đúng; các phương án còn lại trái với kiến thức hoặc điều kiện áp dụng.
}
\end{ex}

% ===== Câu 189 =====
% ID: L12C4B24-70-TN
% Mức: VD
\begin{ex}
% ID: L12C4B24-70-TN
% Mức: VD
Kết luận đúng khi vận dụng là về ứng dụng công nghệ hạt nhân, an toàn và phát triển bền vững?
\choice
{Ứng dụng hạt nhân hoàn toàn không tạo ra rủi ro môi trường.}
{\True Công nghệ hạt nhân có ứng dụng trong năng lượng, y học, nông nghiệp và công nghiệp nhưng phải đánh giá rủi ro, an toàn và chất thải.}
{Có thể bỏ qua điều kiện áp dụng và đơn vị trong mọi trường hợp.}
{Kết quả của bài toán không phụ thuộc dữ kiện đã cho.}
\loigiai{
Công nghệ hạt nhân có ứng dụng trong năng lượng, y học, nông nghiệp và công nghiệp nhưng phải đánh giá rủi ro, an toàn và chất thải. Vì vậy phương án $B$ đúng; các phương án còn lại trái với kiến thức hoặc điều kiện áp dụng.
}
\end{ex}
